% ---------------------------------------------------------------------------
% arch.tex --- SynKit CRN at a glance: three routes in, one canonical object,
% views derived from it, and the analysis outputs that then apply.
%
% Source fragment: loaded directly by main.tex and SI/si.tex.
%
% Width.  The shared \scCanvas is pinned to the sn-jnl class default of 372 pt.
% paper_layout.tex replaces that text block with A4 at 22/24 mm margins, giving
% \textwidth = 466.63 pt = 16.40 cm, so a \scCanvas figure fills only 79.7% of
% the line.  This figure is therefore laid out on its own 16.40 cm canvas and
% included unscaled, so it fills the block at 1:1 and its type prints at the
% same physical size as every other figure's.
%
% Colour is semantic: green = structure, orange = the highlighted thing
% (the canonical object, the catalyst, the certificate), vermillion = a
% violation, and blue = a secondary/reference view.
%
% Text is deliberately sparse: the panels carry glyphs, the caption carries the
% qualifications.  Nothing here is a result.
% ---------------------------------------------------------------------------
% Resolve the shared style from the manuscript root or the SI directory.
\IfFileExists{synstyle.tex}{% ---------------------------------------------------------------------------
% synstyle.tex --- shared look for every TikZ figure of the SynKit CRN paper.
%
% Input it from a figure file:
%     \IfFileExists{synstyle.tex}{% ---------------------------------------------------------------------------
% synstyle.tex --- shared look for every TikZ figure of the SynKit CRN paper.
%
% Input it from a figure file:
%     \IfFileExists{synstyle.tex}{% ---------------------------------------------------------------------------
% synstyle.tex --- shared look for every TikZ figure of the SynKit CRN paper.
%
% Input it from a figure file:
%     \IfFileExists{synstyle.tex}{\input{synstyle}}{\input{Figure/synstyle}}
% so that the figure compiles both from Figure/ and from the paper root.
%
% The visual vocabulary follows ../../../../Style/style.tex: the
% colour-blind-safe Okabe--Ito accents, quiet neutral rules and wash, circular
% panel badges, and a compact technical type scale are shared by the companion
% papers.  Semantic aliases below keep the CRN meaning explicit.
%
% DESIGN CONTRACT -----------------------------------------------------------
% 1. Geometry.  Every figure is laid out on a canvas exactly \scW wide (the
%    sn-jnl text width, 372 pt) and is included unscaled, so one point of type
%    is the same physical size in all five figures.  Open each picture with
%    \scCanvas{<height in cm>} to pin the bounding box.
% 2. Colour is semantic, never decorative.  Four roles only:
%       scAccent (green)       structure, the normal/expected path
%       scMark   (orange)      catalyst, certificate, or highlight
%       scAlert  (vermillion)  failure, violation, or uncovered set
%       scRef    (blue)        secondary or reference series
%    A panel gets at most one accent.  Titles are always ink; panel identity
%    is also encoded by a letter inside a circular badge, not by colour alone.
% 3. Type scale.  Three sizes, no others, and nothing below 6 pt:
%       \scTitle (7 pt bold)  panel titles
%       \scBody  (6.5 pt)     labels, values, in-figure prose
%       \scSmall (6 pt)       units, footnotes, dense graph labels
% ---------------------------------------------------------------------------
% Load guard.  The paper inputs every figure source inline, so this file is
% reached once per figure; without the guard the second \newlength below
% would abort the run.
\ifdefined\SynKitStyleLoaded\endinput\fi
\newcommand{\SynKitStyleLoaded}{}

% Load in the document preamble. XeLaTeX/LuaLaTeX and an installed Arial
% family are required; fontspec reports an error instead of substituting.
\RequirePackage[no-math]{fontspec}
\newfontfamily\scArial{Arial}[NFSSFamily=scArial,Ligatures=TeX]
% Keep manuscript mathematics unchanged; select Arial letters and numerals
% only inside figures. TeX supplies specialist mathematical symbols.
\RequirePackage[subdued,italic]{mathastext}
\MTDeclareVersion{synarial}{TU}{scArial}{m}{n}
\SetMathAlphabet{\Mathsf}{synarial}{TU}{scArial}{m}{n}
\SetMathAlphabet{\Mathtt}{synarial}{TU}{scArial}{m}{n}

\usetikzlibrary{positioning,calc,fit,backgrounds,arrows.meta,shapes.geometric,
                shapes.misc,decorations.pathreplacing,patterns,math}

% --- canvas ---------------------------------------------------------------
% 372 pt text width minus the 1 pt standalone border on each side.
\newlength{\scW}\setlength{\scW}{13.0326cm}
\providecommand{\scWnum}{13.0326}
% \scCanvas{<height>} pins the picture to the full text width, y running down
% from 0, so every figure shares one coordinate convention.  It *sets* the
% bounding box rather than contributing to it, so the emitted page is exactly
% the text width and every figure is included at 1:1 -- otherwise a stray
% invisible extent (pgfplots reserves a little around an axis) rescales one
% figure relative to the others and the same \scSmall prints at two sizes.
% Anything drawn outside the canvas is therefore cropped: keep content inside,
% and check with
%     \providecommand{\scCanvas}[1]{\path[draw=red] (0,0) rectangle (\scWnum,-#1);}
\providecommand{\scCanvas}[1]{%
  \path[use as bounding box] (0,0) rectangle (\scWnum,-#1);}

% --- palette --------------------------------------------------------------
% Canonical companion-paper values (../../../../Style/style.tex).
\definecolor{npgOrange}{HTML}{E69F00}
\definecolor{npgSky}{HTML}{56B4E9}
\definecolor{npgGreen}{HTML}{009E73}
\definecolor{npgYellow}{HTML}{F0E442}
\definecolor{npgBlue}{HTML}{0072B2}
\definecolor{npgVermillion}{HTML}{E64B35}
\definecolor{npgPurple}{HTML}{CC79A7}
\definecolor{figInk}{HTML}{1A1A1A}
\definecolor{figMuted}{HTML}{666666}
\definecolor{figHair}{HTML}{D9D9D9}
\definecolor{figWash}{HTML}{F7F7F7}

\colorlet{scInk}{figInk}              % primary text
\colorlet{scInkM}{figMuted}           % secondary text
\colorlet{scPaper}{figWash}           % quiet cards and note strips
\colorlet{scSpF}{npgSky!12}           % species fill
\colorlet{scSpE}{npgBlue}             % species edge
\colorlet{scRxF}{npgOrange!12}        % reaction fill
\colorlet{scRxE}{npgOrange!88!black}  % reaction edge
\colorlet{scReac}{npgBlue!78!black}   % reactant edge
\colorlet{scProd}{npgGreen!82!black}  % product edge
\colorlet{scHiF}{npgOrange!13}        % highlight fill
\colorlet{scHiE}{npgOrange!82!black}  % highlight edge
% Stage / accent tints.  stDE remains available for the fifth plot series.
\colorlet{stAF}{npgGreen!9}\colorlet{stAE}{npgGreen}
\colorlet{stBF}{npgOrange!11}\colorlet{stBE}{npgOrange}
\colorlet{stCF}{npgBlue!9}\colorlet{stCE}{npgBlue}
\colorlet{stDF}{npgPurple!10}\colorlet{stDE}{npgPurple!85!black}
\colorlet{stEF}{npgVermillion!9}\colorlet{stEE}{npgVermillion}

% Semantic aliases.  Figures should speak in these, not in stAE/stBE/...
\colorlet{scAccent}{stAE}\colorlet{scAccentF}{stAF}   % structure / normal
\colorlet{scMark}{stBE}\colorlet{scMarkF}{stBF}       % highlight / certificate
\colorlet{scRef}{stCE}\colorlet{scRefF}{stCF}         % secondary series
\colorlet{scAlert}{stEE}\colorlet{scAlertF}{stEF}     % failure / violation
\colorlet{scRule}{figHair}                            % card and axis hairlines
\colorlet{scRuleL}{figHair!55}                        % grid, dividers

% --- text styles ----------------------------------------------------------
% Three sizes, fixed in points so they do not drift with the document class.
% Arial is shared by labels, headings, panel badges, and chemistry labels.
\providecommand{\scTitle}{\scArial\fontsize{7}{8.4}\selectfont\bfseries\color{scInk}}
\providecommand{\scBody}{\scArial\fontsize{6.5}{7.8}\selectfont\color{scInkM}}
\providecommand{\scSmall}{\scArial\fontsize{6}{7.2}\selectfont\color{scInkM}}
\providecommand{\scSub}[1]{{\scSmall\itshape #1}}
\providecommand{\scPanel}[1]{{\scTitle(#1)}}
% Numerals that carry a result: ink, not colour, so the eye reads the value.
\providecommand{\scNum}{\scArial\fontsize{10}{11}\selectfont\bfseries\color{scInk}}

% --- shared TikZ styles ---------------------------------------------------
\tikzset{
  % text never hyphenates inside a figure
  every node/.append style={execute at begin node={%
      \hyphenpenalty=10000 \exhyphenpenalty=10000 \hbadness=10000 }},
  syn/.style={font=\scArial,color=scInk,
              execute at begin picture={\MTversion*{synarial}},
              >={Stealth[round,length=2.4mm,width=1.9mm]}},
  % containers ------------------------------------------------------------
  % One card chrome for the whole paper: white, 3 pt radius, one hairline.
  panel/.style={rounded corners=2.5pt,fill=none,draw=scRule,line width=0.45pt},
  card/.style={rounded corners=2.2pt,fill=scPaper,draw=scRule,line width=0.45pt,
               inner sep=5pt,align=left,anchor=north west},
  % a quiet strip for a panel's closing fact
  note/.style={rounded corners=2.5pt,fill=scPaper,draw=scRuleL,line width=0.4pt,
               inner xsep=6pt,inner ysep=3.5pt,align=left,font=\scSmall},
  % small labelled value pill
  chip/.style={rounded corners=2pt,fill=scPaper,draw=scRule,line width=0.4pt,
               inner xsep=4pt,inner ysep=2pt,font=\scSmall},
  % Circular letter badge: the same accessible panel marker used in Style/*.
  hbullet/.style={circle,fill=scInk,text=white,draw=none,
                  font=\scArial\fontsize{6}{6}\selectfont\bfseries,
                  minimum size=4.2mm,inner sep=0pt},
  ctag/.style={font=\scTitle,anchor=north west,text=scInk},
  % CRN primitives --------------------------------------------------------
  species/.style={rounded rectangle,fill=scSpF,draw=scSpE,line width=0.45pt,
                  inner xsep=2.4pt,inner ysep=1.6pt,font=\scSmall,
                  text=scInk,minimum height=3.6mm},
  cofactor/.style={species,fill=scRefF,draw=scRef},
  catalyst/.style={species,fill=scHiF,draw=scHiE},
  rxn/.style={rectangle,fill=scRxF,draw=scRxE,line width=0.45pt,
              inner sep=1.2pt,font=\scSmall,text=scInk,
              minimum width=3.0mm,minimum height=3.0mm},
  place/.style={circle,fill=scSpF,draw=scSpE,line width=0.45pt,
                minimum size=4.0mm,inner sep=0pt},
  trans/.style={rectangle,fill=scRxE,draw=scRxE,minimum width=1.1mm,
                minimum height=3.6mm,inner sep=0pt},
  token/.style={circle,fill=scInk,minimum size=1.1mm,inner sep=0pt},
  % edges -----------------------------------------------------------------
  reac/.style={->,line width=0.5pt,draw=scReac,shorten >=0.8pt,shorten <=0.8pt},
  prod/.style={->,line width=0.5pt,draw=scProd,shorten >=0.8pt,shorten <=0.8pt},
  cat/.style={->,line width=0.5pt,draw=scHiE,shorten >=0.8pt,shorten <=0.8pt},
  flow/.style={-{Stealth[round,length=2.2mm,width=1.8mm]},line width=1.1pt,
               draw=scInkM!45},
  flab/.style={font=\scSmall,text=scInkM,inner sep=2pt},
  stoi/.style={font=\scSmall,text=scInkM,inner sep=0.8pt},
  % verdict marks: shape carries the meaning, colour only reinforces it
  vpass/.style={circle,fill=scAccentF,draw=scAccent,text=scAccent,
                font=\fontsize{6}{6}\selectfont\bfseries,
                minimum size=3.2mm,inner sep=0pt},
  vfail/.style={circle,fill=scAlertF,draw=scAlert,text=scAlert,
                font=\fontsize{6}{6}\selectfont\bfseries,
                minimum size=3.2mm,inner sep=0pt},
  callout/.style={font=\scSmall,text=scInkM,align=left},
}

% \scHead[<badge colour>]{<anchor coord>}{<letter>}{<title>}{<subtitle or empty>}
% One header treatment for every panel in the paper.
\providecommand{\scHead}[5][scInk]{%
  \node[hbullet,fill=#1,anchor=north west] (scbadge#3) at ($#2+(0,-0.015)$)
    {\MakeUppercase{#3}};
  \node[ctag,anchor=west] at ($(scbadge#3.east)+(0.14,0)$) {#4};
  \ifx\relax#5\relax\else
    \node[anchor=north west,font=\scSmall,text=scInkM]
      at ($#2+(0.56,-0.245)$) {#5};
  \fi}

% --- shared pgfplots look (only if the figure loaded pgfplots) ------------
\ifdefined\pgfplotsset
  \pgfplotsset{compat=1.18}
  \pgfplotsset{
    synaxis/.style={
      scale only axis,
      tick align=outside, tick pos=left,
      axis line style={draw=scInkM!55,line width=0.4pt},
      % One grid, on the value axis only.  No minor grid: at these sizes it
      % turns the plot frame into a grey wash.
      ymajorgrids=true, xmajorgrids=false,
      major grid style={draw=scRuleL,line width=0.3pt},
      major tick length=1.6pt,
      label style={font=\scBody,color=scInk},
      tick label style={font=\scSmall,color=scInkM},
      legend style={font=\scSmall,draw=none,fill=none,inner sep=1pt,
                    row sep=0.4pt,legend columns=1},
      legend cell align=left,
      every axis plot/.append style={line width=0.7pt,mark size=1.2pt},
    },
  }
\fi
}{% ---------------------------------------------------------------------------
% synstyle.tex --- shared look for every TikZ figure of the SynKit CRN paper.
%
% Input it from a figure file:
%     \IfFileExists{synstyle.tex}{\input{synstyle}}{\input{Figure/synstyle}}
% so that the figure compiles both from Figure/ and from the paper root.
%
% The visual vocabulary follows ../../../../Style/style.tex: the
% colour-blind-safe Okabe--Ito accents, quiet neutral rules and wash, circular
% panel badges, and a compact technical type scale are shared by the companion
% papers.  Semantic aliases below keep the CRN meaning explicit.
%
% DESIGN CONTRACT -----------------------------------------------------------
% 1. Geometry.  Every figure is laid out on a canvas exactly \scW wide (the
%    sn-jnl text width, 372 pt) and is included unscaled, so one point of type
%    is the same physical size in all five figures.  Open each picture with
%    \scCanvas{<height in cm>} to pin the bounding box.
% 2. Colour is semantic, never decorative.  Four roles only:
%       scAccent (green)       structure, the normal/expected path
%       scMark   (orange)      catalyst, certificate, or highlight
%       scAlert  (vermillion)  failure, violation, or uncovered set
%       scRef    (blue)        secondary or reference series
%    A panel gets at most one accent.  Titles are always ink; panel identity
%    is also encoded by a letter inside a circular badge, not by colour alone.
% 3. Type scale.  Three sizes, no others, and nothing below 6 pt:
%       \scTitle (7 pt bold)  panel titles
%       \scBody  (6.5 pt)     labels, values, in-figure prose
%       \scSmall (6 pt)       units, footnotes, dense graph labels
% ---------------------------------------------------------------------------
% Load guard.  The paper inputs every figure source inline, so this file is
% reached once per figure; without the guard the second \newlength below
% would abort the run.
\ifdefined\SynKitStyleLoaded\endinput\fi
\newcommand{\SynKitStyleLoaded}{}

% Load in the document preamble. XeLaTeX/LuaLaTeX and an installed Arial
% family are required; fontspec reports an error instead of substituting.
\RequirePackage[no-math]{fontspec}
\newfontfamily\scArial{Arial}[NFSSFamily=scArial,Ligatures=TeX]
% Keep manuscript mathematics unchanged; select Arial letters and numerals
% only inside figures. TeX supplies specialist mathematical symbols.
\RequirePackage[subdued,italic]{mathastext}
\MTDeclareVersion{synarial}{TU}{scArial}{m}{n}
\SetMathAlphabet{\Mathsf}{synarial}{TU}{scArial}{m}{n}
\SetMathAlphabet{\Mathtt}{synarial}{TU}{scArial}{m}{n}

\usetikzlibrary{positioning,calc,fit,backgrounds,arrows.meta,shapes.geometric,
                shapes.misc,decorations.pathreplacing,patterns,math}

% --- canvas ---------------------------------------------------------------
% 372 pt text width minus the 1 pt standalone border on each side.
\newlength{\scW}\setlength{\scW}{13.0326cm}
\providecommand{\scWnum}{13.0326}
% \scCanvas{<height>} pins the picture to the full text width, y running down
% from 0, so every figure shares one coordinate convention.  It *sets* the
% bounding box rather than contributing to it, so the emitted page is exactly
% the text width and every figure is included at 1:1 -- otherwise a stray
% invisible extent (pgfplots reserves a little around an axis) rescales one
% figure relative to the others and the same \scSmall prints at two sizes.
% Anything drawn outside the canvas is therefore cropped: keep content inside,
% and check with
%     \providecommand{\scCanvas}[1]{\path[draw=red] (0,0) rectangle (\scWnum,-#1);}
\providecommand{\scCanvas}[1]{%
  \path[use as bounding box] (0,0) rectangle (\scWnum,-#1);}

% --- palette --------------------------------------------------------------
% Canonical companion-paper values (../../../../Style/style.tex).
\definecolor{npgOrange}{HTML}{E69F00}
\definecolor{npgSky}{HTML}{56B4E9}
\definecolor{npgGreen}{HTML}{009E73}
\definecolor{npgYellow}{HTML}{F0E442}
\definecolor{npgBlue}{HTML}{0072B2}
\definecolor{npgVermillion}{HTML}{E64B35}
\definecolor{npgPurple}{HTML}{CC79A7}
\definecolor{figInk}{HTML}{1A1A1A}
\definecolor{figMuted}{HTML}{666666}
\definecolor{figHair}{HTML}{D9D9D9}
\definecolor{figWash}{HTML}{F7F7F7}

\colorlet{scInk}{figInk}              % primary text
\colorlet{scInkM}{figMuted}           % secondary text
\colorlet{scPaper}{figWash}           % quiet cards and note strips
\colorlet{scSpF}{npgSky!12}           % species fill
\colorlet{scSpE}{npgBlue}             % species edge
\colorlet{scRxF}{npgOrange!12}        % reaction fill
\colorlet{scRxE}{npgOrange!88!black}  % reaction edge
\colorlet{scReac}{npgBlue!78!black}   % reactant edge
\colorlet{scProd}{npgGreen!82!black}  % product edge
\colorlet{scHiF}{npgOrange!13}        % highlight fill
\colorlet{scHiE}{npgOrange!82!black}  % highlight edge
% Stage / accent tints.  stDE remains available for the fifth plot series.
\colorlet{stAF}{npgGreen!9}\colorlet{stAE}{npgGreen}
\colorlet{stBF}{npgOrange!11}\colorlet{stBE}{npgOrange}
\colorlet{stCF}{npgBlue!9}\colorlet{stCE}{npgBlue}
\colorlet{stDF}{npgPurple!10}\colorlet{stDE}{npgPurple!85!black}
\colorlet{stEF}{npgVermillion!9}\colorlet{stEE}{npgVermillion}

% Semantic aliases.  Figures should speak in these, not in stAE/stBE/...
\colorlet{scAccent}{stAE}\colorlet{scAccentF}{stAF}   % structure / normal
\colorlet{scMark}{stBE}\colorlet{scMarkF}{stBF}       % highlight / certificate
\colorlet{scRef}{stCE}\colorlet{scRefF}{stCF}         % secondary series
\colorlet{scAlert}{stEE}\colorlet{scAlertF}{stEF}     % failure / violation
\colorlet{scRule}{figHair}                            % card and axis hairlines
\colorlet{scRuleL}{figHair!55}                        % grid, dividers

% --- text styles ----------------------------------------------------------
% Three sizes, fixed in points so they do not drift with the document class.
% Arial is shared by labels, headings, panel badges, and chemistry labels.
\providecommand{\scTitle}{\scArial\fontsize{7}{8.4}\selectfont\bfseries\color{scInk}}
\providecommand{\scBody}{\scArial\fontsize{6.5}{7.8}\selectfont\color{scInkM}}
\providecommand{\scSmall}{\scArial\fontsize{6}{7.2}\selectfont\color{scInkM}}
\providecommand{\scSub}[1]{{\scSmall\itshape #1}}
\providecommand{\scPanel}[1]{{\scTitle(#1)}}
% Numerals that carry a result: ink, not colour, so the eye reads the value.
\providecommand{\scNum}{\scArial\fontsize{10}{11}\selectfont\bfseries\color{scInk}}

% --- shared TikZ styles ---------------------------------------------------
\tikzset{
  % text never hyphenates inside a figure
  every node/.append style={execute at begin node={%
      \hyphenpenalty=10000 \exhyphenpenalty=10000 \hbadness=10000 }},
  syn/.style={font=\scArial,color=scInk,
              execute at begin picture={\MTversion*{synarial}},
              >={Stealth[round,length=2.4mm,width=1.9mm]}},
  % containers ------------------------------------------------------------
  % One card chrome for the whole paper: white, 3 pt radius, one hairline.
  panel/.style={rounded corners=2.5pt,fill=none,draw=scRule,line width=0.45pt},
  card/.style={rounded corners=2.2pt,fill=scPaper,draw=scRule,line width=0.45pt,
               inner sep=5pt,align=left,anchor=north west},
  % a quiet strip for a panel's closing fact
  note/.style={rounded corners=2.5pt,fill=scPaper,draw=scRuleL,line width=0.4pt,
               inner xsep=6pt,inner ysep=3.5pt,align=left,font=\scSmall},
  % small labelled value pill
  chip/.style={rounded corners=2pt,fill=scPaper,draw=scRule,line width=0.4pt,
               inner xsep=4pt,inner ysep=2pt,font=\scSmall},
  % Circular letter badge: the same accessible panel marker used in Style/*.
  hbullet/.style={circle,fill=scInk,text=white,draw=none,
                  font=\scArial\fontsize{6}{6}\selectfont\bfseries,
                  minimum size=4.2mm,inner sep=0pt},
  ctag/.style={font=\scTitle,anchor=north west,text=scInk},
  % CRN primitives --------------------------------------------------------
  species/.style={rounded rectangle,fill=scSpF,draw=scSpE,line width=0.45pt,
                  inner xsep=2.4pt,inner ysep=1.6pt,font=\scSmall,
                  text=scInk,minimum height=3.6mm},
  cofactor/.style={species,fill=scRefF,draw=scRef},
  catalyst/.style={species,fill=scHiF,draw=scHiE},
  rxn/.style={rectangle,fill=scRxF,draw=scRxE,line width=0.45pt,
              inner sep=1.2pt,font=\scSmall,text=scInk,
              minimum width=3.0mm,minimum height=3.0mm},
  place/.style={circle,fill=scSpF,draw=scSpE,line width=0.45pt,
                minimum size=4.0mm,inner sep=0pt},
  trans/.style={rectangle,fill=scRxE,draw=scRxE,minimum width=1.1mm,
                minimum height=3.6mm,inner sep=0pt},
  token/.style={circle,fill=scInk,minimum size=1.1mm,inner sep=0pt},
  % edges -----------------------------------------------------------------
  reac/.style={->,line width=0.5pt,draw=scReac,shorten >=0.8pt,shorten <=0.8pt},
  prod/.style={->,line width=0.5pt,draw=scProd,shorten >=0.8pt,shorten <=0.8pt},
  cat/.style={->,line width=0.5pt,draw=scHiE,shorten >=0.8pt,shorten <=0.8pt},
  flow/.style={-{Stealth[round,length=2.2mm,width=1.8mm]},line width=1.1pt,
               draw=scInkM!45},
  flab/.style={font=\scSmall,text=scInkM,inner sep=2pt},
  stoi/.style={font=\scSmall,text=scInkM,inner sep=0.8pt},
  % verdict marks: shape carries the meaning, colour only reinforces it
  vpass/.style={circle,fill=scAccentF,draw=scAccent,text=scAccent,
                font=\fontsize{6}{6}\selectfont\bfseries,
                minimum size=3.2mm,inner sep=0pt},
  vfail/.style={circle,fill=scAlertF,draw=scAlert,text=scAlert,
                font=\fontsize{6}{6}\selectfont\bfseries,
                minimum size=3.2mm,inner sep=0pt},
  callout/.style={font=\scSmall,text=scInkM,align=left},
}

% \scHead[<badge colour>]{<anchor coord>}{<letter>}{<title>}{<subtitle or empty>}
% One header treatment for every panel in the paper.
\providecommand{\scHead}[5][scInk]{%
  \node[hbullet,fill=#1,anchor=north west] (scbadge#3) at ($#2+(0,-0.015)$)
    {\MakeUppercase{#3}};
  \node[ctag,anchor=west] at ($(scbadge#3.east)+(0.14,0)$) {#4};
  \ifx\relax#5\relax\else
    \node[anchor=north west,font=\scSmall,text=scInkM]
      at ($#2+(0.56,-0.245)$) {#5};
  \fi}

% --- shared pgfplots look (only if the figure loaded pgfplots) ------------
\ifdefined\pgfplotsset
  \pgfplotsset{compat=1.18}
  \pgfplotsset{
    synaxis/.style={
      scale only axis,
      tick align=outside, tick pos=left,
      axis line style={draw=scInkM!55,line width=0.4pt},
      % One grid, on the value axis only.  No minor grid: at these sizes it
      % turns the plot frame into a grey wash.
      ymajorgrids=true, xmajorgrids=false,
      major grid style={draw=scRuleL,line width=0.3pt},
      major tick length=1.6pt,
      label style={font=\scBody,color=scInk},
      tick label style={font=\scSmall,color=scInkM},
      legend style={font=\scSmall,draw=none,fill=none,inner sep=1pt,
                    row sep=0.4pt,legend columns=1},
      legend cell align=left,
      every axis plot/.append style={line width=0.7pt,mark size=1.2pt},
    },
  }
\fi
}
% so that the figure compiles both from Figure/ and from the paper root.
%
% The visual vocabulary follows ../../../../Style/style.tex: the
% colour-blind-safe Okabe--Ito accents, quiet neutral rules and wash, circular
% panel badges, and a compact technical type scale are shared by the companion
% papers.  Semantic aliases below keep the CRN meaning explicit.
%
% DESIGN CONTRACT -----------------------------------------------------------
% 1. Geometry.  Every figure is laid out on a canvas exactly \scW wide (the
%    sn-jnl text width, 372 pt) and is included unscaled, so one point of type
%    is the same physical size in all five figures.  Open each picture with
%    \scCanvas{<height in cm>} to pin the bounding box.
% 2. Colour is semantic, never decorative.  Four roles only:
%       scAccent (green)       structure, the normal/expected path
%       scMark   (orange)      catalyst, certificate, or highlight
%       scAlert  (vermillion)  failure, violation, or uncovered set
%       scRef    (blue)        secondary or reference series
%    A panel gets at most one accent.  Titles are always ink; panel identity
%    is also encoded by a letter inside a circular badge, not by colour alone.
% 3. Type scale.  Three sizes, no others, and nothing below 6 pt:
%       \scTitle (7 pt bold)  panel titles
%       \scBody  (6.5 pt)     labels, values, in-figure prose
%       \scSmall (6 pt)       units, footnotes, dense graph labels
% ---------------------------------------------------------------------------
% Load guard.  The paper inputs every figure source inline, so this file is
% reached once per figure; without the guard the second \newlength below
% would abort the run.
\ifdefined\SynKitStyleLoaded\endinput\fi
\newcommand{\SynKitStyleLoaded}{}

% Load in the document preamble. XeLaTeX/LuaLaTeX and an installed Arial
% family are required; fontspec reports an error instead of substituting.
\RequirePackage[no-math]{fontspec}
\newfontfamily\scArial{Arial}[NFSSFamily=scArial,Ligatures=TeX]
% Keep manuscript mathematics unchanged; select Arial letters and numerals
% only inside figures. TeX supplies specialist mathematical symbols.
\RequirePackage[subdued,italic]{mathastext}
\MTDeclareVersion{synarial}{TU}{scArial}{m}{n}
\SetMathAlphabet{\Mathsf}{synarial}{TU}{scArial}{m}{n}
\SetMathAlphabet{\Mathtt}{synarial}{TU}{scArial}{m}{n}

\usetikzlibrary{positioning,calc,fit,backgrounds,arrows.meta,shapes.geometric,
                shapes.misc,decorations.pathreplacing,patterns,math}

% --- canvas ---------------------------------------------------------------
% 372 pt text width minus the 1 pt standalone border on each side.
\newlength{\scW}\setlength{\scW}{13.0326cm}
\providecommand{\scWnum}{13.0326}
% \scCanvas{<height>} pins the picture to the full text width, y running down
% from 0, so every figure shares one coordinate convention.  It *sets* the
% bounding box rather than contributing to it, so the emitted page is exactly
% the text width and every figure is included at 1:1 -- otherwise a stray
% invisible extent (pgfplots reserves a little around an axis) rescales one
% figure relative to the others and the same \scSmall prints at two sizes.
% Anything drawn outside the canvas is therefore cropped: keep content inside,
% and check with
%     \providecommand{\scCanvas}[1]{\path[draw=red] (0,0) rectangle (\scWnum,-#1);}
\providecommand{\scCanvas}[1]{%
  \path[use as bounding box] (0,0) rectangle (\scWnum,-#1);}

% --- palette --------------------------------------------------------------
% Canonical companion-paper values (../../../../Style/style.tex).
\definecolor{npgOrange}{HTML}{E69F00}
\definecolor{npgSky}{HTML}{56B4E9}
\definecolor{npgGreen}{HTML}{009E73}
\definecolor{npgYellow}{HTML}{F0E442}
\definecolor{npgBlue}{HTML}{0072B2}
\definecolor{npgVermillion}{HTML}{E64B35}
\definecolor{npgPurple}{HTML}{CC79A7}
\definecolor{figInk}{HTML}{1A1A1A}
\definecolor{figMuted}{HTML}{666666}
\definecolor{figHair}{HTML}{D9D9D9}
\definecolor{figWash}{HTML}{F7F7F7}

\colorlet{scInk}{figInk}              % primary text
\colorlet{scInkM}{figMuted}           % secondary text
\colorlet{scPaper}{figWash}           % quiet cards and note strips
\colorlet{scSpF}{npgSky!12}           % species fill
\colorlet{scSpE}{npgBlue}             % species edge
\colorlet{scRxF}{npgOrange!12}        % reaction fill
\colorlet{scRxE}{npgOrange!88!black}  % reaction edge
\colorlet{scReac}{npgBlue!78!black}   % reactant edge
\colorlet{scProd}{npgGreen!82!black}  % product edge
\colorlet{scHiF}{npgOrange!13}        % highlight fill
\colorlet{scHiE}{npgOrange!82!black}  % highlight edge
% Stage / accent tints.  stDE remains available for the fifth plot series.
\colorlet{stAF}{npgGreen!9}\colorlet{stAE}{npgGreen}
\colorlet{stBF}{npgOrange!11}\colorlet{stBE}{npgOrange}
\colorlet{stCF}{npgBlue!9}\colorlet{stCE}{npgBlue}
\colorlet{stDF}{npgPurple!10}\colorlet{stDE}{npgPurple!85!black}
\colorlet{stEF}{npgVermillion!9}\colorlet{stEE}{npgVermillion}

% Semantic aliases.  Figures should speak in these, not in stAE/stBE/...
\colorlet{scAccent}{stAE}\colorlet{scAccentF}{stAF}   % structure / normal
\colorlet{scMark}{stBE}\colorlet{scMarkF}{stBF}       % highlight / certificate
\colorlet{scRef}{stCE}\colorlet{scRefF}{stCF}         % secondary series
\colorlet{scAlert}{stEE}\colorlet{scAlertF}{stEF}     % failure / violation
\colorlet{scRule}{figHair}                            % card and axis hairlines
\colorlet{scRuleL}{figHair!55}                        % grid, dividers

% --- text styles ----------------------------------------------------------
% Three sizes, fixed in points so they do not drift with the document class.
% Arial is shared by labels, headings, panel badges, and chemistry labels.
\providecommand{\scTitle}{\scArial\fontsize{7}{8.4}\selectfont\bfseries\color{scInk}}
\providecommand{\scBody}{\scArial\fontsize{6.5}{7.8}\selectfont\color{scInkM}}
\providecommand{\scSmall}{\scArial\fontsize{6}{7.2}\selectfont\color{scInkM}}
\providecommand{\scSub}[1]{{\scSmall\itshape #1}}
\providecommand{\scPanel}[1]{{\scTitle(#1)}}
% Numerals that carry a result: ink, not colour, so the eye reads the value.
\providecommand{\scNum}{\scArial\fontsize{10}{11}\selectfont\bfseries\color{scInk}}

% --- shared TikZ styles ---------------------------------------------------
\tikzset{
  % text never hyphenates inside a figure
  every node/.append style={execute at begin node={%
      \hyphenpenalty=10000 \exhyphenpenalty=10000 \hbadness=10000 }},
  syn/.style={font=\scArial,color=scInk,
              execute at begin picture={\MTversion*{synarial}},
              >={Stealth[round,length=2.4mm,width=1.9mm]}},
  % containers ------------------------------------------------------------
  % One card chrome for the whole paper: white, 3 pt radius, one hairline.
  panel/.style={rounded corners=2.5pt,fill=none,draw=scRule,line width=0.45pt},
  card/.style={rounded corners=2.2pt,fill=scPaper,draw=scRule,line width=0.45pt,
               inner sep=5pt,align=left,anchor=north west},
  % a quiet strip for a panel's closing fact
  note/.style={rounded corners=2.5pt,fill=scPaper,draw=scRuleL,line width=0.4pt,
               inner xsep=6pt,inner ysep=3.5pt,align=left,font=\scSmall},
  % small labelled value pill
  chip/.style={rounded corners=2pt,fill=scPaper,draw=scRule,line width=0.4pt,
               inner xsep=4pt,inner ysep=2pt,font=\scSmall},
  % Circular letter badge: the same accessible panel marker used in Style/*.
  hbullet/.style={circle,fill=scInk,text=white,draw=none,
                  font=\scArial\fontsize{6}{6}\selectfont\bfseries,
                  minimum size=4.2mm,inner sep=0pt},
  ctag/.style={font=\scTitle,anchor=north west,text=scInk},
  % CRN primitives --------------------------------------------------------
  species/.style={rounded rectangle,fill=scSpF,draw=scSpE,line width=0.45pt,
                  inner xsep=2.4pt,inner ysep=1.6pt,font=\scSmall,
                  text=scInk,minimum height=3.6mm},
  cofactor/.style={species,fill=scRefF,draw=scRef},
  catalyst/.style={species,fill=scHiF,draw=scHiE},
  rxn/.style={rectangle,fill=scRxF,draw=scRxE,line width=0.45pt,
              inner sep=1.2pt,font=\scSmall,text=scInk,
              minimum width=3.0mm,minimum height=3.0mm},
  place/.style={circle,fill=scSpF,draw=scSpE,line width=0.45pt,
                minimum size=4.0mm,inner sep=0pt},
  trans/.style={rectangle,fill=scRxE,draw=scRxE,minimum width=1.1mm,
                minimum height=3.6mm,inner sep=0pt},
  token/.style={circle,fill=scInk,minimum size=1.1mm,inner sep=0pt},
  % edges -----------------------------------------------------------------
  reac/.style={->,line width=0.5pt,draw=scReac,shorten >=0.8pt,shorten <=0.8pt},
  prod/.style={->,line width=0.5pt,draw=scProd,shorten >=0.8pt,shorten <=0.8pt},
  cat/.style={->,line width=0.5pt,draw=scHiE,shorten >=0.8pt,shorten <=0.8pt},
  flow/.style={-{Stealth[round,length=2.2mm,width=1.8mm]},line width=1.1pt,
               draw=scInkM!45},
  flab/.style={font=\scSmall,text=scInkM,inner sep=2pt},
  stoi/.style={font=\scSmall,text=scInkM,inner sep=0.8pt},
  % verdict marks: shape carries the meaning, colour only reinforces it
  vpass/.style={circle,fill=scAccentF,draw=scAccent,text=scAccent,
                font=\fontsize{6}{6}\selectfont\bfseries,
                minimum size=3.2mm,inner sep=0pt},
  vfail/.style={circle,fill=scAlertF,draw=scAlert,text=scAlert,
                font=\fontsize{6}{6}\selectfont\bfseries,
                minimum size=3.2mm,inner sep=0pt},
  callout/.style={font=\scSmall,text=scInkM,align=left},
}

% \scHead[<badge colour>]{<anchor coord>}{<letter>}{<title>}{<subtitle or empty>}
% One header treatment for every panel in the paper.
\providecommand{\scHead}[5][scInk]{%
  \node[hbullet,fill=#1,anchor=north west] (scbadge#3) at ($#2+(0,-0.015)$)
    {\MakeUppercase{#3}};
  \node[ctag,anchor=west] at ($(scbadge#3.east)+(0.14,0)$) {#4};
  \ifx\relax#5\relax\else
    \node[anchor=north west,font=\scSmall,text=scInkM]
      at ($#2+(0.56,-0.245)$) {#5};
  \fi}

% --- shared pgfplots look (only if the figure loaded pgfplots) ------------
\ifdefined\pgfplotsset
  \pgfplotsset{compat=1.18}
  \pgfplotsset{
    synaxis/.style={
      scale only axis,
      tick align=outside, tick pos=left,
      axis line style={draw=scInkM!55,line width=0.4pt},
      % One grid, on the value axis only.  No minor grid: at these sizes it
      % turns the plot frame into a grey wash.
      ymajorgrids=true, xmajorgrids=false,
      major grid style={draw=scRuleL,line width=0.3pt},
      major tick length=1.6pt,
      label style={font=\scBody,color=scInk},
      tick label style={font=\scSmall,color=scInkM},
      legend style={font=\scSmall,draw=none,fill=none,inner sep=1pt,
                    row sep=0.4pt,legend columns=1},
      legend cell align=left,
      every axis plot/.append style={line width=0.7pt,mark size=1.2pt},
    },
  }
\fi
}{% ---------------------------------------------------------------------------
% synstyle.tex --- shared look for every TikZ figure of the SynKit CRN paper.
%
% Input it from a figure file:
%     \IfFileExists{synstyle.tex}{% ---------------------------------------------------------------------------
% synstyle.tex --- shared look for every TikZ figure of the SynKit CRN paper.
%
% Input it from a figure file:
%     \IfFileExists{synstyle.tex}{\input{synstyle}}{\input{Figure/synstyle}}
% so that the figure compiles both from Figure/ and from the paper root.
%
% The visual vocabulary follows ../../../../Style/style.tex: the
% colour-blind-safe Okabe--Ito accents, quiet neutral rules and wash, circular
% panel badges, and a compact technical type scale are shared by the companion
% papers.  Semantic aliases below keep the CRN meaning explicit.
%
% DESIGN CONTRACT -----------------------------------------------------------
% 1. Geometry.  Every figure is laid out on a canvas exactly \scW wide (the
%    sn-jnl text width, 372 pt) and is included unscaled, so one point of type
%    is the same physical size in all five figures.  Open each picture with
%    \scCanvas{<height in cm>} to pin the bounding box.
% 2. Colour is semantic, never decorative.  Four roles only:
%       scAccent (green)       structure, the normal/expected path
%       scMark   (orange)      catalyst, certificate, or highlight
%       scAlert  (vermillion)  failure, violation, or uncovered set
%       scRef    (blue)        secondary or reference series
%    A panel gets at most one accent.  Titles are always ink; panel identity
%    is also encoded by a letter inside a circular badge, not by colour alone.
% 3. Type scale.  Three sizes, no others, and nothing below 6 pt:
%       \scTitle (7 pt bold)  panel titles
%       \scBody  (6.5 pt)     labels, values, in-figure prose
%       \scSmall (6 pt)       units, footnotes, dense graph labels
% ---------------------------------------------------------------------------
% Load guard.  The paper inputs every figure source inline, so this file is
% reached once per figure; without the guard the second \newlength below
% would abort the run.
\ifdefined\SynKitStyleLoaded\endinput\fi
\newcommand{\SynKitStyleLoaded}{}

% Load in the document preamble. XeLaTeX/LuaLaTeX and an installed Arial
% family are required; fontspec reports an error instead of substituting.
\RequirePackage[no-math]{fontspec}
\newfontfamily\scArial{Arial}[NFSSFamily=scArial,Ligatures=TeX]
% Keep manuscript mathematics unchanged; select Arial letters and numerals
% only inside figures. TeX supplies specialist mathematical symbols.
\RequirePackage[subdued,italic]{mathastext}
\MTDeclareVersion{synarial}{TU}{scArial}{m}{n}
\SetMathAlphabet{\Mathsf}{synarial}{TU}{scArial}{m}{n}
\SetMathAlphabet{\Mathtt}{synarial}{TU}{scArial}{m}{n}

\usetikzlibrary{positioning,calc,fit,backgrounds,arrows.meta,shapes.geometric,
                shapes.misc,decorations.pathreplacing,patterns,math}

% --- canvas ---------------------------------------------------------------
% 372 pt text width minus the 1 pt standalone border on each side.
\newlength{\scW}\setlength{\scW}{13.0326cm}
\providecommand{\scWnum}{13.0326}
% \scCanvas{<height>} pins the picture to the full text width, y running down
% from 0, so every figure shares one coordinate convention.  It *sets* the
% bounding box rather than contributing to it, so the emitted page is exactly
% the text width and every figure is included at 1:1 -- otherwise a stray
% invisible extent (pgfplots reserves a little around an axis) rescales one
% figure relative to the others and the same \scSmall prints at two sizes.
% Anything drawn outside the canvas is therefore cropped: keep content inside,
% and check with
%     \providecommand{\scCanvas}[1]{\path[draw=red] (0,0) rectangle (\scWnum,-#1);}
\providecommand{\scCanvas}[1]{%
  \path[use as bounding box] (0,0) rectangle (\scWnum,-#1);}

% --- palette --------------------------------------------------------------
% Canonical companion-paper values (../../../../Style/style.tex).
\definecolor{npgOrange}{HTML}{E69F00}
\definecolor{npgSky}{HTML}{56B4E9}
\definecolor{npgGreen}{HTML}{009E73}
\definecolor{npgYellow}{HTML}{F0E442}
\definecolor{npgBlue}{HTML}{0072B2}
\definecolor{npgVermillion}{HTML}{E64B35}
\definecolor{npgPurple}{HTML}{CC79A7}
\definecolor{figInk}{HTML}{1A1A1A}
\definecolor{figMuted}{HTML}{666666}
\definecolor{figHair}{HTML}{D9D9D9}
\definecolor{figWash}{HTML}{F7F7F7}

\colorlet{scInk}{figInk}              % primary text
\colorlet{scInkM}{figMuted}           % secondary text
\colorlet{scPaper}{figWash}           % quiet cards and note strips
\colorlet{scSpF}{npgSky!12}           % species fill
\colorlet{scSpE}{npgBlue}             % species edge
\colorlet{scRxF}{npgOrange!12}        % reaction fill
\colorlet{scRxE}{npgOrange!88!black}  % reaction edge
\colorlet{scReac}{npgBlue!78!black}   % reactant edge
\colorlet{scProd}{npgGreen!82!black}  % product edge
\colorlet{scHiF}{npgOrange!13}        % highlight fill
\colorlet{scHiE}{npgOrange!82!black}  % highlight edge
% Stage / accent tints.  stDE remains available for the fifth plot series.
\colorlet{stAF}{npgGreen!9}\colorlet{stAE}{npgGreen}
\colorlet{stBF}{npgOrange!11}\colorlet{stBE}{npgOrange}
\colorlet{stCF}{npgBlue!9}\colorlet{stCE}{npgBlue}
\colorlet{stDF}{npgPurple!10}\colorlet{stDE}{npgPurple!85!black}
\colorlet{stEF}{npgVermillion!9}\colorlet{stEE}{npgVermillion}

% Semantic aliases.  Figures should speak in these, not in stAE/stBE/...
\colorlet{scAccent}{stAE}\colorlet{scAccentF}{stAF}   % structure / normal
\colorlet{scMark}{stBE}\colorlet{scMarkF}{stBF}       % highlight / certificate
\colorlet{scRef}{stCE}\colorlet{scRefF}{stCF}         % secondary series
\colorlet{scAlert}{stEE}\colorlet{scAlertF}{stEF}     % failure / violation
\colorlet{scRule}{figHair}                            % card and axis hairlines
\colorlet{scRuleL}{figHair!55}                        % grid, dividers

% --- text styles ----------------------------------------------------------
% Three sizes, fixed in points so they do not drift with the document class.
% Arial is shared by labels, headings, panel badges, and chemistry labels.
\providecommand{\scTitle}{\scArial\fontsize{7}{8.4}\selectfont\bfseries\color{scInk}}
\providecommand{\scBody}{\scArial\fontsize{6.5}{7.8}\selectfont\color{scInkM}}
\providecommand{\scSmall}{\scArial\fontsize{6}{7.2}\selectfont\color{scInkM}}
\providecommand{\scSub}[1]{{\scSmall\itshape #1}}
\providecommand{\scPanel}[1]{{\scTitle(#1)}}
% Numerals that carry a result: ink, not colour, so the eye reads the value.
\providecommand{\scNum}{\scArial\fontsize{10}{11}\selectfont\bfseries\color{scInk}}

% --- shared TikZ styles ---------------------------------------------------
\tikzset{
  % text never hyphenates inside a figure
  every node/.append style={execute at begin node={%
      \hyphenpenalty=10000 \exhyphenpenalty=10000 \hbadness=10000 }},
  syn/.style={font=\scArial,color=scInk,
              execute at begin picture={\MTversion*{synarial}},
              >={Stealth[round,length=2.4mm,width=1.9mm]}},
  % containers ------------------------------------------------------------
  % One card chrome for the whole paper: white, 3 pt radius, one hairline.
  panel/.style={rounded corners=2.5pt,fill=none,draw=scRule,line width=0.45pt},
  card/.style={rounded corners=2.2pt,fill=scPaper,draw=scRule,line width=0.45pt,
               inner sep=5pt,align=left,anchor=north west},
  % a quiet strip for a panel's closing fact
  note/.style={rounded corners=2.5pt,fill=scPaper,draw=scRuleL,line width=0.4pt,
               inner xsep=6pt,inner ysep=3.5pt,align=left,font=\scSmall},
  % small labelled value pill
  chip/.style={rounded corners=2pt,fill=scPaper,draw=scRule,line width=0.4pt,
               inner xsep=4pt,inner ysep=2pt,font=\scSmall},
  % Circular letter badge: the same accessible panel marker used in Style/*.
  hbullet/.style={circle,fill=scInk,text=white,draw=none,
                  font=\scArial\fontsize{6}{6}\selectfont\bfseries,
                  minimum size=4.2mm,inner sep=0pt},
  ctag/.style={font=\scTitle,anchor=north west,text=scInk},
  % CRN primitives --------------------------------------------------------
  species/.style={rounded rectangle,fill=scSpF,draw=scSpE,line width=0.45pt,
                  inner xsep=2.4pt,inner ysep=1.6pt,font=\scSmall,
                  text=scInk,minimum height=3.6mm},
  cofactor/.style={species,fill=scRefF,draw=scRef},
  catalyst/.style={species,fill=scHiF,draw=scHiE},
  rxn/.style={rectangle,fill=scRxF,draw=scRxE,line width=0.45pt,
              inner sep=1.2pt,font=\scSmall,text=scInk,
              minimum width=3.0mm,minimum height=3.0mm},
  place/.style={circle,fill=scSpF,draw=scSpE,line width=0.45pt,
                minimum size=4.0mm,inner sep=0pt},
  trans/.style={rectangle,fill=scRxE,draw=scRxE,minimum width=1.1mm,
                minimum height=3.6mm,inner sep=0pt},
  token/.style={circle,fill=scInk,minimum size=1.1mm,inner sep=0pt},
  % edges -----------------------------------------------------------------
  reac/.style={->,line width=0.5pt,draw=scReac,shorten >=0.8pt,shorten <=0.8pt},
  prod/.style={->,line width=0.5pt,draw=scProd,shorten >=0.8pt,shorten <=0.8pt},
  cat/.style={->,line width=0.5pt,draw=scHiE,shorten >=0.8pt,shorten <=0.8pt},
  flow/.style={-{Stealth[round,length=2.2mm,width=1.8mm]},line width=1.1pt,
               draw=scInkM!45},
  flab/.style={font=\scSmall,text=scInkM,inner sep=2pt},
  stoi/.style={font=\scSmall,text=scInkM,inner sep=0.8pt},
  % verdict marks: shape carries the meaning, colour only reinforces it
  vpass/.style={circle,fill=scAccentF,draw=scAccent,text=scAccent,
                font=\fontsize{6}{6}\selectfont\bfseries,
                minimum size=3.2mm,inner sep=0pt},
  vfail/.style={circle,fill=scAlertF,draw=scAlert,text=scAlert,
                font=\fontsize{6}{6}\selectfont\bfseries,
                minimum size=3.2mm,inner sep=0pt},
  callout/.style={font=\scSmall,text=scInkM,align=left},
}

% \scHead[<badge colour>]{<anchor coord>}{<letter>}{<title>}{<subtitle or empty>}
% One header treatment for every panel in the paper.
\providecommand{\scHead}[5][scInk]{%
  \node[hbullet,fill=#1,anchor=north west] (scbadge#3) at ($#2+(0,-0.015)$)
    {\MakeUppercase{#3}};
  \node[ctag,anchor=west] at ($(scbadge#3.east)+(0.14,0)$) {#4};
  \ifx\relax#5\relax\else
    \node[anchor=north west,font=\scSmall,text=scInkM]
      at ($#2+(0.56,-0.245)$) {#5};
  \fi}

% --- shared pgfplots look (only if the figure loaded pgfplots) ------------
\ifdefined\pgfplotsset
  \pgfplotsset{compat=1.18}
  \pgfplotsset{
    synaxis/.style={
      scale only axis,
      tick align=outside, tick pos=left,
      axis line style={draw=scInkM!55,line width=0.4pt},
      % One grid, on the value axis only.  No minor grid: at these sizes it
      % turns the plot frame into a grey wash.
      ymajorgrids=true, xmajorgrids=false,
      major grid style={draw=scRuleL,line width=0.3pt},
      major tick length=1.6pt,
      label style={font=\scBody,color=scInk},
      tick label style={font=\scSmall,color=scInkM},
      legend style={font=\scSmall,draw=none,fill=none,inner sep=1pt,
                    row sep=0.4pt,legend columns=1},
      legend cell align=left,
      every axis plot/.append style={line width=0.7pt,mark size=1.2pt},
    },
  }
\fi
}{% ---------------------------------------------------------------------------
% synstyle.tex --- shared look for every TikZ figure of the SynKit CRN paper.
%
% Input it from a figure file:
%     \IfFileExists{synstyle.tex}{\input{synstyle}}{\input{Figure/synstyle}}
% so that the figure compiles both from Figure/ and from the paper root.
%
% The visual vocabulary follows ../../../../Style/style.tex: the
% colour-blind-safe Okabe--Ito accents, quiet neutral rules and wash, circular
% panel badges, and a compact technical type scale are shared by the companion
% papers.  Semantic aliases below keep the CRN meaning explicit.
%
% DESIGN CONTRACT -----------------------------------------------------------
% 1. Geometry.  Every figure is laid out on a canvas exactly \scW wide (the
%    sn-jnl text width, 372 pt) and is included unscaled, so one point of type
%    is the same physical size in all five figures.  Open each picture with
%    \scCanvas{<height in cm>} to pin the bounding box.
% 2. Colour is semantic, never decorative.  Four roles only:
%       scAccent (green)       structure, the normal/expected path
%       scMark   (orange)      catalyst, certificate, or highlight
%       scAlert  (vermillion)  failure, violation, or uncovered set
%       scRef    (blue)        secondary or reference series
%    A panel gets at most one accent.  Titles are always ink; panel identity
%    is also encoded by a letter inside a circular badge, not by colour alone.
% 3. Type scale.  Three sizes, no others, and nothing below 6 pt:
%       \scTitle (7 pt bold)  panel titles
%       \scBody  (6.5 pt)     labels, values, in-figure prose
%       \scSmall (6 pt)       units, footnotes, dense graph labels
% ---------------------------------------------------------------------------
% Load guard.  The paper inputs every figure source inline, so this file is
% reached once per figure; without the guard the second \newlength below
% would abort the run.
\ifdefined\SynKitStyleLoaded\endinput\fi
\newcommand{\SynKitStyleLoaded}{}

% Load in the document preamble. XeLaTeX/LuaLaTeX and an installed Arial
% family are required; fontspec reports an error instead of substituting.
\RequirePackage[no-math]{fontspec}
\newfontfamily\scArial{Arial}[NFSSFamily=scArial,Ligatures=TeX]
% Keep manuscript mathematics unchanged; select Arial letters and numerals
% only inside figures. TeX supplies specialist mathematical symbols.
\RequirePackage[subdued,italic]{mathastext}
\MTDeclareVersion{synarial}{TU}{scArial}{m}{n}
\SetMathAlphabet{\Mathsf}{synarial}{TU}{scArial}{m}{n}
\SetMathAlphabet{\Mathtt}{synarial}{TU}{scArial}{m}{n}

\usetikzlibrary{positioning,calc,fit,backgrounds,arrows.meta,shapes.geometric,
                shapes.misc,decorations.pathreplacing,patterns,math}

% --- canvas ---------------------------------------------------------------
% 372 pt text width minus the 1 pt standalone border on each side.
\newlength{\scW}\setlength{\scW}{13.0326cm}
\providecommand{\scWnum}{13.0326}
% \scCanvas{<height>} pins the picture to the full text width, y running down
% from 0, so every figure shares one coordinate convention.  It *sets* the
% bounding box rather than contributing to it, so the emitted page is exactly
% the text width and every figure is included at 1:1 -- otherwise a stray
% invisible extent (pgfplots reserves a little around an axis) rescales one
% figure relative to the others and the same \scSmall prints at two sizes.
% Anything drawn outside the canvas is therefore cropped: keep content inside,
% and check with
%     \providecommand{\scCanvas}[1]{\path[draw=red] (0,0) rectangle (\scWnum,-#1);}
\providecommand{\scCanvas}[1]{%
  \path[use as bounding box] (0,0) rectangle (\scWnum,-#1);}

% --- palette --------------------------------------------------------------
% Canonical companion-paper values (../../../../Style/style.tex).
\definecolor{npgOrange}{HTML}{E69F00}
\definecolor{npgSky}{HTML}{56B4E9}
\definecolor{npgGreen}{HTML}{009E73}
\definecolor{npgYellow}{HTML}{F0E442}
\definecolor{npgBlue}{HTML}{0072B2}
\definecolor{npgVermillion}{HTML}{E64B35}
\definecolor{npgPurple}{HTML}{CC79A7}
\definecolor{figInk}{HTML}{1A1A1A}
\definecolor{figMuted}{HTML}{666666}
\definecolor{figHair}{HTML}{D9D9D9}
\definecolor{figWash}{HTML}{F7F7F7}

\colorlet{scInk}{figInk}              % primary text
\colorlet{scInkM}{figMuted}           % secondary text
\colorlet{scPaper}{figWash}           % quiet cards and note strips
\colorlet{scSpF}{npgSky!12}           % species fill
\colorlet{scSpE}{npgBlue}             % species edge
\colorlet{scRxF}{npgOrange!12}        % reaction fill
\colorlet{scRxE}{npgOrange!88!black}  % reaction edge
\colorlet{scReac}{npgBlue!78!black}   % reactant edge
\colorlet{scProd}{npgGreen!82!black}  % product edge
\colorlet{scHiF}{npgOrange!13}        % highlight fill
\colorlet{scHiE}{npgOrange!82!black}  % highlight edge
% Stage / accent tints.  stDE remains available for the fifth plot series.
\colorlet{stAF}{npgGreen!9}\colorlet{stAE}{npgGreen}
\colorlet{stBF}{npgOrange!11}\colorlet{stBE}{npgOrange}
\colorlet{stCF}{npgBlue!9}\colorlet{stCE}{npgBlue}
\colorlet{stDF}{npgPurple!10}\colorlet{stDE}{npgPurple!85!black}
\colorlet{stEF}{npgVermillion!9}\colorlet{stEE}{npgVermillion}

% Semantic aliases.  Figures should speak in these, not in stAE/stBE/...
\colorlet{scAccent}{stAE}\colorlet{scAccentF}{stAF}   % structure / normal
\colorlet{scMark}{stBE}\colorlet{scMarkF}{stBF}       % highlight / certificate
\colorlet{scRef}{stCE}\colorlet{scRefF}{stCF}         % secondary series
\colorlet{scAlert}{stEE}\colorlet{scAlertF}{stEF}     % failure / violation
\colorlet{scRule}{figHair}                            % card and axis hairlines
\colorlet{scRuleL}{figHair!55}                        % grid, dividers

% --- text styles ----------------------------------------------------------
% Three sizes, fixed in points so they do not drift with the document class.
% Arial is shared by labels, headings, panel badges, and chemistry labels.
\providecommand{\scTitle}{\scArial\fontsize{7}{8.4}\selectfont\bfseries\color{scInk}}
\providecommand{\scBody}{\scArial\fontsize{6.5}{7.8}\selectfont\color{scInkM}}
\providecommand{\scSmall}{\scArial\fontsize{6}{7.2}\selectfont\color{scInkM}}
\providecommand{\scSub}[1]{{\scSmall\itshape #1}}
\providecommand{\scPanel}[1]{{\scTitle(#1)}}
% Numerals that carry a result: ink, not colour, so the eye reads the value.
\providecommand{\scNum}{\scArial\fontsize{10}{11}\selectfont\bfseries\color{scInk}}

% --- shared TikZ styles ---------------------------------------------------
\tikzset{
  % text never hyphenates inside a figure
  every node/.append style={execute at begin node={%
      \hyphenpenalty=10000 \exhyphenpenalty=10000 \hbadness=10000 }},
  syn/.style={font=\scArial,color=scInk,
              execute at begin picture={\MTversion*{synarial}},
              >={Stealth[round,length=2.4mm,width=1.9mm]}},
  % containers ------------------------------------------------------------
  % One card chrome for the whole paper: white, 3 pt radius, one hairline.
  panel/.style={rounded corners=2.5pt,fill=none,draw=scRule,line width=0.45pt},
  card/.style={rounded corners=2.2pt,fill=scPaper,draw=scRule,line width=0.45pt,
               inner sep=5pt,align=left,anchor=north west},
  % a quiet strip for a panel's closing fact
  note/.style={rounded corners=2.5pt,fill=scPaper,draw=scRuleL,line width=0.4pt,
               inner xsep=6pt,inner ysep=3.5pt,align=left,font=\scSmall},
  % small labelled value pill
  chip/.style={rounded corners=2pt,fill=scPaper,draw=scRule,line width=0.4pt,
               inner xsep=4pt,inner ysep=2pt,font=\scSmall},
  % Circular letter badge: the same accessible panel marker used in Style/*.
  hbullet/.style={circle,fill=scInk,text=white,draw=none,
                  font=\scArial\fontsize{6}{6}\selectfont\bfseries,
                  minimum size=4.2mm,inner sep=0pt},
  ctag/.style={font=\scTitle,anchor=north west,text=scInk},
  % CRN primitives --------------------------------------------------------
  species/.style={rounded rectangle,fill=scSpF,draw=scSpE,line width=0.45pt,
                  inner xsep=2.4pt,inner ysep=1.6pt,font=\scSmall,
                  text=scInk,minimum height=3.6mm},
  cofactor/.style={species,fill=scRefF,draw=scRef},
  catalyst/.style={species,fill=scHiF,draw=scHiE},
  rxn/.style={rectangle,fill=scRxF,draw=scRxE,line width=0.45pt,
              inner sep=1.2pt,font=\scSmall,text=scInk,
              minimum width=3.0mm,minimum height=3.0mm},
  place/.style={circle,fill=scSpF,draw=scSpE,line width=0.45pt,
                minimum size=4.0mm,inner sep=0pt},
  trans/.style={rectangle,fill=scRxE,draw=scRxE,minimum width=1.1mm,
                minimum height=3.6mm,inner sep=0pt},
  token/.style={circle,fill=scInk,minimum size=1.1mm,inner sep=0pt},
  % edges -----------------------------------------------------------------
  reac/.style={->,line width=0.5pt,draw=scReac,shorten >=0.8pt,shorten <=0.8pt},
  prod/.style={->,line width=0.5pt,draw=scProd,shorten >=0.8pt,shorten <=0.8pt},
  cat/.style={->,line width=0.5pt,draw=scHiE,shorten >=0.8pt,shorten <=0.8pt},
  flow/.style={-{Stealth[round,length=2.2mm,width=1.8mm]},line width=1.1pt,
               draw=scInkM!45},
  flab/.style={font=\scSmall,text=scInkM,inner sep=2pt},
  stoi/.style={font=\scSmall,text=scInkM,inner sep=0.8pt},
  % verdict marks: shape carries the meaning, colour only reinforces it
  vpass/.style={circle,fill=scAccentF,draw=scAccent,text=scAccent,
                font=\fontsize{6}{6}\selectfont\bfseries,
                minimum size=3.2mm,inner sep=0pt},
  vfail/.style={circle,fill=scAlertF,draw=scAlert,text=scAlert,
                font=\fontsize{6}{6}\selectfont\bfseries,
                minimum size=3.2mm,inner sep=0pt},
  callout/.style={font=\scSmall,text=scInkM,align=left},
}

% \scHead[<badge colour>]{<anchor coord>}{<letter>}{<title>}{<subtitle or empty>}
% One header treatment for every panel in the paper.
\providecommand{\scHead}[5][scInk]{%
  \node[hbullet,fill=#1,anchor=north west] (scbadge#3) at ($#2+(0,-0.015)$)
    {\MakeUppercase{#3}};
  \node[ctag,anchor=west] at ($(scbadge#3.east)+(0.14,0)$) {#4};
  \ifx\relax#5\relax\else
    \node[anchor=north west,font=\scSmall,text=scInkM]
      at ($#2+(0.56,-0.245)$) {#5};
  \fi}

% --- shared pgfplots look (only if the figure loaded pgfplots) ------------
\ifdefined\pgfplotsset
  \pgfplotsset{compat=1.18}
  \pgfplotsset{
    synaxis/.style={
      scale only axis,
      tick align=outside, tick pos=left,
      axis line style={draw=scInkM!55,line width=0.4pt},
      % One grid, on the value axis only.  No minor grid: at these sizes it
      % turns the plot frame into a grey wash.
      ymajorgrids=true, xmajorgrids=false,
      major grid style={draw=scRuleL,line width=0.3pt},
      major tick length=1.6pt,
      label style={font=\scBody,color=scInk},
      tick label style={font=\scSmall,color=scInkM},
      legend style={font=\scSmall,draw=none,fill=none,inner sep=1pt,
                    row sep=0.4pt,legend columns=1},
      legend cell align=left,
      every axis plot/.append style={line width=0.7pt,mark size=1.2pt},
    },
  }
\fi
}
% so that the figure compiles both from Figure/ and from the paper root.
%
% The visual vocabulary follows ../../../../Style/style.tex: the
% colour-blind-safe Okabe--Ito accents, quiet neutral rules and wash, circular
% panel badges, and a compact technical type scale are shared by the companion
% papers.  Semantic aliases below keep the CRN meaning explicit.
%
% DESIGN CONTRACT -----------------------------------------------------------
% 1. Geometry.  Every figure is laid out on a canvas exactly \scW wide (the
%    sn-jnl text width, 372 pt) and is included unscaled, so one point of type
%    is the same physical size in all five figures.  Open each picture with
%    \scCanvas{<height in cm>} to pin the bounding box.
% 2. Colour is semantic, never decorative.  Four roles only:
%       scAccent (green)       structure, the normal/expected path
%       scMark   (orange)      catalyst, certificate, or highlight
%       scAlert  (vermillion)  failure, violation, or uncovered set
%       scRef    (blue)        secondary or reference series
%    A panel gets at most one accent.  Titles are always ink; panel identity
%    is also encoded by a letter inside a circular badge, not by colour alone.
% 3. Type scale.  Three sizes, no others, and nothing below 6 pt:
%       \scTitle (7 pt bold)  panel titles
%       \scBody  (6.5 pt)     labels, values, in-figure prose
%       \scSmall (6 pt)       units, footnotes, dense graph labels
% ---------------------------------------------------------------------------
% Load guard.  The paper inputs every figure source inline, so this file is
% reached once per figure; without the guard the second \newlength below
% would abort the run.
\ifdefined\SynKitStyleLoaded\endinput\fi
\newcommand{\SynKitStyleLoaded}{}

% Load in the document preamble. XeLaTeX/LuaLaTeX and an installed Arial
% family are required; fontspec reports an error instead of substituting.
\RequirePackage[no-math]{fontspec}
\newfontfamily\scArial{Arial}[NFSSFamily=scArial,Ligatures=TeX]
% Keep manuscript mathematics unchanged; select Arial letters and numerals
% only inside figures. TeX supplies specialist mathematical symbols.
\RequirePackage[subdued,italic]{mathastext}
\MTDeclareVersion{synarial}{TU}{scArial}{m}{n}
\SetMathAlphabet{\Mathsf}{synarial}{TU}{scArial}{m}{n}
\SetMathAlphabet{\Mathtt}{synarial}{TU}{scArial}{m}{n}

\usetikzlibrary{positioning,calc,fit,backgrounds,arrows.meta,shapes.geometric,
                shapes.misc,decorations.pathreplacing,patterns,math}

% --- canvas ---------------------------------------------------------------
% 372 pt text width minus the 1 pt standalone border on each side.
\newlength{\scW}\setlength{\scW}{13.0326cm}
\providecommand{\scWnum}{13.0326}
% \scCanvas{<height>} pins the picture to the full text width, y running down
% from 0, so every figure shares one coordinate convention.  It *sets* the
% bounding box rather than contributing to it, so the emitted page is exactly
% the text width and every figure is included at 1:1 -- otherwise a stray
% invisible extent (pgfplots reserves a little around an axis) rescales one
% figure relative to the others and the same \scSmall prints at two sizes.
% Anything drawn outside the canvas is therefore cropped: keep content inside,
% and check with
%     \providecommand{\scCanvas}[1]{\path[draw=red] (0,0) rectangle (\scWnum,-#1);}
\providecommand{\scCanvas}[1]{%
  \path[use as bounding box] (0,0) rectangle (\scWnum,-#1);}

% --- palette --------------------------------------------------------------
% Canonical companion-paper values (../../../../Style/style.tex).
\definecolor{npgOrange}{HTML}{E69F00}
\definecolor{npgSky}{HTML}{56B4E9}
\definecolor{npgGreen}{HTML}{009E73}
\definecolor{npgYellow}{HTML}{F0E442}
\definecolor{npgBlue}{HTML}{0072B2}
\definecolor{npgVermillion}{HTML}{E64B35}
\definecolor{npgPurple}{HTML}{CC79A7}
\definecolor{figInk}{HTML}{1A1A1A}
\definecolor{figMuted}{HTML}{666666}
\definecolor{figHair}{HTML}{D9D9D9}
\definecolor{figWash}{HTML}{F7F7F7}

\colorlet{scInk}{figInk}              % primary text
\colorlet{scInkM}{figMuted}           % secondary text
\colorlet{scPaper}{figWash}           % quiet cards and note strips
\colorlet{scSpF}{npgSky!12}           % species fill
\colorlet{scSpE}{npgBlue}             % species edge
\colorlet{scRxF}{npgOrange!12}        % reaction fill
\colorlet{scRxE}{npgOrange!88!black}  % reaction edge
\colorlet{scReac}{npgBlue!78!black}   % reactant edge
\colorlet{scProd}{npgGreen!82!black}  % product edge
\colorlet{scHiF}{npgOrange!13}        % highlight fill
\colorlet{scHiE}{npgOrange!82!black}  % highlight edge
% Stage / accent tints.  stDE remains available for the fifth plot series.
\colorlet{stAF}{npgGreen!9}\colorlet{stAE}{npgGreen}
\colorlet{stBF}{npgOrange!11}\colorlet{stBE}{npgOrange}
\colorlet{stCF}{npgBlue!9}\colorlet{stCE}{npgBlue}
\colorlet{stDF}{npgPurple!10}\colorlet{stDE}{npgPurple!85!black}
\colorlet{stEF}{npgVermillion!9}\colorlet{stEE}{npgVermillion}

% Semantic aliases.  Figures should speak in these, not in stAE/stBE/...
\colorlet{scAccent}{stAE}\colorlet{scAccentF}{stAF}   % structure / normal
\colorlet{scMark}{stBE}\colorlet{scMarkF}{stBF}       % highlight / certificate
\colorlet{scRef}{stCE}\colorlet{scRefF}{stCF}         % secondary series
\colorlet{scAlert}{stEE}\colorlet{scAlertF}{stEF}     % failure / violation
\colorlet{scRule}{figHair}                            % card and axis hairlines
\colorlet{scRuleL}{figHair!55}                        % grid, dividers

% --- text styles ----------------------------------------------------------
% Three sizes, fixed in points so they do not drift with the document class.
% Arial is shared by labels, headings, panel badges, and chemistry labels.
\providecommand{\scTitle}{\scArial\fontsize{7}{8.4}\selectfont\bfseries\color{scInk}}
\providecommand{\scBody}{\scArial\fontsize{6.5}{7.8}\selectfont\color{scInkM}}
\providecommand{\scSmall}{\scArial\fontsize{6}{7.2}\selectfont\color{scInkM}}
\providecommand{\scSub}[1]{{\scSmall\itshape #1}}
\providecommand{\scPanel}[1]{{\scTitle(#1)}}
% Numerals that carry a result: ink, not colour, so the eye reads the value.
\providecommand{\scNum}{\scArial\fontsize{10}{11}\selectfont\bfseries\color{scInk}}

% --- shared TikZ styles ---------------------------------------------------
\tikzset{
  % text never hyphenates inside a figure
  every node/.append style={execute at begin node={%
      \hyphenpenalty=10000 \exhyphenpenalty=10000 \hbadness=10000 }},
  syn/.style={font=\scArial,color=scInk,
              execute at begin picture={\MTversion*{synarial}},
              >={Stealth[round,length=2.4mm,width=1.9mm]}},
  % containers ------------------------------------------------------------
  % One card chrome for the whole paper: white, 3 pt radius, one hairline.
  panel/.style={rounded corners=2.5pt,fill=none,draw=scRule,line width=0.45pt},
  card/.style={rounded corners=2.2pt,fill=scPaper,draw=scRule,line width=0.45pt,
               inner sep=5pt,align=left,anchor=north west},
  % a quiet strip for a panel's closing fact
  note/.style={rounded corners=2.5pt,fill=scPaper,draw=scRuleL,line width=0.4pt,
               inner xsep=6pt,inner ysep=3.5pt,align=left,font=\scSmall},
  % small labelled value pill
  chip/.style={rounded corners=2pt,fill=scPaper,draw=scRule,line width=0.4pt,
               inner xsep=4pt,inner ysep=2pt,font=\scSmall},
  % Circular letter badge: the same accessible panel marker used in Style/*.
  hbullet/.style={circle,fill=scInk,text=white,draw=none,
                  font=\scArial\fontsize{6}{6}\selectfont\bfseries,
                  minimum size=4.2mm,inner sep=0pt},
  ctag/.style={font=\scTitle,anchor=north west,text=scInk},
  % CRN primitives --------------------------------------------------------
  species/.style={rounded rectangle,fill=scSpF,draw=scSpE,line width=0.45pt,
                  inner xsep=2.4pt,inner ysep=1.6pt,font=\scSmall,
                  text=scInk,minimum height=3.6mm},
  cofactor/.style={species,fill=scRefF,draw=scRef},
  catalyst/.style={species,fill=scHiF,draw=scHiE},
  rxn/.style={rectangle,fill=scRxF,draw=scRxE,line width=0.45pt,
              inner sep=1.2pt,font=\scSmall,text=scInk,
              minimum width=3.0mm,minimum height=3.0mm},
  place/.style={circle,fill=scSpF,draw=scSpE,line width=0.45pt,
                minimum size=4.0mm,inner sep=0pt},
  trans/.style={rectangle,fill=scRxE,draw=scRxE,minimum width=1.1mm,
                minimum height=3.6mm,inner sep=0pt},
  token/.style={circle,fill=scInk,minimum size=1.1mm,inner sep=0pt},
  % edges -----------------------------------------------------------------
  reac/.style={->,line width=0.5pt,draw=scReac,shorten >=0.8pt,shorten <=0.8pt},
  prod/.style={->,line width=0.5pt,draw=scProd,shorten >=0.8pt,shorten <=0.8pt},
  cat/.style={->,line width=0.5pt,draw=scHiE,shorten >=0.8pt,shorten <=0.8pt},
  flow/.style={-{Stealth[round,length=2.2mm,width=1.8mm]},line width=1.1pt,
               draw=scInkM!45},
  flab/.style={font=\scSmall,text=scInkM,inner sep=2pt},
  stoi/.style={font=\scSmall,text=scInkM,inner sep=0.8pt},
  % verdict marks: shape carries the meaning, colour only reinforces it
  vpass/.style={circle,fill=scAccentF,draw=scAccent,text=scAccent,
                font=\fontsize{6}{6}\selectfont\bfseries,
                minimum size=3.2mm,inner sep=0pt},
  vfail/.style={circle,fill=scAlertF,draw=scAlert,text=scAlert,
                font=\fontsize{6}{6}\selectfont\bfseries,
                minimum size=3.2mm,inner sep=0pt},
  callout/.style={font=\scSmall,text=scInkM,align=left},
}

% \scHead[<badge colour>]{<anchor coord>}{<letter>}{<title>}{<subtitle or empty>}
% One header treatment for every panel in the paper.
\providecommand{\scHead}[5][scInk]{%
  \node[hbullet,fill=#1,anchor=north west] (scbadge#3) at ($#2+(0,-0.015)$)
    {\MakeUppercase{#3}};
  \node[ctag,anchor=west] at ($(scbadge#3.east)+(0.14,0)$) {#4};
  \ifx\relax#5\relax\else
    \node[anchor=north west,font=\scSmall,text=scInkM]
      at ($#2+(0.56,-0.245)$) {#5};
  \fi}

% --- shared pgfplots look (only if the figure loaded pgfplots) ------------
\ifdefined\pgfplotsset
  \pgfplotsset{compat=1.18}
  \pgfplotsset{
    synaxis/.style={
      scale only axis,
      tick align=outside, tick pos=left,
      axis line style={draw=scInkM!55,line width=0.4pt},
      % One grid, on the value axis only.  No minor grid: at these sizes it
      % turns the plot frame into a grey wash.
      ymajorgrids=true, xmajorgrids=false,
      major grid style={draw=scRuleL,line width=0.3pt},
      major tick length=1.6pt,
      label style={font=\scBody,color=scInk},
      tick label style={font=\scSmall,color=scInkM},
      legend style={font=\scSmall,draw=none,fill=none,inner sep=1pt,
                    row sep=0.4pt,legend columns=1},
      legend cell align=left,
      every axis plot/.append style={line width=0.7pt,mark size=1.2pt},
    },
  }
\fi
}
% so that the figure compiles both from Figure/ and from the paper root.
%
% The visual vocabulary follows ../../../../Style/style.tex: the
% colour-blind-safe Okabe--Ito accents, quiet neutral rules and wash, circular
% panel badges, and a compact technical type scale are shared by the companion
% papers.  Semantic aliases below keep the CRN meaning explicit.
%
% DESIGN CONTRACT -----------------------------------------------------------
% 1. Geometry.  Every figure is laid out on a canvas exactly \scW wide (the
%    sn-jnl text width, 372 pt) and is included unscaled, so one point of type
%    is the same physical size in all five figures.  Open each picture with
%    \scCanvas{<height in cm>} to pin the bounding box.
% 2. Colour is semantic, never decorative.  Four roles only:
%       scAccent (green)       structure, the normal/expected path
%       scMark   (orange)      catalyst, certificate, or highlight
%       scAlert  (vermillion)  failure, violation, or uncovered set
%       scRef    (blue)        secondary or reference series
%    A panel gets at most one accent.  Titles are always ink; panel identity
%    is also encoded by a letter inside a circular badge, not by colour alone.
% 3. Type scale.  Three sizes, no others, and nothing below 6 pt:
%       \scTitle (7 pt bold)  panel titles
%       \scBody  (6.5 pt)     labels, values, in-figure prose
%       \scSmall (6 pt)       units, footnotes, dense graph labels
% ---------------------------------------------------------------------------
% Load guard.  The paper inputs every figure source inline, so this file is
% reached once per figure; without the guard the second \newlength below
% would abort the run.
\ifdefined\SynKitStyleLoaded\endinput\fi
\newcommand{\SynKitStyleLoaded}{}

% Load in the document preamble. XeLaTeX/LuaLaTeX and an installed Arial
% family are required; fontspec reports an error instead of substituting.
\RequirePackage[no-math]{fontspec}
\newfontfamily\scArial{Arial}[NFSSFamily=scArial,Ligatures=TeX]
% Keep manuscript mathematics unchanged; select Arial letters and numerals
% only inside figures. TeX supplies specialist mathematical symbols.
\RequirePackage[subdued,italic]{mathastext}
\MTDeclareVersion{synarial}{TU}{scArial}{m}{n}
\SetMathAlphabet{\Mathsf}{synarial}{TU}{scArial}{m}{n}
\SetMathAlphabet{\Mathtt}{synarial}{TU}{scArial}{m}{n}

\usetikzlibrary{positioning,calc,fit,backgrounds,arrows.meta,shapes.geometric,
                shapes.misc,decorations.pathreplacing,patterns,math}

% --- canvas ---------------------------------------------------------------
% 372 pt text width minus the 1 pt standalone border on each side.
\newlength{\scW}\setlength{\scW}{13.0326cm}
\providecommand{\scWnum}{13.0326}
% \scCanvas{<height>} pins the picture to the full text width, y running down
% from 0, so every figure shares one coordinate convention.  It *sets* the
% bounding box rather than contributing to it, so the emitted page is exactly
% the text width and every figure is included at 1:1 -- otherwise a stray
% invisible extent (pgfplots reserves a little around an axis) rescales one
% figure relative to the others and the same \scSmall prints at two sizes.
% Anything drawn outside the canvas is therefore cropped: keep content inside,
% and check with
%     \providecommand{\scCanvas}[1]{\path[draw=red] (0,0) rectangle (\scWnum,-#1);}
\providecommand{\scCanvas}[1]{%
  \path[use as bounding box] (0,0) rectangle (\scWnum,-#1);}

% --- palette --------------------------------------------------------------
% Canonical companion-paper values (../../../../Style/style.tex).
\definecolor{npgOrange}{HTML}{E69F00}
\definecolor{npgSky}{HTML}{56B4E9}
\definecolor{npgGreen}{HTML}{009E73}
\definecolor{npgYellow}{HTML}{F0E442}
\definecolor{npgBlue}{HTML}{0072B2}
\definecolor{npgVermillion}{HTML}{E64B35}
\definecolor{npgPurple}{HTML}{CC79A7}
\definecolor{figInk}{HTML}{1A1A1A}
\definecolor{figMuted}{HTML}{666666}
\definecolor{figHair}{HTML}{D9D9D9}
\definecolor{figWash}{HTML}{F7F7F7}

\colorlet{scInk}{figInk}              % primary text
\colorlet{scInkM}{figMuted}           % secondary text
\colorlet{scPaper}{figWash}           % quiet cards and note strips
\colorlet{scSpF}{npgSky!12}           % species fill
\colorlet{scSpE}{npgBlue}             % species edge
\colorlet{scRxF}{npgOrange!12}        % reaction fill
\colorlet{scRxE}{npgOrange!88!black}  % reaction edge
\colorlet{scReac}{npgBlue!78!black}   % reactant edge
\colorlet{scProd}{npgGreen!82!black}  % product edge
\colorlet{scHiF}{npgOrange!13}        % highlight fill
\colorlet{scHiE}{npgOrange!82!black}  % highlight edge
% Stage / accent tints.  stDE remains available for the fifth plot series.
\colorlet{stAF}{npgGreen!9}\colorlet{stAE}{npgGreen}
\colorlet{stBF}{npgOrange!11}\colorlet{stBE}{npgOrange}
\colorlet{stCF}{npgBlue!9}\colorlet{stCE}{npgBlue}
\colorlet{stDF}{npgPurple!10}\colorlet{stDE}{npgPurple!85!black}
\colorlet{stEF}{npgVermillion!9}\colorlet{stEE}{npgVermillion}

% Semantic aliases.  Figures should speak in these, not in stAE/stBE/...
\colorlet{scAccent}{stAE}\colorlet{scAccentF}{stAF}   % structure / normal
\colorlet{scMark}{stBE}\colorlet{scMarkF}{stBF}       % highlight / certificate
\colorlet{scRef}{stCE}\colorlet{scRefF}{stCF}         % secondary series
\colorlet{scAlert}{stEE}\colorlet{scAlertF}{stEF}     % failure / violation
\colorlet{scRule}{figHair}                            % card and axis hairlines
\colorlet{scRuleL}{figHair!55}                        % grid, dividers

% --- text styles ----------------------------------------------------------
% Three sizes, fixed in points so they do not drift with the document class.
% Arial is shared by labels, headings, panel badges, and chemistry labels.
\providecommand{\scTitle}{\scArial\fontsize{7}{8.4}\selectfont\bfseries\color{scInk}}
\providecommand{\scBody}{\scArial\fontsize{6.5}{7.8}\selectfont\color{scInkM}}
\providecommand{\scSmall}{\scArial\fontsize{6}{7.2}\selectfont\color{scInkM}}
\providecommand{\scSub}[1]{{\scSmall\itshape #1}}
\providecommand{\scPanel}[1]{{\scTitle(#1)}}
% Numerals that carry a result: ink, not colour, so the eye reads the value.
\providecommand{\scNum}{\scArial\fontsize{10}{11}\selectfont\bfseries\color{scInk}}

% --- shared TikZ styles ---------------------------------------------------
\tikzset{
  % text never hyphenates inside a figure
  every node/.append style={execute at begin node={%
      \hyphenpenalty=10000 \exhyphenpenalty=10000 \hbadness=10000 }},
  syn/.style={font=\scArial,color=scInk,
              execute at begin picture={\MTversion*{synarial}},
              >={Stealth[round,length=2.4mm,width=1.9mm]}},
  % containers ------------------------------------------------------------
  % One card chrome for the whole paper: white, 3 pt radius, one hairline.
  panel/.style={rounded corners=2.5pt,fill=none,draw=scRule,line width=0.45pt},
  card/.style={rounded corners=2.2pt,fill=scPaper,draw=scRule,line width=0.45pt,
               inner sep=5pt,align=left,anchor=north west},
  % a quiet strip for a panel's closing fact
  note/.style={rounded corners=2.5pt,fill=scPaper,draw=scRuleL,line width=0.4pt,
               inner xsep=6pt,inner ysep=3.5pt,align=left,font=\scSmall},
  % small labelled value pill
  chip/.style={rounded corners=2pt,fill=scPaper,draw=scRule,line width=0.4pt,
               inner xsep=4pt,inner ysep=2pt,font=\scSmall},
  % Circular letter badge: the same accessible panel marker used in Style/*.
  hbullet/.style={circle,fill=scInk,text=white,draw=none,
                  font=\scArial\fontsize{6}{6}\selectfont\bfseries,
                  minimum size=4.2mm,inner sep=0pt},
  ctag/.style={font=\scTitle,anchor=north west,text=scInk},
  % CRN primitives --------------------------------------------------------
  species/.style={rounded rectangle,fill=scSpF,draw=scSpE,line width=0.45pt,
                  inner xsep=2.4pt,inner ysep=1.6pt,font=\scSmall,
                  text=scInk,minimum height=3.6mm},
  cofactor/.style={species,fill=scRefF,draw=scRef},
  catalyst/.style={species,fill=scHiF,draw=scHiE},
  rxn/.style={rectangle,fill=scRxF,draw=scRxE,line width=0.45pt,
              inner sep=1.2pt,font=\scSmall,text=scInk,
              minimum width=3.0mm,minimum height=3.0mm},
  place/.style={circle,fill=scSpF,draw=scSpE,line width=0.45pt,
                minimum size=4.0mm,inner sep=0pt},
  trans/.style={rectangle,fill=scRxE,draw=scRxE,minimum width=1.1mm,
                minimum height=3.6mm,inner sep=0pt},
  token/.style={circle,fill=scInk,minimum size=1.1mm,inner sep=0pt},
  % edges -----------------------------------------------------------------
  reac/.style={->,line width=0.5pt,draw=scReac,shorten >=0.8pt,shorten <=0.8pt},
  prod/.style={->,line width=0.5pt,draw=scProd,shorten >=0.8pt,shorten <=0.8pt},
  cat/.style={->,line width=0.5pt,draw=scHiE,shorten >=0.8pt,shorten <=0.8pt},
  flow/.style={-{Stealth[round,length=2.2mm,width=1.8mm]},line width=1.1pt,
               draw=scInkM!45},
  flab/.style={font=\scSmall,text=scInkM,inner sep=2pt},
  stoi/.style={font=\scSmall,text=scInkM,inner sep=0.8pt},
  % verdict marks: shape carries the meaning, colour only reinforces it
  vpass/.style={circle,fill=scAccentF,draw=scAccent,text=scAccent,
                font=\fontsize{6}{6}\selectfont\bfseries,
                minimum size=3.2mm,inner sep=0pt},
  vfail/.style={circle,fill=scAlertF,draw=scAlert,text=scAlert,
                font=\fontsize{6}{6}\selectfont\bfseries,
                minimum size=3.2mm,inner sep=0pt},
  callout/.style={font=\scSmall,text=scInkM,align=left},
}

% \scHead[<badge colour>]{<anchor coord>}{<letter>}{<title>}{<subtitle or empty>}
% One header treatment for every panel in the paper.
\providecommand{\scHead}[5][scInk]{%
  \node[hbullet,fill=#1,anchor=north west] (scbadge#3) at ($#2+(0,-0.015)$)
    {\MakeUppercase{#3}};
  \node[ctag,anchor=west] at ($(scbadge#3.east)+(0.14,0)$) {#4};
  \ifx\relax#5\relax\else
    \node[anchor=north west,font=\scSmall,text=scInkM]
      at ($#2+(0.56,-0.245)$) {#5};
  \fi}

% --- shared pgfplots look (only if the figure loaded pgfplots) ------------
\ifdefined\pgfplotsset
  \pgfplotsset{compat=1.18}
  \pgfplotsset{
    synaxis/.style={
      scale only axis,
      tick align=outside, tick pos=left,
      axis line style={draw=scInkM!55,line width=0.4pt},
      % One grid, on the value axis only.  No minor grid: at these sizes it
      % turns the plot frame into a grey wash.
      ymajorgrids=true, xmajorgrids=false,
      major grid style={draw=scRuleL,line width=0.3pt},
      major tick length=1.6pt,
      label style={font=\scBody,color=scInk},
      tick label style={font=\scSmall,color=scInkM},
      legend style={font=\scSmall,draw=none,fill=none,inner sep=1pt,
                    row sep=0.4pt,legend columns=1},
      legend cell align=left,
      every axis plot/.append style={line width=0.7pt,mark size=1.2pt},
    },
  }
\fi
}{%
  \IfFileExists{Figure/synstyle.tex}{% ---------------------------------------------------------------------------
% synstyle.tex --- shared look for every TikZ figure of the SynKit CRN paper.
%
% Input it from a figure file:
%     \IfFileExists{synstyle.tex}{% ---------------------------------------------------------------------------
% synstyle.tex --- shared look for every TikZ figure of the SynKit CRN paper.
%
% Input it from a figure file:
%     \IfFileExists{synstyle.tex}{% ---------------------------------------------------------------------------
% synstyle.tex --- shared look for every TikZ figure of the SynKit CRN paper.
%
% Input it from a figure file:
%     \IfFileExists{synstyle.tex}{\input{synstyle}}{\input{Figure/synstyle}}
% so that the figure compiles both from Figure/ and from the paper root.
%
% The visual vocabulary follows ../../../../Style/style.tex: the
% colour-blind-safe Okabe--Ito accents, quiet neutral rules and wash, circular
% panel badges, and a compact technical type scale are shared by the companion
% papers.  Semantic aliases below keep the CRN meaning explicit.
%
% DESIGN CONTRACT -----------------------------------------------------------
% 1. Geometry.  Every figure is laid out on a canvas exactly \scW wide (the
%    sn-jnl text width, 372 pt) and is included unscaled, so one point of type
%    is the same physical size in all five figures.  Open each picture with
%    \scCanvas{<height in cm>} to pin the bounding box.
% 2. Colour is semantic, never decorative.  Four roles only:
%       scAccent (green)       structure, the normal/expected path
%       scMark   (orange)      catalyst, certificate, or highlight
%       scAlert  (vermillion)  failure, violation, or uncovered set
%       scRef    (blue)        secondary or reference series
%    A panel gets at most one accent.  Titles are always ink; panel identity
%    is also encoded by a letter inside a circular badge, not by colour alone.
% 3. Type scale.  Three sizes, no others, and nothing below 6 pt:
%       \scTitle (7 pt bold)  panel titles
%       \scBody  (6.5 pt)     labels, values, in-figure prose
%       \scSmall (6 pt)       units, footnotes, dense graph labels
% ---------------------------------------------------------------------------
% Load guard.  The paper inputs every figure source inline, so this file is
% reached once per figure; without the guard the second \newlength below
% would abort the run.
\ifdefined\SynKitStyleLoaded\endinput\fi
\newcommand{\SynKitStyleLoaded}{}

% Load in the document preamble. XeLaTeX/LuaLaTeX and an installed Arial
% family are required; fontspec reports an error instead of substituting.
\RequirePackage[no-math]{fontspec}
\newfontfamily\scArial{Arial}[NFSSFamily=scArial,Ligatures=TeX]
% Keep manuscript mathematics unchanged; select Arial letters and numerals
% only inside figures. TeX supplies specialist mathematical symbols.
\RequirePackage[subdued,italic]{mathastext}
\MTDeclareVersion{synarial}{TU}{scArial}{m}{n}
\SetMathAlphabet{\Mathsf}{synarial}{TU}{scArial}{m}{n}
\SetMathAlphabet{\Mathtt}{synarial}{TU}{scArial}{m}{n}

\usetikzlibrary{positioning,calc,fit,backgrounds,arrows.meta,shapes.geometric,
                shapes.misc,decorations.pathreplacing,patterns,math}

% --- canvas ---------------------------------------------------------------
% 372 pt text width minus the 1 pt standalone border on each side.
\newlength{\scW}\setlength{\scW}{13.0326cm}
\providecommand{\scWnum}{13.0326}
% \scCanvas{<height>} pins the picture to the full text width, y running down
% from 0, so every figure shares one coordinate convention.  It *sets* the
% bounding box rather than contributing to it, so the emitted page is exactly
% the text width and every figure is included at 1:1 -- otherwise a stray
% invisible extent (pgfplots reserves a little around an axis) rescales one
% figure relative to the others and the same \scSmall prints at two sizes.
% Anything drawn outside the canvas is therefore cropped: keep content inside,
% and check with
%     \providecommand{\scCanvas}[1]{\path[draw=red] (0,0) rectangle (\scWnum,-#1);}
\providecommand{\scCanvas}[1]{%
  \path[use as bounding box] (0,0) rectangle (\scWnum,-#1);}

% --- palette --------------------------------------------------------------
% Canonical companion-paper values (../../../../Style/style.tex).
\definecolor{npgOrange}{HTML}{E69F00}
\definecolor{npgSky}{HTML}{56B4E9}
\definecolor{npgGreen}{HTML}{009E73}
\definecolor{npgYellow}{HTML}{F0E442}
\definecolor{npgBlue}{HTML}{0072B2}
\definecolor{npgVermillion}{HTML}{E64B35}
\definecolor{npgPurple}{HTML}{CC79A7}
\definecolor{figInk}{HTML}{1A1A1A}
\definecolor{figMuted}{HTML}{666666}
\definecolor{figHair}{HTML}{D9D9D9}
\definecolor{figWash}{HTML}{F7F7F7}

\colorlet{scInk}{figInk}              % primary text
\colorlet{scInkM}{figMuted}           % secondary text
\colorlet{scPaper}{figWash}           % quiet cards and note strips
\colorlet{scSpF}{npgSky!12}           % species fill
\colorlet{scSpE}{npgBlue}             % species edge
\colorlet{scRxF}{npgOrange!12}        % reaction fill
\colorlet{scRxE}{npgOrange!88!black}  % reaction edge
\colorlet{scReac}{npgBlue!78!black}   % reactant edge
\colorlet{scProd}{npgGreen!82!black}  % product edge
\colorlet{scHiF}{npgOrange!13}        % highlight fill
\colorlet{scHiE}{npgOrange!82!black}  % highlight edge
% Stage / accent tints.  stDE remains available for the fifth plot series.
\colorlet{stAF}{npgGreen!9}\colorlet{stAE}{npgGreen}
\colorlet{stBF}{npgOrange!11}\colorlet{stBE}{npgOrange}
\colorlet{stCF}{npgBlue!9}\colorlet{stCE}{npgBlue}
\colorlet{stDF}{npgPurple!10}\colorlet{stDE}{npgPurple!85!black}
\colorlet{stEF}{npgVermillion!9}\colorlet{stEE}{npgVermillion}

% Semantic aliases.  Figures should speak in these, not in stAE/stBE/...
\colorlet{scAccent}{stAE}\colorlet{scAccentF}{stAF}   % structure / normal
\colorlet{scMark}{stBE}\colorlet{scMarkF}{stBF}       % highlight / certificate
\colorlet{scRef}{stCE}\colorlet{scRefF}{stCF}         % secondary series
\colorlet{scAlert}{stEE}\colorlet{scAlertF}{stEF}     % failure / violation
\colorlet{scRule}{figHair}                            % card and axis hairlines
\colorlet{scRuleL}{figHair!55}                        % grid, dividers

% --- text styles ----------------------------------------------------------
% Three sizes, fixed in points so they do not drift with the document class.
% Arial is shared by labels, headings, panel badges, and chemistry labels.
\providecommand{\scTitle}{\scArial\fontsize{7}{8.4}\selectfont\bfseries\color{scInk}}
\providecommand{\scBody}{\scArial\fontsize{6.5}{7.8}\selectfont\color{scInkM}}
\providecommand{\scSmall}{\scArial\fontsize{6}{7.2}\selectfont\color{scInkM}}
\providecommand{\scSub}[1]{{\scSmall\itshape #1}}
\providecommand{\scPanel}[1]{{\scTitle(#1)}}
% Numerals that carry a result: ink, not colour, so the eye reads the value.
\providecommand{\scNum}{\scArial\fontsize{10}{11}\selectfont\bfseries\color{scInk}}

% --- shared TikZ styles ---------------------------------------------------
\tikzset{
  % text never hyphenates inside a figure
  every node/.append style={execute at begin node={%
      \hyphenpenalty=10000 \exhyphenpenalty=10000 \hbadness=10000 }},
  syn/.style={font=\scArial,color=scInk,
              execute at begin picture={\MTversion*{synarial}},
              >={Stealth[round,length=2.4mm,width=1.9mm]}},
  % containers ------------------------------------------------------------
  % One card chrome for the whole paper: white, 3 pt radius, one hairline.
  panel/.style={rounded corners=2.5pt,fill=none,draw=scRule,line width=0.45pt},
  card/.style={rounded corners=2.2pt,fill=scPaper,draw=scRule,line width=0.45pt,
               inner sep=5pt,align=left,anchor=north west},
  % a quiet strip for a panel's closing fact
  note/.style={rounded corners=2.5pt,fill=scPaper,draw=scRuleL,line width=0.4pt,
               inner xsep=6pt,inner ysep=3.5pt,align=left,font=\scSmall},
  % small labelled value pill
  chip/.style={rounded corners=2pt,fill=scPaper,draw=scRule,line width=0.4pt,
               inner xsep=4pt,inner ysep=2pt,font=\scSmall},
  % Circular letter badge: the same accessible panel marker used in Style/*.
  hbullet/.style={circle,fill=scInk,text=white,draw=none,
                  font=\scArial\fontsize{6}{6}\selectfont\bfseries,
                  minimum size=4.2mm,inner sep=0pt},
  ctag/.style={font=\scTitle,anchor=north west,text=scInk},
  % CRN primitives --------------------------------------------------------
  species/.style={rounded rectangle,fill=scSpF,draw=scSpE,line width=0.45pt,
                  inner xsep=2.4pt,inner ysep=1.6pt,font=\scSmall,
                  text=scInk,minimum height=3.6mm},
  cofactor/.style={species,fill=scRefF,draw=scRef},
  catalyst/.style={species,fill=scHiF,draw=scHiE},
  rxn/.style={rectangle,fill=scRxF,draw=scRxE,line width=0.45pt,
              inner sep=1.2pt,font=\scSmall,text=scInk,
              minimum width=3.0mm,minimum height=3.0mm},
  place/.style={circle,fill=scSpF,draw=scSpE,line width=0.45pt,
                minimum size=4.0mm,inner sep=0pt},
  trans/.style={rectangle,fill=scRxE,draw=scRxE,minimum width=1.1mm,
                minimum height=3.6mm,inner sep=0pt},
  token/.style={circle,fill=scInk,minimum size=1.1mm,inner sep=0pt},
  % edges -----------------------------------------------------------------
  reac/.style={->,line width=0.5pt,draw=scReac,shorten >=0.8pt,shorten <=0.8pt},
  prod/.style={->,line width=0.5pt,draw=scProd,shorten >=0.8pt,shorten <=0.8pt},
  cat/.style={->,line width=0.5pt,draw=scHiE,shorten >=0.8pt,shorten <=0.8pt},
  flow/.style={-{Stealth[round,length=2.2mm,width=1.8mm]},line width=1.1pt,
               draw=scInkM!45},
  flab/.style={font=\scSmall,text=scInkM,inner sep=2pt},
  stoi/.style={font=\scSmall,text=scInkM,inner sep=0.8pt},
  % verdict marks: shape carries the meaning, colour only reinforces it
  vpass/.style={circle,fill=scAccentF,draw=scAccent,text=scAccent,
                font=\fontsize{6}{6}\selectfont\bfseries,
                minimum size=3.2mm,inner sep=0pt},
  vfail/.style={circle,fill=scAlertF,draw=scAlert,text=scAlert,
                font=\fontsize{6}{6}\selectfont\bfseries,
                minimum size=3.2mm,inner sep=0pt},
  callout/.style={font=\scSmall,text=scInkM,align=left},
}

% \scHead[<badge colour>]{<anchor coord>}{<letter>}{<title>}{<subtitle or empty>}
% One header treatment for every panel in the paper.
\providecommand{\scHead}[5][scInk]{%
  \node[hbullet,fill=#1,anchor=north west] (scbadge#3) at ($#2+(0,-0.015)$)
    {\MakeUppercase{#3}};
  \node[ctag,anchor=west] at ($(scbadge#3.east)+(0.14,0)$) {#4};
  \ifx\relax#5\relax\else
    \node[anchor=north west,font=\scSmall,text=scInkM]
      at ($#2+(0.56,-0.245)$) {#5};
  \fi}

% --- shared pgfplots look (only if the figure loaded pgfplots) ------------
\ifdefined\pgfplotsset
  \pgfplotsset{compat=1.18}
  \pgfplotsset{
    synaxis/.style={
      scale only axis,
      tick align=outside, tick pos=left,
      axis line style={draw=scInkM!55,line width=0.4pt},
      % One grid, on the value axis only.  No minor grid: at these sizes it
      % turns the plot frame into a grey wash.
      ymajorgrids=true, xmajorgrids=false,
      major grid style={draw=scRuleL,line width=0.3pt},
      major tick length=1.6pt,
      label style={font=\scBody,color=scInk},
      tick label style={font=\scSmall,color=scInkM},
      legend style={font=\scSmall,draw=none,fill=none,inner sep=1pt,
                    row sep=0.4pt,legend columns=1},
      legend cell align=left,
      every axis plot/.append style={line width=0.7pt,mark size=1.2pt},
    },
  }
\fi
}{% ---------------------------------------------------------------------------
% synstyle.tex --- shared look for every TikZ figure of the SynKit CRN paper.
%
% Input it from a figure file:
%     \IfFileExists{synstyle.tex}{\input{synstyle}}{\input{Figure/synstyle}}
% so that the figure compiles both from Figure/ and from the paper root.
%
% The visual vocabulary follows ../../../../Style/style.tex: the
% colour-blind-safe Okabe--Ito accents, quiet neutral rules and wash, circular
% panel badges, and a compact technical type scale are shared by the companion
% papers.  Semantic aliases below keep the CRN meaning explicit.
%
% DESIGN CONTRACT -----------------------------------------------------------
% 1. Geometry.  Every figure is laid out on a canvas exactly \scW wide (the
%    sn-jnl text width, 372 pt) and is included unscaled, so one point of type
%    is the same physical size in all five figures.  Open each picture with
%    \scCanvas{<height in cm>} to pin the bounding box.
% 2. Colour is semantic, never decorative.  Four roles only:
%       scAccent (green)       structure, the normal/expected path
%       scMark   (orange)      catalyst, certificate, or highlight
%       scAlert  (vermillion)  failure, violation, or uncovered set
%       scRef    (blue)        secondary or reference series
%    A panel gets at most one accent.  Titles are always ink; panel identity
%    is also encoded by a letter inside a circular badge, not by colour alone.
% 3. Type scale.  Three sizes, no others, and nothing below 6 pt:
%       \scTitle (7 pt bold)  panel titles
%       \scBody  (6.5 pt)     labels, values, in-figure prose
%       \scSmall (6 pt)       units, footnotes, dense graph labels
% ---------------------------------------------------------------------------
% Load guard.  The paper inputs every figure source inline, so this file is
% reached once per figure; without the guard the second \newlength below
% would abort the run.
\ifdefined\SynKitStyleLoaded\endinput\fi
\newcommand{\SynKitStyleLoaded}{}

% Load in the document preamble. XeLaTeX/LuaLaTeX and an installed Arial
% family are required; fontspec reports an error instead of substituting.
\RequirePackage[no-math]{fontspec}
\newfontfamily\scArial{Arial}[NFSSFamily=scArial,Ligatures=TeX]
% Keep manuscript mathematics unchanged; select Arial letters and numerals
% only inside figures. TeX supplies specialist mathematical symbols.
\RequirePackage[subdued,italic]{mathastext}
\MTDeclareVersion{synarial}{TU}{scArial}{m}{n}
\SetMathAlphabet{\Mathsf}{synarial}{TU}{scArial}{m}{n}
\SetMathAlphabet{\Mathtt}{synarial}{TU}{scArial}{m}{n}

\usetikzlibrary{positioning,calc,fit,backgrounds,arrows.meta,shapes.geometric,
                shapes.misc,decorations.pathreplacing,patterns,math}

% --- canvas ---------------------------------------------------------------
% 372 pt text width minus the 1 pt standalone border on each side.
\newlength{\scW}\setlength{\scW}{13.0326cm}
\providecommand{\scWnum}{13.0326}
% \scCanvas{<height>} pins the picture to the full text width, y running down
% from 0, so every figure shares one coordinate convention.  It *sets* the
% bounding box rather than contributing to it, so the emitted page is exactly
% the text width and every figure is included at 1:1 -- otherwise a stray
% invisible extent (pgfplots reserves a little around an axis) rescales one
% figure relative to the others and the same \scSmall prints at two sizes.
% Anything drawn outside the canvas is therefore cropped: keep content inside,
% and check with
%     \providecommand{\scCanvas}[1]{\path[draw=red] (0,0) rectangle (\scWnum,-#1);}
\providecommand{\scCanvas}[1]{%
  \path[use as bounding box] (0,0) rectangle (\scWnum,-#1);}

% --- palette --------------------------------------------------------------
% Canonical companion-paper values (../../../../Style/style.tex).
\definecolor{npgOrange}{HTML}{E69F00}
\definecolor{npgSky}{HTML}{56B4E9}
\definecolor{npgGreen}{HTML}{009E73}
\definecolor{npgYellow}{HTML}{F0E442}
\definecolor{npgBlue}{HTML}{0072B2}
\definecolor{npgVermillion}{HTML}{E64B35}
\definecolor{npgPurple}{HTML}{CC79A7}
\definecolor{figInk}{HTML}{1A1A1A}
\definecolor{figMuted}{HTML}{666666}
\definecolor{figHair}{HTML}{D9D9D9}
\definecolor{figWash}{HTML}{F7F7F7}

\colorlet{scInk}{figInk}              % primary text
\colorlet{scInkM}{figMuted}           % secondary text
\colorlet{scPaper}{figWash}           % quiet cards and note strips
\colorlet{scSpF}{npgSky!12}           % species fill
\colorlet{scSpE}{npgBlue}             % species edge
\colorlet{scRxF}{npgOrange!12}        % reaction fill
\colorlet{scRxE}{npgOrange!88!black}  % reaction edge
\colorlet{scReac}{npgBlue!78!black}   % reactant edge
\colorlet{scProd}{npgGreen!82!black}  % product edge
\colorlet{scHiF}{npgOrange!13}        % highlight fill
\colorlet{scHiE}{npgOrange!82!black}  % highlight edge
% Stage / accent tints.  stDE remains available for the fifth plot series.
\colorlet{stAF}{npgGreen!9}\colorlet{stAE}{npgGreen}
\colorlet{stBF}{npgOrange!11}\colorlet{stBE}{npgOrange}
\colorlet{stCF}{npgBlue!9}\colorlet{stCE}{npgBlue}
\colorlet{stDF}{npgPurple!10}\colorlet{stDE}{npgPurple!85!black}
\colorlet{stEF}{npgVermillion!9}\colorlet{stEE}{npgVermillion}

% Semantic aliases.  Figures should speak in these, not in stAE/stBE/...
\colorlet{scAccent}{stAE}\colorlet{scAccentF}{stAF}   % structure / normal
\colorlet{scMark}{stBE}\colorlet{scMarkF}{stBF}       % highlight / certificate
\colorlet{scRef}{stCE}\colorlet{scRefF}{stCF}         % secondary series
\colorlet{scAlert}{stEE}\colorlet{scAlertF}{stEF}     % failure / violation
\colorlet{scRule}{figHair}                            % card and axis hairlines
\colorlet{scRuleL}{figHair!55}                        % grid, dividers

% --- text styles ----------------------------------------------------------
% Three sizes, fixed in points so they do not drift with the document class.
% Arial is shared by labels, headings, panel badges, and chemistry labels.
\providecommand{\scTitle}{\scArial\fontsize{7}{8.4}\selectfont\bfseries\color{scInk}}
\providecommand{\scBody}{\scArial\fontsize{6.5}{7.8}\selectfont\color{scInkM}}
\providecommand{\scSmall}{\scArial\fontsize{6}{7.2}\selectfont\color{scInkM}}
\providecommand{\scSub}[1]{{\scSmall\itshape #1}}
\providecommand{\scPanel}[1]{{\scTitle(#1)}}
% Numerals that carry a result: ink, not colour, so the eye reads the value.
\providecommand{\scNum}{\scArial\fontsize{10}{11}\selectfont\bfseries\color{scInk}}

% --- shared TikZ styles ---------------------------------------------------
\tikzset{
  % text never hyphenates inside a figure
  every node/.append style={execute at begin node={%
      \hyphenpenalty=10000 \exhyphenpenalty=10000 \hbadness=10000 }},
  syn/.style={font=\scArial,color=scInk,
              execute at begin picture={\MTversion*{synarial}},
              >={Stealth[round,length=2.4mm,width=1.9mm]}},
  % containers ------------------------------------------------------------
  % One card chrome for the whole paper: white, 3 pt radius, one hairline.
  panel/.style={rounded corners=2.5pt,fill=none,draw=scRule,line width=0.45pt},
  card/.style={rounded corners=2.2pt,fill=scPaper,draw=scRule,line width=0.45pt,
               inner sep=5pt,align=left,anchor=north west},
  % a quiet strip for a panel's closing fact
  note/.style={rounded corners=2.5pt,fill=scPaper,draw=scRuleL,line width=0.4pt,
               inner xsep=6pt,inner ysep=3.5pt,align=left,font=\scSmall},
  % small labelled value pill
  chip/.style={rounded corners=2pt,fill=scPaper,draw=scRule,line width=0.4pt,
               inner xsep=4pt,inner ysep=2pt,font=\scSmall},
  % Circular letter badge: the same accessible panel marker used in Style/*.
  hbullet/.style={circle,fill=scInk,text=white,draw=none,
                  font=\scArial\fontsize{6}{6}\selectfont\bfseries,
                  minimum size=4.2mm,inner sep=0pt},
  ctag/.style={font=\scTitle,anchor=north west,text=scInk},
  % CRN primitives --------------------------------------------------------
  species/.style={rounded rectangle,fill=scSpF,draw=scSpE,line width=0.45pt,
                  inner xsep=2.4pt,inner ysep=1.6pt,font=\scSmall,
                  text=scInk,minimum height=3.6mm},
  cofactor/.style={species,fill=scRefF,draw=scRef},
  catalyst/.style={species,fill=scHiF,draw=scHiE},
  rxn/.style={rectangle,fill=scRxF,draw=scRxE,line width=0.45pt,
              inner sep=1.2pt,font=\scSmall,text=scInk,
              minimum width=3.0mm,minimum height=3.0mm},
  place/.style={circle,fill=scSpF,draw=scSpE,line width=0.45pt,
                minimum size=4.0mm,inner sep=0pt},
  trans/.style={rectangle,fill=scRxE,draw=scRxE,minimum width=1.1mm,
                minimum height=3.6mm,inner sep=0pt},
  token/.style={circle,fill=scInk,minimum size=1.1mm,inner sep=0pt},
  % edges -----------------------------------------------------------------
  reac/.style={->,line width=0.5pt,draw=scReac,shorten >=0.8pt,shorten <=0.8pt},
  prod/.style={->,line width=0.5pt,draw=scProd,shorten >=0.8pt,shorten <=0.8pt},
  cat/.style={->,line width=0.5pt,draw=scHiE,shorten >=0.8pt,shorten <=0.8pt},
  flow/.style={-{Stealth[round,length=2.2mm,width=1.8mm]},line width=1.1pt,
               draw=scInkM!45},
  flab/.style={font=\scSmall,text=scInkM,inner sep=2pt},
  stoi/.style={font=\scSmall,text=scInkM,inner sep=0.8pt},
  % verdict marks: shape carries the meaning, colour only reinforces it
  vpass/.style={circle,fill=scAccentF,draw=scAccent,text=scAccent,
                font=\fontsize{6}{6}\selectfont\bfseries,
                minimum size=3.2mm,inner sep=0pt},
  vfail/.style={circle,fill=scAlertF,draw=scAlert,text=scAlert,
                font=\fontsize{6}{6}\selectfont\bfseries,
                minimum size=3.2mm,inner sep=0pt},
  callout/.style={font=\scSmall,text=scInkM,align=left},
}

% \scHead[<badge colour>]{<anchor coord>}{<letter>}{<title>}{<subtitle or empty>}
% One header treatment for every panel in the paper.
\providecommand{\scHead}[5][scInk]{%
  \node[hbullet,fill=#1,anchor=north west] (scbadge#3) at ($#2+(0,-0.015)$)
    {\MakeUppercase{#3}};
  \node[ctag,anchor=west] at ($(scbadge#3.east)+(0.14,0)$) {#4};
  \ifx\relax#5\relax\else
    \node[anchor=north west,font=\scSmall,text=scInkM]
      at ($#2+(0.56,-0.245)$) {#5};
  \fi}

% --- shared pgfplots look (only if the figure loaded pgfplots) ------------
\ifdefined\pgfplotsset
  \pgfplotsset{compat=1.18}
  \pgfplotsset{
    synaxis/.style={
      scale only axis,
      tick align=outside, tick pos=left,
      axis line style={draw=scInkM!55,line width=0.4pt},
      % One grid, on the value axis only.  No minor grid: at these sizes it
      % turns the plot frame into a grey wash.
      ymajorgrids=true, xmajorgrids=false,
      major grid style={draw=scRuleL,line width=0.3pt},
      major tick length=1.6pt,
      label style={font=\scBody,color=scInk},
      tick label style={font=\scSmall,color=scInkM},
      legend style={font=\scSmall,draw=none,fill=none,inner sep=1pt,
                    row sep=0.4pt,legend columns=1},
      legend cell align=left,
      every axis plot/.append style={line width=0.7pt,mark size=1.2pt},
    },
  }
\fi
}
% so that the figure compiles both from Figure/ and from the paper root.
%
% The visual vocabulary follows ../../../../Style/style.tex: the
% colour-blind-safe Okabe--Ito accents, quiet neutral rules and wash, circular
% panel badges, and a compact technical type scale are shared by the companion
% papers.  Semantic aliases below keep the CRN meaning explicit.
%
% DESIGN CONTRACT -----------------------------------------------------------
% 1. Geometry.  Every figure is laid out on a canvas exactly \scW wide (the
%    sn-jnl text width, 372 pt) and is included unscaled, so one point of type
%    is the same physical size in all five figures.  Open each picture with
%    \scCanvas{<height in cm>} to pin the bounding box.
% 2. Colour is semantic, never decorative.  Four roles only:
%       scAccent (green)       structure, the normal/expected path
%       scMark   (orange)      catalyst, certificate, or highlight
%       scAlert  (vermillion)  failure, violation, or uncovered set
%       scRef    (blue)        secondary or reference series
%    A panel gets at most one accent.  Titles are always ink; panel identity
%    is also encoded by a letter inside a circular badge, not by colour alone.
% 3. Type scale.  Three sizes, no others, and nothing below 6 pt:
%       \scTitle (7 pt bold)  panel titles
%       \scBody  (6.5 pt)     labels, values, in-figure prose
%       \scSmall (6 pt)       units, footnotes, dense graph labels
% ---------------------------------------------------------------------------
% Load guard.  The paper inputs every figure source inline, so this file is
% reached once per figure; without the guard the second \newlength below
% would abort the run.
\ifdefined\SynKitStyleLoaded\endinput\fi
\newcommand{\SynKitStyleLoaded}{}

% Load in the document preamble. XeLaTeX/LuaLaTeX and an installed Arial
% family are required; fontspec reports an error instead of substituting.
\RequirePackage[no-math]{fontspec}
\newfontfamily\scArial{Arial}[NFSSFamily=scArial,Ligatures=TeX]
% Keep manuscript mathematics unchanged; select Arial letters and numerals
% only inside figures. TeX supplies specialist mathematical symbols.
\RequirePackage[subdued,italic]{mathastext}
\MTDeclareVersion{synarial}{TU}{scArial}{m}{n}
\SetMathAlphabet{\Mathsf}{synarial}{TU}{scArial}{m}{n}
\SetMathAlphabet{\Mathtt}{synarial}{TU}{scArial}{m}{n}

\usetikzlibrary{positioning,calc,fit,backgrounds,arrows.meta,shapes.geometric,
                shapes.misc,decorations.pathreplacing,patterns,math}

% --- canvas ---------------------------------------------------------------
% 372 pt text width minus the 1 pt standalone border on each side.
\newlength{\scW}\setlength{\scW}{13.0326cm}
\providecommand{\scWnum}{13.0326}
% \scCanvas{<height>} pins the picture to the full text width, y running down
% from 0, so every figure shares one coordinate convention.  It *sets* the
% bounding box rather than contributing to it, so the emitted page is exactly
% the text width and every figure is included at 1:1 -- otherwise a stray
% invisible extent (pgfplots reserves a little around an axis) rescales one
% figure relative to the others and the same \scSmall prints at two sizes.
% Anything drawn outside the canvas is therefore cropped: keep content inside,
% and check with
%     \providecommand{\scCanvas}[1]{\path[draw=red] (0,0) rectangle (\scWnum,-#1);}
\providecommand{\scCanvas}[1]{%
  \path[use as bounding box] (0,0) rectangle (\scWnum,-#1);}

% --- palette --------------------------------------------------------------
% Canonical companion-paper values (../../../../Style/style.tex).
\definecolor{npgOrange}{HTML}{E69F00}
\definecolor{npgSky}{HTML}{56B4E9}
\definecolor{npgGreen}{HTML}{009E73}
\definecolor{npgYellow}{HTML}{F0E442}
\definecolor{npgBlue}{HTML}{0072B2}
\definecolor{npgVermillion}{HTML}{E64B35}
\definecolor{npgPurple}{HTML}{CC79A7}
\definecolor{figInk}{HTML}{1A1A1A}
\definecolor{figMuted}{HTML}{666666}
\definecolor{figHair}{HTML}{D9D9D9}
\definecolor{figWash}{HTML}{F7F7F7}

\colorlet{scInk}{figInk}              % primary text
\colorlet{scInkM}{figMuted}           % secondary text
\colorlet{scPaper}{figWash}           % quiet cards and note strips
\colorlet{scSpF}{npgSky!12}           % species fill
\colorlet{scSpE}{npgBlue}             % species edge
\colorlet{scRxF}{npgOrange!12}        % reaction fill
\colorlet{scRxE}{npgOrange!88!black}  % reaction edge
\colorlet{scReac}{npgBlue!78!black}   % reactant edge
\colorlet{scProd}{npgGreen!82!black}  % product edge
\colorlet{scHiF}{npgOrange!13}        % highlight fill
\colorlet{scHiE}{npgOrange!82!black}  % highlight edge
% Stage / accent tints.  stDE remains available for the fifth plot series.
\colorlet{stAF}{npgGreen!9}\colorlet{stAE}{npgGreen}
\colorlet{stBF}{npgOrange!11}\colorlet{stBE}{npgOrange}
\colorlet{stCF}{npgBlue!9}\colorlet{stCE}{npgBlue}
\colorlet{stDF}{npgPurple!10}\colorlet{stDE}{npgPurple!85!black}
\colorlet{stEF}{npgVermillion!9}\colorlet{stEE}{npgVermillion}

% Semantic aliases.  Figures should speak in these, not in stAE/stBE/...
\colorlet{scAccent}{stAE}\colorlet{scAccentF}{stAF}   % structure / normal
\colorlet{scMark}{stBE}\colorlet{scMarkF}{stBF}       % highlight / certificate
\colorlet{scRef}{stCE}\colorlet{scRefF}{stCF}         % secondary series
\colorlet{scAlert}{stEE}\colorlet{scAlertF}{stEF}     % failure / violation
\colorlet{scRule}{figHair}                            % card and axis hairlines
\colorlet{scRuleL}{figHair!55}                        % grid, dividers

% --- text styles ----------------------------------------------------------
% Three sizes, fixed in points so they do not drift with the document class.
% Arial is shared by labels, headings, panel badges, and chemistry labels.
\providecommand{\scTitle}{\scArial\fontsize{7}{8.4}\selectfont\bfseries\color{scInk}}
\providecommand{\scBody}{\scArial\fontsize{6.5}{7.8}\selectfont\color{scInkM}}
\providecommand{\scSmall}{\scArial\fontsize{6}{7.2}\selectfont\color{scInkM}}
\providecommand{\scSub}[1]{{\scSmall\itshape #1}}
\providecommand{\scPanel}[1]{{\scTitle(#1)}}
% Numerals that carry a result: ink, not colour, so the eye reads the value.
\providecommand{\scNum}{\scArial\fontsize{10}{11}\selectfont\bfseries\color{scInk}}

% --- shared TikZ styles ---------------------------------------------------
\tikzset{
  % text never hyphenates inside a figure
  every node/.append style={execute at begin node={%
      \hyphenpenalty=10000 \exhyphenpenalty=10000 \hbadness=10000 }},
  syn/.style={font=\scArial,color=scInk,
              execute at begin picture={\MTversion*{synarial}},
              >={Stealth[round,length=2.4mm,width=1.9mm]}},
  % containers ------------------------------------------------------------
  % One card chrome for the whole paper: white, 3 pt radius, one hairline.
  panel/.style={rounded corners=2.5pt,fill=none,draw=scRule,line width=0.45pt},
  card/.style={rounded corners=2.2pt,fill=scPaper,draw=scRule,line width=0.45pt,
               inner sep=5pt,align=left,anchor=north west},
  % a quiet strip for a panel's closing fact
  note/.style={rounded corners=2.5pt,fill=scPaper,draw=scRuleL,line width=0.4pt,
               inner xsep=6pt,inner ysep=3.5pt,align=left,font=\scSmall},
  % small labelled value pill
  chip/.style={rounded corners=2pt,fill=scPaper,draw=scRule,line width=0.4pt,
               inner xsep=4pt,inner ysep=2pt,font=\scSmall},
  % Circular letter badge: the same accessible panel marker used in Style/*.
  hbullet/.style={circle,fill=scInk,text=white,draw=none,
                  font=\scArial\fontsize{6}{6}\selectfont\bfseries,
                  minimum size=4.2mm,inner sep=0pt},
  ctag/.style={font=\scTitle,anchor=north west,text=scInk},
  % CRN primitives --------------------------------------------------------
  species/.style={rounded rectangle,fill=scSpF,draw=scSpE,line width=0.45pt,
                  inner xsep=2.4pt,inner ysep=1.6pt,font=\scSmall,
                  text=scInk,minimum height=3.6mm},
  cofactor/.style={species,fill=scRefF,draw=scRef},
  catalyst/.style={species,fill=scHiF,draw=scHiE},
  rxn/.style={rectangle,fill=scRxF,draw=scRxE,line width=0.45pt,
              inner sep=1.2pt,font=\scSmall,text=scInk,
              minimum width=3.0mm,minimum height=3.0mm},
  place/.style={circle,fill=scSpF,draw=scSpE,line width=0.45pt,
                minimum size=4.0mm,inner sep=0pt},
  trans/.style={rectangle,fill=scRxE,draw=scRxE,minimum width=1.1mm,
                minimum height=3.6mm,inner sep=0pt},
  token/.style={circle,fill=scInk,minimum size=1.1mm,inner sep=0pt},
  % edges -----------------------------------------------------------------
  reac/.style={->,line width=0.5pt,draw=scReac,shorten >=0.8pt,shorten <=0.8pt},
  prod/.style={->,line width=0.5pt,draw=scProd,shorten >=0.8pt,shorten <=0.8pt},
  cat/.style={->,line width=0.5pt,draw=scHiE,shorten >=0.8pt,shorten <=0.8pt},
  flow/.style={-{Stealth[round,length=2.2mm,width=1.8mm]},line width=1.1pt,
               draw=scInkM!45},
  flab/.style={font=\scSmall,text=scInkM,inner sep=2pt},
  stoi/.style={font=\scSmall,text=scInkM,inner sep=0.8pt},
  % verdict marks: shape carries the meaning, colour only reinforces it
  vpass/.style={circle,fill=scAccentF,draw=scAccent,text=scAccent,
                font=\fontsize{6}{6}\selectfont\bfseries,
                minimum size=3.2mm,inner sep=0pt},
  vfail/.style={circle,fill=scAlertF,draw=scAlert,text=scAlert,
                font=\fontsize{6}{6}\selectfont\bfseries,
                minimum size=3.2mm,inner sep=0pt},
  callout/.style={font=\scSmall,text=scInkM,align=left},
}

% \scHead[<badge colour>]{<anchor coord>}{<letter>}{<title>}{<subtitle or empty>}
% One header treatment for every panel in the paper.
\providecommand{\scHead}[5][scInk]{%
  \node[hbullet,fill=#1,anchor=north west] (scbadge#3) at ($#2+(0,-0.015)$)
    {\MakeUppercase{#3}};
  \node[ctag,anchor=west] at ($(scbadge#3.east)+(0.14,0)$) {#4};
  \ifx\relax#5\relax\else
    \node[anchor=north west,font=\scSmall,text=scInkM]
      at ($#2+(0.56,-0.245)$) {#5};
  \fi}

% --- shared pgfplots look (only if the figure loaded pgfplots) ------------
\ifdefined\pgfplotsset
  \pgfplotsset{compat=1.18}
  \pgfplotsset{
    synaxis/.style={
      scale only axis,
      tick align=outside, tick pos=left,
      axis line style={draw=scInkM!55,line width=0.4pt},
      % One grid, on the value axis only.  No minor grid: at these sizes it
      % turns the plot frame into a grey wash.
      ymajorgrids=true, xmajorgrids=false,
      major grid style={draw=scRuleL,line width=0.3pt},
      major tick length=1.6pt,
      label style={font=\scBody,color=scInk},
      tick label style={font=\scSmall,color=scInkM},
      legend style={font=\scSmall,draw=none,fill=none,inner sep=1pt,
                    row sep=0.4pt,legend columns=1},
      legend cell align=left,
      every axis plot/.append style={line width=0.7pt,mark size=1.2pt},
    },
  }
\fi
}{% ---------------------------------------------------------------------------
% synstyle.tex --- shared look for every TikZ figure of the SynKit CRN paper.
%
% Input it from a figure file:
%     \IfFileExists{synstyle.tex}{% ---------------------------------------------------------------------------
% synstyle.tex --- shared look for every TikZ figure of the SynKit CRN paper.
%
% Input it from a figure file:
%     \IfFileExists{synstyle.tex}{\input{synstyle}}{\input{Figure/synstyle}}
% so that the figure compiles both from Figure/ and from the paper root.
%
% The visual vocabulary follows ../../../../Style/style.tex: the
% colour-blind-safe Okabe--Ito accents, quiet neutral rules and wash, circular
% panel badges, and a compact technical type scale are shared by the companion
% papers.  Semantic aliases below keep the CRN meaning explicit.
%
% DESIGN CONTRACT -----------------------------------------------------------
% 1. Geometry.  Every figure is laid out on a canvas exactly \scW wide (the
%    sn-jnl text width, 372 pt) and is included unscaled, so one point of type
%    is the same physical size in all five figures.  Open each picture with
%    \scCanvas{<height in cm>} to pin the bounding box.
% 2. Colour is semantic, never decorative.  Four roles only:
%       scAccent (green)       structure, the normal/expected path
%       scMark   (orange)      catalyst, certificate, or highlight
%       scAlert  (vermillion)  failure, violation, or uncovered set
%       scRef    (blue)        secondary or reference series
%    A panel gets at most one accent.  Titles are always ink; panel identity
%    is also encoded by a letter inside a circular badge, not by colour alone.
% 3. Type scale.  Three sizes, no others, and nothing below 6 pt:
%       \scTitle (7 pt bold)  panel titles
%       \scBody  (6.5 pt)     labels, values, in-figure prose
%       \scSmall (6 pt)       units, footnotes, dense graph labels
% ---------------------------------------------------------------------------
% Load guard.  The paper inputs every figure source inline, so this file is
% reached once per figure; without the guard the second \newlength below
% would abort the run.
\ifdefined\SynKitStyleLoaded\endinput\fi
\newcommand{\SynKitStyleLoaded}{}

% Load in the document preamble. XeLaTeX/LuaLaTeX and an installed Arial
% family are required; fontspec reports an error instead of substituting.
\RequirePackage[no-math]{fontspec}
\newfontfamily\scArial{Arial}[NFSSFamily=scArial,Ligatures=TeX]
% Keep manuscript mathematics unchanged; select Arial letters and numerals
% only inside figures. TeX supplies specialist mathematical symbols.
\RequirePackage[subdued,italic]{mathastext}
\MTDeclareVersion{synarial}{TU}{scArial}{m}{n}
\SetMathAlphabet{\Mathsf}{synarial}{TU}{scArial}{m}{n}
\SetMathAlphabet{\Mathtt}{synarial}{TU}{scArial}{m}{n}

\usetikzlibrary{positioning,calc,fit,backgrounds,arrows.meta,shapes.geometric,
                shapes.misc,decorations.pathreplacing,patterns,math}

% --- canvas ---------------------------------------------------------------
% 372 pt text width minus the 1 pt standalone border on each side.
\newlength{\scW}\setlength{\scW}{13.0326cm}
\providecommand{\scWnum}{13.0326}
% \scCanvas{<height>} pins the picture to the full text width, y running down
% from 0, so every figure shares one coordinate convention.  It *sets* the
% bounding box rather than contributing to it, so the emitted page is exactly
% the text width and every figure is included at 1:1 -- otherwise a stray
% invisible extent (pgfplots reserves a little around an axis) rescales one
% figure relative to the others and the same \scSmall prints at two sizes.
% Anything drawn outside the canvas is therefore cropped: keep content inside,
% and check with
%     \providecommand{\scCanvas}[1]{\path[draw=red] (0,0) rectangle (\scWnum,-#1);}
\providecommand{\scCanvas}[1]{%
  \path[use as bounding box] (0,0) rectangle (\scWnum,-#1);}

% --- palette --------------------------------------------------------------
% Canonical companion-paper values (../../../../Style/style.tex).
\definecolor{npgOrange}{HTML}{E69F00}
\definecolor{npgSky}{HTML}{56B4E9}
\definecolor{npgGreen}{HTML}{009E73}
\definecolor{npgYellow}{HTML}{F0E442}
\definecolor{npgBlue}{HTML}{0072B2}
\definecolor{npgVermillion}{HTML}{E64B35}
\definecolor{npgPurple}{HTML}{CC79A7}
\definecolor{figInk}{HTML}{1A1A1A}
\definecolor{figMuted}{HTML}{666666}
\definecolor{figHair}{HTML}{D9D9D9}
\definecolor{figWash}{HTML}{F7F7F7}

\colorlet{scInk}{figInk}              % primary text
\colorlet{scInkM}{figMuted}           % secondary text
\colorlet{scPaper}{figWash}           % quiet cards and note strips
\colorlet{scSpF}{npgSky!12}           % species fill
\colorlet{scSpE}{npgBlue}             % species edge
\colorlet{scRxF}{npgOrange!12}        % reaction fill
\colorlet{scRxE}{npgOrange!88!black}  % reaction edge
\colorlet{scReac}{npgBlue!78!black}   % reactant edge
\colorlet{scProd}{npgGreen!82!black}  % product edge
\colorlet{scHiF}{npgOrange!13}        % highlight fill
\colorlet{scHiE}{npgOrange!82!black}  % highlight edge
% Stage / accent tints.  stDE remains available for the fifth plot series.
\colorlet{stAF}{npgGreen!9}\colorlet{stAE}{npgGreen}
\colorlet{stBF}{npgOrange!11}\colorlet{stBE}{npgOrange}
\colorlet{stCF}{npgBlue!9}\colorlet{stCE}{npgBlue}
\colorlet{stDF}{npgPurple!10}\colorlet{stDE}{npgPurple!85!black}
\colorlet{stEF}{npgVermillion!9}\colorlet{stEE}{npgVermillion}

% Semantic aliases.  Figures should speak in these, not in stAE/stBE/...
\colorlet{scAccent}{stAE}\colorlet{scAccentF}{stAF}   % structure / normal
\colorlet{scMark}{stBE}\colorlet{scMarkF}{stBF}       % highlight / certificate
\colorlet{scRef}{stCE}\colorlet{scRefF}{stCF}         % secondary series
\colorlet{scAlert}{stEE}\colorlet{scAlertF}{stEF}     % failure / violation
\colorlet{scRule}{figHair}                            % card and axis hairlines
\colorlet{scRuleL}{figHair!55}                        % grid, dividers

% --- text styles ----------------------------------------------------------
% Three sizes, fixed in points so they do not drift with the document class.
% Arial is shared by labels, headings, panel badges, and chemistry labels.
\providecommand{\scTitle}{\scArial\fontsize{7}{8.4}\selectfont\bfseries\color{scInk}}
\providecommand{\scBody}{\scArial\fontsize{6.5}{7.8}\selectfont\color{scInkM}}
\providecommand{\scSmall}{\scArial\fontsize{6}{7.2}\selectfont\color{scInkM}}
\providecommand{\scSub}[1]{{\scSmall\itshape #1}}
\providecommand{\scPanel}[1]{{\scTitle(#1)}}
% Numerals that carry a result: ink, not colour, so the eye reads the value.
\providecommand{\scNum}{\scArial\fontsize{10}{11}\selectfont\bfseries\color{scInk}}

% --- shared TikZ styles ---------------------------------------------------
\tikzset{
  % text never hyphenates inside a figure
  every node/.append style={execute at begin node={%
      \hyphenpenalty=10000 \exhyphenpenalty=10000 \hbadness=10000 }},
  syn/.style={font=\scArial,color=scInk,
              execute at begin picture={\MTversion*{synarial}},
              >={Stealth[round,length=2.4mm,width=1.9mm]}},
  % containers ------------------------------------------------------------
  % One card chrome for the whole paper: white, 3 pt radius, one hairline.
  panel/.style={rounded corners=2.5pt,fill=none,draw=scRule,line width=0.45pt},
  card/.style={rounded corners=2.2pt,fill=scPaper,draw=scRule,line width=0.45pt,
               inner sep=5pt,align=left,anchor=north west},
  % a quiet strip for a panel's closing fact
  note/.style={rounded corners=2.5pt,fill=scPaper,draw=scRuleL,line width=0.4pt,
               inner xsep=6pt,inner ysep=3.5pt,align=left,font=\scSmall},
  % small labelled value pill
  chip/.style={rounded corners=2pt,fill=scPaper,draw=scRule,line width=0.4pt,
               inner xsep=4pt,inner ysep=2pt,font=\scSmall},
  % Circular letter badge: the same accessible panel marker used in Style/*.
  hbullet/.style={circle,fill=scInk,text=white,draw=none,
                  font=\scArial\fontsize{6}{6}\selectfont\bfseries,
                  minimum size=4.2mm,inner sep=0pt},
  ctag/.style={font=\scTitle,anchor=north west,text=scInk},
  % CRN primitives --------------------------------------------------------
  species/.style={rounded rectangle,fill=scSpF,draw=scSpE,line width=0.45pt,
                  inner xsep=2.4pt,inner ysep=1.6pt,font=\scSmall,
                  text=scInk,minimum height=3.6mm},
  cofactor/.style={species,fill=scRefF,draw=scRef},
  catalyst/.style={species,fill=scHiF,draw=scHiE},
  rxn/.style={rectangle,fill=scRxF,draw=scRxE,line width=0.45pt,
              inner sep=1.2pt,font=\scSmall,text=scInk,
              minimum width=3.0mm,minimum height=3.0mm},
  place/.style={circle,fill=scSpF,draw=scSpE,line width=0.45pt,
                minimum size=4.0mm,inner sep=0pt},
  trans/.style={rectangle,fill=scRxE,draw=scRxE,minimum width=1.1mm,
                minimum height=3.6mm,inner sep=0pt},
  token/.style={circle,fill=scInk,minimum size=1.1mm,inner sep=0pt},
  % edges -----------------------------------------------------------------
  reac/.style={->,line width=0.5pt,draw=scReac,shorten >=0.8pt,shorten <=0.8pt},
  prod/.style={->,line width=0.5pt,draw=scProd,shorten >=0.8pt,shorten <=0.8pt},
  cat/.style={->,line width=0.5pt,draw=scHiE,shorten >=0.8pt,shorten <=0.8pt},
  flow/.style={-{Stealth[round,length=2.2mm,width=1.8mm]},line width=1.1pt,
               draw=scInkM!45},
  flab/.style={font=\scSmall,text=scInkM,inner sep=2pt},
  stoi/.style={font=\scSmall,text=scInkM,inner sep=0.8pt},
  % verdict marks: shape carries the meaning, colour only reinforces it
  vpass/.style={circle,fill=scAccentF,draw=scAccent,text=scAccent,
                font=\fontsize{6}{6}\selectfont\bfseries,
                minimum size=3.2mm,inner sep=0pt},
  vfail/.style={circle,fill=scAlertF,draw=scAlert,text=scAlert,
                font=\fontsize{6}{6}\selectfont\bfseries,
                minimum size=3.2mm,inner sep=0pt},
  callout/.style={font=\scSmall,text=scInkM,align=left},
}

% \scHead[<badge colour>]{<anchor coord>}{<letter>}{<title>}{<subtitle or empty>}
% One header treatment for every panel in the paper.
\providecommand{\scHead}[5][scInk]{%
  \node[hbullet,fill=#1,anchor=north west] (scbadge#3) at ($#2+(0,-0.015)$)
    {\MakeUppercase{#3}};
  \node[ctag,anchor=west] at ($(scbadge#3.east)+(0.14,0)$) {#4};
  \ifx\relax#5\relax\else
    \node[anchor=north west,font=\scSmall,text=scInkM]
      at ($#2+(0.56,-0.245)$) {#5};
  \fi}

% --- shared pgfplots look (only if the figure loaded pgfplots) ------------
\ifdefined\pgfplotsset
  \pgfplotsset{compat=1.18}
  \pgfplotsset{
    synaxis/.style={
      scale only axis,
      tick align=outside, tick pos=left,
      axis line style={draw=scInkM!55,line width=0.4pt},
      % One grid, on the value axis only.  No minor grid: at these sizes it
      % turns the plot frame into a grey wash.
      ymajorgrids=true, xmajorgrids=false,
      major grid style={draw=scRuleL,line width=0.3pt},
      major tick length=1.6pt,
      label style={font=\scBody,color=scInk},
      tick label style={font=\scSmall,color=scInkM},
      legend style={font=\scSmall,draw=none,fill=none,inner sep=1pt,
                    row sep=0.4pt,legend columns=1},
      legend cell align=left,
      every axis plot/.append style={line width=0.7pt,mark size=1.2pt},
    },
  }
\fi
}{% ---------------------------------------------------------------------------
% synstyle.tex --- shared look for every TikZ figure of the SynKit CRN paper.
%
% Input it from a figure file:
%     \IfFileExists{synstyle.tex}{\input{synstyle}}{\input{Figure/synstyle}}
% so that the figure compiles both from Figure/ and from the paper root.
%
% The visual vocabulary follows ../../../../Style/style.tex: the
% colour-blind-safe Okabe--Ito accents, quiet neutral rules and wash, circular
% panel badges, and a compact technical type scale are shared by the companion
% papers.  Semantic aliases below keep the CRN meaning explicit.
%
% DESIGN CONTRACT -----------------------------------------------------------
% 1. Geometry.  Every figure is laid out on a canvas exactly \scW wide (the
%    sn-jnl text width, 372 pt) and is included unscaled, so one point of type
%    is the same physical size in all five figures.  Open each picture with
%    \scCanvas{<height in cm>} to pin the bounding box.
% 2. Colour is semantic, never decorative.  Four roles only:
%       scAccent (green)       structure, the normal/expected path
%       scMark   (orange)      catalyst, certificate, or highlight
%       scAlert  (vermillion)  failure, violation, or uncovered set
%       scRef    (blue)        secondary or reference series
%    A panel gets at most one accent.  Titles are always ink; panel identity
%    is also encoded by a letter inside a circular badge, not by colour alone.
% 3. Type scale.  Three sizes, no others, and nothing below 6 pt:
%       \scTitle (7 pt bold)  panel titles
%       \scBody  (6.5 pt)     labels, values, in-figure prose
%       \scSmall (6 pt)       units, footnotes, dense graph labels
% ---------------------------------------------------------------------------
% Load guard.  The paper inputs every figure source inline, so this file is
% reached once per figure; without the guard the second \newlength below
% would abort the run.
\ifdefined\SynKitStyleLoaded\endinput\fi
\newcommand{\SynKitStyleLoaded}{}

% Load in the document preamble. XeLaTeX/LuaLaTeX and an installed Arial
% family are required; fontspec reports an error instead of substituting.
\RequirePackage[no-math]{fontspec}
\newfontfamily\scArial{Arial}[NFSSFamily=scArial,Ligatures=TeX]
% Keep manuscript mathematics unchanged; select Arial letters and numerals
% only inside figures. TeX supplies specialist mathematical symbols.
\RequirePackage[subdued,italic]{mathastext}
\MTDeclareVersion{synarial}{TU}{scArial}{m}{n}
\SetMathAlphabet{\Mathsf}{synarial}{TU}{scArial}{m}{n}
\SetMathAlphabet{\Mathtt}{synarial}{TU}{scArial}{m}{n}

\usetikzlibrary{positioning,calc,fit,backgrounds,arrows.meta,shapes.geometric,
                shapes.misc,decorations.pathreplacing,patterns,math}

% --- canvas ---------------------------------------------------------------
% 372 pt text width minus the 1 pt standalone border on each side.
\newlength{\scW}\setlength{\scW}{13.0326cm}
\providecommand{\scWnum}{13.0326}
% \scCanvas{<height>} pins the picture to the full text width, y running down
% from 0, so every figure shares one coordinate convention.  It *sets* the
% bounding box rather than contributing to it, so the emitted page is exactly
% the text width and every figure is included at 1:1 -- otherwise a stray
% invisible extent (pgfplots reserves a little around an axis) rescales one
% figure relative to the others and the same \scSmall prints at two sizes.
% Anything drawn outside the canvas is therefore cropped: keep content inside,
% and check with
%     \providecommand{\scCanvas}[1]{\path[draw=red] (0,0) rectangle (\scWnum,-#1);}
\providecommand{\scCanvas}[1]{%
  \path[use as bounding box] (0,0) rectangle (\scWnum,-#1);}

% --- palette --------------------------------------------------------------
% Canonical companion-paper values (../../../../Style/style.tex).
\definecolor{npgOrange}{HTML}{E69F00}
\definecolor{npgSky}{HTML}{56B4E9}
\definecolor{npgGreen}{HTML}{009E73}
\definecolor{npgYellow}{HTML}{F0E442}
\definecolor{npgBlue}{HTML}{0072B2}
\definecolor{npgVermillion}{HTML}{E64B35}
\definecolor{npgPurple}{HTML}{CC79A7}
\definecolor{figInk}{HTML}{1A1A1A}
\definecolor{figMuted}{HTML}{666666}
\definecolor{figHair}{HTML}{D9D9D9}
\definecolor{figWash}{HTML}{F7F7F7}

\colorlet{scInk}{figInk}              % primary text
\colorlet{scInkM}{figMuted}           % secondary text
\colorlet{scPaper}{figWash}           % quiet cards and note strips
\colorlet{scSpF}{npgSky!12}           % species fill
\colorlet{scSpE}{npgBlue}             % species edge
\colorlet{scRxF}{npgOrange!12}        % reaction fill
\colorlet{scRxE}{npgOrange!88!black}  % reaction edge
\colorlet{scReac}{npgBlue!78!black}   % reactant edge
\colorlet{scProd}{npgGreen!82!black}  % product edge
\colorlet{scHiF}{npgOrange!13}        % highlight fill
\colorlet{scHiE}{npgOrange!82!black}  % highlight edge
% Stage / accent tints.  stDE remains available for the fifth plot series.
\colorlet{stAF}{npgGreen!9}\colorlet{stAE}{npgGreen}
\colorlet{stBF}{npgOrange!11}\colorlet{stBE}{npgOrange}
\colorlet{stCF}{npgBlue!9}\colorlet{stCE}{npgBlue}
\colorlet{stDF}{npgPurple!10}\colorlet{stDE}{npgPurple!85!black}
\colorlet{stEF}{npgVermillion!9}\colorlet{stEE}{npgVermillion}

% Semantic aliases.  Figures should speak in these, not in stAE/stBE/...
\colorlet{scAccent}{stAE}\colorlet{scAccentF}{stAF}   % structure / normal
\colorlet{scMark}{stBE}\colorlet{scMarkF}{stBF}       % highlight / certificate
\colorlet{scRef}{stCE}\colorlet{scRefF}{stCF}         % secondary series
\colorlet{scAlert}{stEE}\colorlet{scAlertF}{stEF}     % failure / violation
\colorlet{scRule}{figHair}                            % card and axis hairlines
\colorlet{scRuleL}{figHair!55}                        % grid, dividers

% --- text styles ----------------------------------------------------------
% Three sizes, fixed in points so they do not drift with the document class.
% Arial is shared by labels, headings, panel badges, and chemistry labels.
\providecommand{\scTitle}{\scArial\fontsize{7}{8.4}\selectfont\bfseries\color{scInk}}
\providecommand{\scBody}{\scArial\fontsize{6.5}{7.8}\selectfont\color{scInkM}}
\providecommand{\scSmall}{\scArial\fontsize{6}{7.2}\selectfont\color{scInkM}}
\providecommand{\scSub}[1]{{\scSmall\itshape #1}}
\providecommand{\scPanel}[1]{{\scTitle(#1)}}
% Numerals that carry a result: ink, not colour, so the eye reads the value.
\providecommand{\scNum}{\scArial\fontsize{10}{11}\selectfont\bfseries\color{scInk}}

% --- shared TikZ styles ---------------------------------------------------
\tikzset{
  % text never hyphenates inside a figure
  every node/.append style={execute at begin node={%
      \hyphenpenalty=10000 \exhyphenpenalty=10000 \hbadness=10000 }},
  syn/.style={font=\scArial,color=scInk,
              execute at begin picture={\MTversion*{synarial}},
              >={Stealth[round,length=2.4mm,width=1.9mm]}},
  % containers ------------------------------------------------------------
  % One card chrome for the whole paper: white, 3 pt radius, one hairline.
  panel/.style={rounded corners=2.5pt,fill=none,draw=scRule,line width=0.45pt},
  card/.style={rounded corners=2.2pt,fill=scPaper,draw=scRule,line width=0.45pt,
               inner sep=5pt,align=left,anchor=north west},
  % a quiet strip for a panel's closing fact
  note/.style={rounded corners=2.5pt,fill=scPaper,draw=scRuleL,line width=0.4pt,
               inner xsep=6pt,inner ysep=3.5pt,align=left,font=\scSmall},
  % small labelled value pill
  chip/.style={rounded corners=2pt,fill=scPaper,draw=scRule,line width=0.4pt,
               inner xsep=4pt,inner ysep=2pt,font=\scSmall},
  % Circular letter badge: the same accessible panel marker used in Style/*.
  hbullet/.style={circle,fill=scInk,text=white,draw=none,
                  font=\scArial\fontsize{6}{6}\selectfont\bfseries,
                  minimum size=4.2mm,inner sep=0pt},
  ctag/.style={font=\scTitle,anchor=north west,text=scInk},
  % CRN primitives --------------------------------------------------------
  species/.style={rounded rectangle,fill=scSpF,draw=scSpE,line width=0.45pt,
                  inner xsep=2.4pt,inner ysep=1.6pt,font=\scSmall,
                  text=scInk,minimum height=3.6mm},
  cofactor/.style={species,fill=scRefF,draw=scRef},
  catalyst/.style={species,fill=scHiF,draw=scHiE},
  rxn/.style={rectangle,fill=scRxF,draw=scRxE,line width=0.45pt,
              inner sep=1.2pt,font=\scSmall,text=scInk,
              minimum width=3.0mm,minimum height=3.0mm},
  place/.style={circle,fill=scSpF,draw=scSpE,line width=0.45pt,
                minimum size=4.0mm,inner sep=0pt},
  trans/.style={rectangle,fill=scRxE,draw=scRxE,minimum width=1.1mm,
                minimum height=3.6mm,inner sep=0pt},
  token/.style={circle,fill=scInk,minimum size=1.1mm,inner sep=0pt},
  % edges -----------------------------------------------------------------
  reac/.style={->,line width=0.5pt,draw=scReac,shorten >=0.8pt,shorten <=0.8pt},
  prod/.style={->,line width=0.5pt,draw=scProd,shorten >=0.8pt,shorten <=0.8pt},
  cat/.style={->,line width=0.5pt,draw=scHiE,shorten >=0.8pt,shorten <=0.8pt},
  flow/.style={-{Stealth[round,length=2.2mm,width=1.8mm]},line width=1.1pt,
               draw=scInkM!45},
  flab/.style={font=\scSmall,text=scInkM,inner sep=2pt},
  stoi/.style={font=\scSmall,text=scInkM,inner sep=0.8pt},
  % verdict marks: shape carries the meaning, colour only reinforces it
  vpass/.style={circle,fill=scAccentF,draw=scAccent,text=scAccent,
                font=\fontsize{6}{6}\selectfont\bfseries,
                minimum size=3.2mm,inner sep=0pt},
  vfail/.style={circle,fill=scAlertF,draw=scAlert,text=scAlert,
                font=\fontsize{6}{6}\selectfont\bfseries,
                minimum size=3.2mm,inner sep=0pt},
  callout/.style={font=\scSmall,text=scInkM,align=left},
}

% \scHead[<badge colour>]{<anchor coord>}{<letter>}{<title>}{<subtitle or empty>}
% One header treatment for every panel in the paper.
\providecommand{\scHead}[5][scInk]{%
  \node[hbullet,fill=#1,anchor=north west] (scbadge#3) at ($#2+(0,-0.015)$)
    {\MakeUppercase{#3}};
  \node[ctag,anchor=west] at ($(scbadge#3.east)+(0.14,0)$) {#4};
  \ifx\relax#5\relax\else
    \node[anchor=north west,font=\scSmall,text=scInkM]
      at ($#2+(0.56,-0.245)$) {#5};
  \fi}

% --- shared pgfplots look (only if the figure loaded pgfplots) ------------
\ifdefined\pgfplotsset
  \pgfplotsset{compat=1.18}
  \pgfplotsset{
    synaxis/.style={
      scale only axis,
      tick align=outside, tick pos=left,
      axis line style={draw=scInkM!55,line width=0.4pt},
      % One grid, on the value axis only.  No minor grid: at these sizes it
      % turns the plot frame into a grey wash.
      ymajorgrids=true, xmajorgrids=false,
      major grid style={draw=scRuleL,line width=0.3pt},
      major tick length=1.6pt,
      label style={font=\scBody,color=scInk},
      tick label style={font=\scSmall,color=scInkM},
      legend style={font=\scSmall,draw=none,fill=none,inner sep=1pt,
                    row sep=0.4pt,legend columns=1},
      legend cell align=left,
      every axis plot/.append style={line width=0.7pt,mark size=1.2pt},
    },
  }
\fi
}
% so that the figure compiles both from Figure/ and from the paper root.
%
% The visual vocabulary follows ../../../../Style/style.tex: the
% colour-blind-safe Okabe--Ito accents, quiet neutral rules and wash, circular
% panel badges, and a compact technical type scale are shared by the companion
% papers.  Semantic aliases below keep the CRN meaning explicit.
%
% DESIGN CONTRACT -----------------------------------------------------------
% 1. Geometry.  Every figure is laid out on a canvas exactly \scW wide (the
%    sn-jnl text width, 372 pt) and is included unscaled, so one point of type
%    is the same physical size in all five figures.  Open each picture with
%    \scCanvas{<height in cm>} to pin the bounding box.
% 2. Colour is semantic, never decorative.  Four roles only:
%       scAccent (green)       structure, the normal/expected path
%       scMark   (orange)      catalyst, certificate, or highlight
%       scAlert  (vermillion)  failure, violation, or uncovered set
%       scRef    (blue)        secondary or reference series
%    A panel gets at most one accent.  Titles are always ink; panel identity
%    is also encoded by a letter inside a circular badge, not by colour alone.
% 3. Type scale.  Three sizes, no others, and nothing below 6 pt:
%       \scTitle (7 pt bold)  panel titles
%       \scBody  (6.5 pt)     labels, values, in-figure prose
%       \scSmall (6 pt)       units, footnotes, dense graph labels
% ---------------------------------------------------------------------------
% Load guard.  The paper inputs every figure source inline, so this file is
% reached once per figure; without the guard the second \newlength below
% would abort the run.
\ifdefined\SynKitStyleLoaded\endinput\fi
\newcommand{\SynKitStyleLoaded}{}

% Load in the document preamble. XeLaTeX/LuaLaTeX and an installed Arial
% family are required; fontspec reports an error instead of substituting.
\RequirePackage[no-math]{fontspec}
\newfontfamily\scArial{Arial}[NFSSFamily=scArial,Ligatures=TeX]
% Keep manuscript mathematics unchanged; select Arial letters and numerals
% only inside figures. TeX supplies specialist mathematical symbols.
\RequirePackage[subdued,italic]{mathastext}
\MTDeclareVersion{synarial}{TU}{scArial}{m}{n}
\SetMathAlphabet{\Mathsf}{synarial}{TU}{scArial}{m}{n}
\SetMathAlphabet{\Mathtt}{synarial}{TU}{scArial}{m}{n}

\usetikzlibrary{positioning,calc,fit,backgrounds,arrows.meta,shapes.geometric,
                shapes.misc,decorations.pathreplacing,patterns,math}

% --- canvas ---------------------------------------------------------------
% 372 pt text width minus the 1 pt standalone border on each side.
\newlength{\scW}\setlength{\scW}{13.0326cm}
\providecommand{\scWnum}{13.0326}
% \scCanvas{<height>} pins the picture to the full text width, y running down
% from 0, so every figure shares one coordinate convention.  It *sets* the
% bounding box rather than contributing to it, so the emitted page is exactly
% the text width and every figure is included at 1:1 -- otherwise a stray
% invisible extent (pgfplots reserves a little around an axis) rescales one
% figure relative to the others and the same \scSmall prints at two sizes.
% Anything drawn outside the canvas is therefore cropped: keep content inside,
% and check with
%     \providecommand{\scCanvas}[1]{\path[draw=red] (0,0) rectangle (\scWnum,-#1);}
\providecommand{\scCanvas}[1]{%
  \path[use as bounding box] (0,0) rectangle (\scWnum,-#1);}

% --- palette --------------------------------------------------------------
% Canonical companion-paper values (../../../../Style/style.tex).
\definecolor{npgOrange}{HTML}{E69F00}
\definecolor{npgSky}{HTML}{56B4E9}
\definecolor{npgGreen}{HTML}{009E73}
\definecolor{npgYellow}{HTML}{F0E442}
\definecolor{npgBlue}{HTML}{0072B2}
\definecolor{npgVermillion}{HTML}{E64B35}
\definecolor{npgPurple}{HTML}{CC79A7}
\definecolor{figInk}{HTML}{1A1A1A}
\definecolor{figMuted}{HTML}{666666}
\definecolor{figHair}{HTML}{D9D9D9}
\definecolor{figWash}{HTML}{F7F7F7}

\colorlet{scInk}{figInk}              % primary text
\colorlet{scInkM}{figMuted}           % secondary text
\colorlet{scPaper}{figWash}           % quiet cards and note strips
\colorlet{scSpF}{npgSky!12}           % species fill
\colorlet{scSpE}{npgBlue}             % species edge
\colorlet{scRxF}{npgOrange!12}        % reaction fill
\colorlet{scRxE}{npgOrange!88!black}  % reaction edge
\colorlet{scReac}{npgBlue!78!black}   % reactant edge
\colorlet{scProd}{npgGreen!82!black}  % product edge
\colorlet{scHiF}{npgOrange!13}        % highlight fill
\colorlet{scHiE}{npgOrange!82!black}  % highlight edge
% Stage / accent tints.  stDE remains available for the fifth plot series.
\colorlet{stAF}{npgGreen!9}\colorlet{stAE}{npgGreen}
\colorlet{stBF}{npgOrange!11}\colorlet{stBE}{npgOrange}
\colorlet{stCF}{npgBlue!9}\colorlet{stCE}{npgBlue}
\colorlet{stDF}{npgPurple!10}\colorlet{stDE}{npgPurple!85!black}
\colorlet{stEF}{npgVermillion!9}\colorlet{stEE}{npgVermillion}

% Semantic aliases.  Figures should speak in these, not in stAE/stBE/...
\colorlet{scAccent}{stAE}\colorlet{scAccentF}{stAF}   % structure / normal
\colorlet{scMark}{stBE}\colorlet{scMarkF}{stBF}       % highlight / certificate
\colorlet{scRef}{stCE}\colorlet{scRefF}{stCF}         % secondary series
\colorlet{scAlert}{stEE}\colorlet{scAlertF}{stEF}     % failure / violation
\colorlet{scRule}{figHair}                            % card and axis hairlines
\colorlet{scRuleL}{figHair!55}                        % grid, dividers

% --- text styles ----------------------------------------------------------
% Three sizes, fixed in points so they do not drift with the document class.
% Arial is shared by labels, headings, panel badges, and chemistry labels.
\providecommand{\scTitle}{\scArial\fontsize{7}{8.4}\selectfont\bfseries\color{scInk}}
\providecommand{\scBody}{\scArial\fontsize{6.5}{7.8}\selectfont\color{scInkM}}
\providecommand{\scSmall}{\scArial\fontsize{6}{7.2}\selectfont\color{scInkM}}
\providecommand{\scSub}[1]{{\scSmall\itshape #1}}
\providecommand{\scPanel}[1]{{\scTitle(#1)}}
% Numerals that carry a result: ink, not colour, so the eye reads the value.
\providecommand{\scNum}{\scArial\fontsize{10}{11}\selectfont\bfseries\color{scInk}}

% --- shared TikZ styles ---------------------------------------------------
\tikzset{
  % text never hyphenates inside a figure
  every node/.append style={execute at begin node={%
      \hyphenpenalty=10000 \exhyphenpenalty=10000 \hbadness=10000 }},
  syn/.style={font=\scArial,color=scInk,
              execute at begin picture={\MTversion*{synarial}},
              >={Stealth[round,length=2.4mm,width=1.9mm]}},
  % containers ------------------------------------------------------------
  % One card chrome for the whole paper: white, 3 pt radius, one hairline.
  panel/.style={rounded corners=2.5pt,fill=none,draw=scRule,line width=0.45pt},
  card/.style={rounded corners=2.2pt,fill=scPaper,draw=scRule,line width=0.45pt,
               inner sep=5pt,align=left,anchor=north west},
  % a quiet strip for a panel's closing fact
  note/.style={rounded corners=2.5pt,fill=scPaper,draw=scRuleL,line width=0.4pt,
               inner xsep=6pt,inner ysep=3.5pt,align=left,font=\scSmall},
  % small labelled value pill
  chip/.style={rounded corners=2pt,fill=scPaper,draw=scRule,line width=0.4pt,
               inner xsep=4pt,inner ysep=2pt,font=\scSmall},
  % Circular letter badge: the same accessible panel marker used in Style/*.
  hbullet/.style={circle,fill=scInk,text=white,draw=none,
                  font=\scArial\fontsize{6}{6}\selectfont\bfseries,
                  minimum size=4.2mm,inner sep=0pt},
  ctag/.style={font=\scTitle,anchor=north west,text=scInk},
  % CRN primitives --------------------------------------------------------
  species/.style={rounded rectangle,fill=scSpF,draw=scSpE,line width=0.45pt,
                  inner xsep=2.4pt,inner ysep=1.6pt,font=\scSmall,
                  text=scInk,minimum height=3.6mm},
  cofactor/.style={species,fill=scRefF,draw=scRef},
  catalyst/.style={species,fill=scHiF,draw=scHiE},
  rxn/.style={rectangle,fill=scRxF,draw=scRxE,line width=0.45pt,
              inner sep=1.2pt,font=\scSmall,text=scInk,
              minimum width=3.0mm,minimum height=3.0mm},
  place/.style={circle,fill=scSpF,draw=scSpE,line width=0.45pt,
                minimum size=4.0mm,inner sep=0pt},
  trans/.style={rectangle,fill=scRxE,draw=scRxE,minimum width=1.1mm,
                minimum height=3.6mm,inner sep=0pt},
  token/.style={circle,fill=scInk,minimum size=1.1mm,inner sep=0pt},
  % edges -----------------------------------------------------------------
  reac/.style={->,line width=0.5pt,draw=scReac,shorten >=0.8pt,shorten <=0.8pt},
  prod/.style={->,line width=0.5pt,draw=scProd,shorten >=0.8pt,shorten <=0.8pt},
  cat/.style={->,line width=0.5pt,draw=scHiE,shorten >=0.8pt,shorten <=0.8pt},
  flow/.style={-{Stealth[round,length=2.2mm,width=1.8mm]},line width=1.1pt,
               draw=scInkM!45},
  flab/.style={font=\scSmall,text=scInkM,inner sep=2pt},
  stoi/.style={font=\scSmall,text=scInkM,inner sep=0.8pt},
  % verdict marks: shape carries the meaning, colour only reinforces it
  vpass/.style={circle,fill=scAccentF,draw=scAccent,text=scAccent,
                font=\fontsize{6}{6}\selectfont\bfseries,
                minimum size=3.2mm,inner sep=0pt},
  vfail/.style={circle,fill=scAlertF,draw=scAlert,text=scAlert,
                font=\fontsize{6}{6}\selectfont\bfseries,
                minimum size=3.2mm,inner sep=0pt},
  callout/.style={font=\scSmall,text=scInkM,align=left},
}

% \scHead[<badge colour>]{<anchor coord>}{<letter>}{<title>}{<subtitle or empty>}
% One header treatment for every panel in the paper.
\providecommand{\scHead}[5][scInk]{%
  \node[hbullet,fill=#1,anchor=north west] (scbadge#3) at ($#2+(0,-0.015)$)
    {\MakeUppercase{#3}};
  \node[ctag,anchor=west] at ($(scbadge#3.east)+(0.14,0)$) {#4};
  \ifx\relax#5\relax\else
    \node[anchor=north west,font=\scSmall,text=scInkM]
      at ($#2+(0.56,-0.245)$) {#5};
  \fi}

% --- shared pgfplots look (only if the figure loaded pgfplots) ------------
\ifdefined\pgfplotsset
  \pgfplotsset{compat=1.18}
  \pgfplotsset{
    synaxis/.style={
      scale only axis,
      tick align=outside, tick pos=left,
      axis line style={draw=scInkM!55,line width=0.4pt},
      % One grid, on the value axis only.  No minor grid: at these sizes it
      % turns the plot frame into a grey wash.
      ymajorgrids=true, xmajorgrids=false,
      major grid style={draw=scRuleL,line width=0.3pt},
      major tick length=1.6pt,
      label style={font=\scBody,color=scInk},
      tick label style={font=\scSmall,color=scInkM},
      legend style={font=\scSmall,draw=none,fill=none,inner sep=1pt,
                    row sep=0.4pt,legend columns=1},
      legend cell align=left,
      every axis plot/.append style={line width=0.7pt,mark size=1.2pt},
    },
  }
\fi
}
% so that the figure compiles both from Figure/ and from the paper root.
%
% The visual vocabulary follows ../../../../Style/style.tex: the
% colour-blind-safe Okabe--Ito accents, quiet neutral rules and wash, circular
% panel badges, and a compact technical type scale are shared by the companion
% papers.  Semantic aliases below keep the CRN meaning explicit.
%
% DESIGN CONTRACT -----------------------------------------------------------
% 1. Geometry.  Every figure is laid out on a canvas exactly \scW wide (the
%    sn-jnl text width, 372 pt) and is included unscaled, so one point of type
%    is the same physical size in all five figures.  Open each picture with
%    \scCanvas{<height in cm>} to pin the bounding box.
% 2. Colour is semantic, never decorative.  Four roles only:
%       scAccent (green)       structure, the normal/expected path
%       scMark   (orange)      catalyst, certificate, or highlight
%       scAlert  (vermillion)  failure, violation, or uncovered set
%       scRef    (blue)        secondary or reference series
%    A panel gets at most one accent.  Titles are always ink; panel identity
%    is also encoded by a letter inside a circular badge, not by colour alone.
% 3. Type scale.  Three sizes, no others, and nothing below 6 pt:
%       \scTitle (7 pt bold)  panel titles
%       \scBody  (6.5 pt)     labels, values, in-figure prose
%       \scSmall (6 pt)       units, footnotes, dense graph labels
% ---------------------------------------------------------------------------
% Load guard.  The paper inputs every figure source inline, so this file is
% reached once per figure; without the guard the second \newlength below
% would abort the run.
\ifdefined\SynKitStyleLoaded\endinput\fi
\newcommand{\SynKitStyleLoaded}{}

% Load in the document preamble. XeLaTeX/LuaLaTeX and an installed Arial
% family are required; fontspec reports an error instead of substituting.
\RequirePackage[no-math]{fontspec}
\newfontfamily\scArial{Arial}[NFSSFamily=scArial,Ligatures=TeX]
% Keep manuscript mathematics unchanged; select Arial letters and numerals
% only inside figures. TeX supplies specialist mathematical symbols.
\RequirePackage[subdued,italic]{mathastext}
\MTDeclareVersion{synarial}{TU}{scArial}{m}{n}
\SetMathAlphabet{\Mathsf}{synarial}{TU}{scArial}{m}{n}
\SetMathAlphabet{\Mathtt}{synarial}{TU}{scArial}{m}{n}

\usetikzlibrary{positioning,calc,fit,backgrounds,arrows.meta,shapes.geometric,
                shapes.misc,decorations.pathreplacing,patterns,math}

% --- canvas ---------------------------------------------------------------
% 372 pt text width minus the 1 pt standalone border on each side.
\newlength{\scW}\setlength{\scW}{13.0326cm}
\providecommand{\scWnum}{13.0326}
% \scCanvas{<height>} pins the picture to the full text width, y running down
% from 0, so every figure shares one coordinate convention.  It *sets* the
% bounding box rather than contributing to it, so the emitted page is exactly
% the text width and every figure is included at 1:1 -- otherwise a stray
% invisible extent (pgfplots reserves a little around an axis) rescales one
% figure relative to the others and the same \scSmall prints at two sizes.
% Anything drawn outside the canvas is therefore cropped: keep content inside,
% and check with
%     \providecommand{\scCanvas}[1]{\path[draw=red] (0,0) rectangle (\scWnum,-#1);}
\providecommand{\scCanvas}[1]{%
  \path[use as bounding box] (0,0) rectangle (\scWnum,-#1);}

% --- palette --------------------------------------------------------------
% Canonical companion-paper values (../../../../Style/style.tex).
\definecolor{npgOrange}{HTML}{E69F00}
\definecolor{npgSky}{HTML}{56B4E9}
\definecolor{npgGreen}{HTML}{009E73}
\definecolor{npgYellow}{HTML}{F0E442}
\definecolor{npgBlue}{HTML}{0072B2}
\definecolor{npgVermillion}{HTML}{E64B35}
\definecolor{npgPurple}{HTML}{CC79A7}
\definecolor{figInk}{HTML}{1A1A1A}
\definecolor{figMuted}{HTML}{666666}
\definecolor{figHair}{HTML}{D9D9D9}
\definecolor{figWash}{HTML}{F7F7F7}

\colorlet{scInk}{figInk}              % primary text
\colorlet{scInkM}{figMuted}           % secondary text
\colorlet{scPaper}{figWash}           % quiet cards and note strips
\colorlet{scSpF}{npgSky!12}           % species fill
\colorlet{scSpE}{npgBlue}             % species edge
\colorlet{scRxF}{npgOrange!12}        % reaction fill
\colorlet{scRxE}{npgOrange!88!black}  % reaction edge
\colorlet{scReac}{npgBlue!78!black}   % reactant edge
\colorlet{scProd}{npgGreen!82!black}  % product edge
\colorlet{scHiF}{npgOrange!13}        % highlight fill
\colorlet{scHiE}{npgOrange!82!black}  % highlight edge
% Stage / accent tints.  stDE remains available for the fifth plot series.
\colorlet{stAF}{npgGreen!9}\colorlet{stAE}{npgGreen}
\colorlet{stBF}{npgOrange!11}\colorlet{stBE}{npgOrange}
\colorlet{stCF}{npgBlue!9}\colorlet{stCE}{npgBlue}
\colorlet{stDF}{npgPurple!10}\colorlet{stDE}{npgPurple!85!black}
\colorlet{stEF}{npgVermillion!9}\colorlet{stEE}{npgVermillion}

% Semantic aliases.  Figures should speak in these, not in stAE/stBE/...
\colorlet{scAccent}{stAE}\colorlet{scAccentF}{stAF}   % structure / normal
\colorlet{scMark}{stBE}\colorlet{scMarkF}{stBF}       % highlight / certificate
\colorlet{scRef}{stCE}\colorlet{scRefF}{stCF}         % secondary series
\colorlet{scAlert}{stEE}\colorlet{scAlertF}{stEF}     % failure / violation
\colorlet{scRule}{figHair}                            % card and axis hairlines
\colorlet{scRuleL}{figHair!55}                        % grid, dividers

% --- text styles ----------------------------------------------------------
% Three sizes, fixed in points so they do not drift with the document class.
% Arial is shared by labels, headings, panel badges, and chemistry labels.
\providecommand{\scTitle}{\scArial\fontsize{7}{8.4}\selectfont\bfseries\color{scInk}}
\providecommand{\scBody}{\scArial\fontsize{6.5}{7.8}\selectfont\color{scInkM}}
\providecommand{\scSmall}{\scArial\fontsize{6}{7.2}\selectfont\color{scInkM}}
\providecommand{\scSub}[1]{{\scSmall\itshape #1}}
\providecommand{\scPanel}[1]{{\scTitle(#1)}}
% Numerals that carry a result: ink, not colour, so the eye reads the value.
\providecommand{\scNum}{\scArial\fontsize{10}{11}\selectfont\bfseries\color{scInk}}

% --- shared TikZ styles ---------------------------------------------------
\tikzset{
  % text never hyphenates inside a figure
  every node/.append style={execute at begin node={%
      \hyphenpenalty=10000 \exhyphenpenalty=10000 \hbadness=10000 }},
  syn/.style={font=\scArial,color=scInk,
              execute at begin picture={\MTversion*{synarial}},
              >={Stealth[round,length=2.4mm,width=1.9mm]}},
  % containers ------------------------------------------------------------
  % One card chrome for the whole paper: white, 3 pt radius, one hairline.
  panel/.style={rounded corners=2.5pt,fill=none,draw=scRule,line width=0.45pt},
  card/.style={rounded corners=2.2pt,fill=scPaper,draw=scRule,line width=0.45pt,
               inner sep=5pt,align=left,anchor=north west},
  % a quiet strip for a panel's closing fact
  note/.style={rounded corners=2.5pt,fill=scPaper,draw=scRuleL,line width=0.4pt,
               inner xsep=6pt,inner ysep=3.5pt,align=left,font=\scSmall},
  % small labelled value pill
  chip/.style={rounded corners=2pt,fill=scPaper,draw=scRule,line width=0.4pt,
               inner xsep=4pt,inner ysep=2pt,font=\scSmall},
  % Circular letter badge: the same accessible panel marker used in Style/*.
  hbullet/.style={circle,fill=scInk,text=white,draw=none,
                  font=\scArial\fontsize{6}{6}\selectfont\bfseries,
                  minimum size=4.2mm,inner sep=0pt},
  ctag/.style={font=\scTitle,anchor=north west,text=scInk},
  % CRN primitives --------------------------------------------------------
  species/.style={rounded rectangle,fill=scSpF,draw=scSpE,line width=0.45pt,
                  inner xsep=2.4pt,inner ysep=1.6pt,font=\scSmall,
                  text=scInk,minimum height=3.6mm},
  cofactor/.style={species,fill=scRefF,draw=scRef},
  catalyst/.style={species,fill=scHiF,draw=scHiE},
  rxn/.style={rectangle,fill=scRxF,draw=scRxE,line width=0.45pt,
              inner sep=1.2pt,font=\scSmall,text=scInk,
              minimum width=3.0mm,minimum height=3.0mm},
  place/.style={circle,fill=scSpF,draw=scSpE,line width=0.45pt,
                minimum size=4.0mm,inner sep=0pt},
  trans/.style={rectangle,fill=scRxE,draw=scRxE,minimum width=1.1mm,
                minimum height=3.6mm,inner sep=0pt},
  token/.style={circle,fill=scInk,minimum size=1.1mm,inner sep=0pt},
  % edges -----------------------------------------------------------------
  reac/.style={->,line width=0.5pt,draw=scReac,shorten >=0.8pt,shorten <=0.8pt},
  prod/.style={->,line width=0.5pt,draw=scProd,shorten >=0.8pt,shorten <=0.8pt},
  cat/.style={->,line width=0.5pt,draw=scHiE,shorten >=0.8pt,shorten <=0.8pt},
  flow/.style={-{Stealth[round,length=2.2mm,width=1.8mm]},line width=1.1pt,
               draw=scInkM!45},
  flab/.style={font=\scSmall,text=scInkM,inner sep=2pt},
  stoi/.style={font=\scSmall,text=scInkM,inner sep=0.8pt},
  % verdict marks: shape carries the meaning, colour only reinforces it
  vpass/.style={circle,fill=scAccentF,draw=scAccent,text=scAccent,
                font=\fontsize{6}{6}\selectfont\bfseries,
                minimum size=3.2mm,inner sep=0pt},
  vfail/.style={circle,fill=scAlertF,draw=scAlert,text=scAlert,
                font=\fontsize{6}{6}\selectfont\bfseries,
                minimum size=3.2mm,inner sep=0pt},
  callout/.style={font=\scSmall,text=scInkM,align=left},
}

% \scHead[<badge colour>]{<anchor coord>}{<letter>}{<title>}{<subtitle or empty>}
% One header treatment for every panel in the paper.
\providecommand{\scHead}[5][scInk]{%
  \node[hbullet,fill=#1,anchor=north west] (scbadge#3) at ($#2+(0,-0.015)$)
    {\MakeUppercase{#3}};
  \node[ctag,anchor=west] at ($(scbadge#3.east)+(0.14,0)$) {#4};
  \ifx\relax#5\relax\else
    \node[anchor=north west,font=\scSmall,text=scInkM]
      at ($#2+(0.56,-0.245)$) {#5};
  \fi}

% --- shared pgfplots look (only if the figure loaded pgfplots) ------------
\ifdefined\pgfplotsset
  \pgfplotsset{compat=1.18}
  \pgfplotsset{
    synaxis/.style={
      scale only axis,
      tick align=outside, tick pos=left,
      axis line style={draw=scInkM!55,line width=0.4pt},
      % One grid, on the value axis only.  No minor grid: at these sizes it
      % turns the plot frame into a grey wash.
      ymajorgrids=true, xmajorgrids=false,
      major grid style={draw=scRuleL,line width=0.3pt},
      major tick length=1.6pt,
      label style={font=\scBody,color=scInk},
      tick label style={font=\scSmall,color=scInkM},
      legend style={font=\scSmall,draw=none,fill=none,inner sep=1pt,
                    row sep=0.4pt,legend columns=1},
      legend cell align=left,
      every axis plot/.append style={line width=0.7pt,mark size=1.2pt},
    },
  }
\fi
}{%
    % ---------------------------------------------------------------------------
% synstyle.tex --- shared look for every TikZ figure of the SynKit CRN paper.
%
% Input it from a figure file:
%     \IfFileExists{synstyle.tex}{% ---------------------------------------------------------------------------
% synstyle.tex --- shared look for every TikZ figure of the SynKit CRN paper.
%
% Input it from a figure file:
%     \IfFileExists{synstyle.tex}{% ---------------------------------------------------------------------------
% synstyle.tex --- shared look for every TikZ figure of the SynKit CRN paper.
%
% Input it from a figure file:
%     \IfFileExists{synstyle.tex}{\input{synstyle}}{\input{Figure/synstyle}}
% so that the figure compiles both from Figure/ and from the paper root.
%
% The visual vocabulary follows ../../../../Style/style.tex: the
% colour-blind-safe Okabe--Ito accents, quiet neutral rules and wash, circular
% panel badges, and a compact technical type scale are shared by the companion
% papers.  Semantic aliases below keep the CRN meaning explicit.
%
% DESIGN CONTRACT -----------------------------------------------------------
% 1. Geometry.  Every figure is laid out on a canvas exactly \scW wide (the
%    sn-jnl text width, 372 pt) and is included unscaled, so one point of type
%    is the same physical size in all five figures.  Open each picture with
%    \scCanvas{<height in cm>} to pin the bounding box.
% 2. Colour is semantic, never decorative.  Four roles only:
%       scAccent (green)       structure, the normal/expected path
%       scMark   (orange)      catalyst, certificate, or highlight
%       scAlert  (vermillion)  failure, violation, or uncovered set
%       scRef    (blue)        secondary or reference series
%    A panel gets at most one accent.  Titles are always ink; panel identity
%    is also encoded by a letter inside a circular badge, not by colour alone.
% 3. Type scale.  Three sizes, no others, and nothing below 6 pt:
%       \scTitle (7 pt bold)  panel titles
%       \scBody  (6.5 pt)     labels, values, in-figure prose
%       \scSmall (6 pt)       units, footnotes, dense graph labels
% ---------------------------------------------------------------------------
% Load guard.  The paper inputs every figure source inline, so this file is
% reached once per figure; without the guard the second \newlength below
% would abort the run.
\ifdefined\SynKitStyleLoaded\endinput\fi
\newcommand{\SynKitStyleLoaded}{}

% Load in the document preamble. XeLaTeX/LuaLaTeX and an installed Arial
% family are required; fontspec reports an error instead of substituting.
\RequirePackage[no-math]{fontspec}
\newfontfamily\scArial{Arial}[NFSSFamily=scArial,Ligatures=TeX]
% Keep manuscript mathematics unchanged; select Arial letters and numerals
% only inside figures. TeX supplies specialist mathematical symbols.
\RequirePackage[subdued,italic]{mathastext}
\MTDeclareVersion{synarial}{TU}{scArial}{m}{n}
\SetMathAlphabet{\Mathsf}{synarial}{TU}{scArial}{m}{n}
\SetMathAlphabet{\Mathtt}{synarial}{TU}{scArial}{m}{n}

\usetikzlibrary{positioning,calc,fit,backgrounds,arrows.meta,shapes.geometric,
                shapes.misc,decorations.pathreplacing,patterns,math}

% --- canvas ---------------------------------------------------------------
% 372 pt text width minus the 1 pt standalone border on each side.
\newlength{\scW}\setlength{\scW}{13.0326cm}
\providecommand{\scWnum}{13.0326}
% \scCanvas{<height>} pins the picture to the full text width, y running down
% from 0, so every figure shares one coordinate convention.  It *sets* the
% bounding box rather than contributing to it, so the emitted page is exactly
% the text width and every figure is included at 1:1 -- otherwise a stray
% invisible extent (pgfplots reserves a little around an axis) rescales one
% figure relative to the others and the same \scSmall prints at two sizes.
% Anything drawn outside the canvas is therefore cropped: keep content inside,
% and check with
%     \providecommand{\scCanvas}[1]{\path[draw=red] (0,0) rectangle (\scWnum,-#1);}
\providecommand{\scCanvas}[1]{%
  \path[use as bounding box] (0,0) rectangle (\scWnum,-#1);}

% --- palette --------------------------------------------------------------
% Canonical companion-paper values (../../../../Style/style.tex).
\definecolor{npgOrange}{HTML}{E69F00}
\definecolor{npgSky}{HTML}{56B4E9}
\definecolor{npgGreen}{HTML}{009E73}
\definecolor{npgYellow}{HTML}{F0E442}
\definecolor{npgBlue}{HTML}{0072B2}
\definecolor{npgVermillion}{HTML}{E64B35}
\definecolor{npgPurple}{HTML}{CC79A7}
\definecolor{figInk}{HTML}{1A1A1A}
\definecolor{figMuted}{HTML}{666666}
\definecolor{figHair}{HTML}{D9D9D9}
\definecolor{figWash}{HTML}{F7F7F7}

\colorlet{scInk}{figInk}              % primary text
\colorlet{scInkM}{figMuted}           % secondary text
\colorlet{scPaper}{figWash}           % quiet cards and note strips
\colorlet{scSpF}{npgSky!12}           % species fill
\colorlet{scSpE}{npgBlue}             % species edge
\colorlet{scRxF}{npgOrange!12}        % reaction fill
\colorlet{scRxE}{npgOrange!88!black}  % reaction edge
\colorlet{scReac}{npgBlue!78!black}   % reactant edge
\colorlet{scProd}{npgGreen!82!black}  % product edge
\colorlet{scHiF}{npgOrange!13}        % highlight fill
\colorlet{scHiE}{npgOrange!82!black}  % highlight edge
% Stage / accent tints.  stDE remains available for the fifth plot series.
\colorlet{stAF}{npgGreen!9}\colorlet{stAE}{npgGreen}
\colorlet{stBF}{npgOrange!11}\colorlet{stBE}{npgOrange}
\colorlet{stCF}{npgBlue!9}\colorlet{stCE}{npgBlue}
\colorlet{stDF}{npgPurple!10}\colorlet{stDE}{npgPurple!85!black}
\colorlet{stEF}{npgVermillion!9}\colorlet{stEE}{npgVermillion}

% Semantic aliases.  Figures should speak in these, not in stAE/stBE/...
\colorlet{scAccent}{stAE}\colorlet{scAccentF}{stAF}   % structure / normal
\colorlet{scMark}{stBE}\colorlet{scMarkF}{stBF}       % highlight / certificate
\colorlet{scRef}{stCE}\colorlet{scRefF}{stCF}         % secondary series
\colorlet{scAlert}{stEE}\colorlet{scAlertF}{stEF}     % failure / violation
\colorlet{scRule}{figHair}                            % card and axis hairlines
\colorlet{scRuleL}{figHair!55}                        % grid, dividers

% --- text styles ----------------------------------------------------------
% Three sizes, fixed in points so they do not drift with the document class.
% Arial is shared by labels, headings, panel badges, and chemistry labels.
\providecommand{\scTitle}{\scArial\fontsize{7}{8.4}\selectfont\bfseries\color{scInk}}
\providecommand{\scBody}{\scArial\fontsize{6.5}{7.8}\selectfont\color{scInkM}}
\providecommand{\scSmall}{\scArial\fontsize{6}{7.2}\selectfont\color{scInkM}}
\providecommand{\scSub}[1]{{\scSmall\itshape #1}}
\providecommand{\scPanel}[1]{{\scTitle(#1)}}
% Numerals that carry a result: ink, not colour, so the eye reads the value.
\providecommand{\scNum}{\scArial\fontsize{10}{11}\selectfont\bfseries\color{scInk}}

% --- shared TikZ styles ---------------------------------------------------
\tikzset{
  % text never hyphenates inside a figure
  every node/.append style={execute at begin node={%
      \hyphenpenalty=10000 \exhyphenpenalty=10000 \hbadness=10000 }},
  syn/.style={font=\scArial,color=scInk,
              execute at begin picture={\MTversion*{synarial}},
              >={Stealth[round,length=2.4mm,width=1.9mm]}},
  % containers ------------------------------------------------------------
  % One card chrome for the whole paper: white, 3 pt radius, one hairline.
  panel/.style={rounded corners=2.5pt,fill=none,draw=scRule,line width=0.45pt},
  card/.style={rounded corners=2.2pt,fill=scPaper,draw=scRule,line width=0.45pt,
               inner sep=5pt,align=left,anchor=north west},
  % a quiet strip for a panel's closing fact
  note/.style={rounded corners=2.5pt,fill=scPaper,draw=scRuleL,line width=0.4pt,
               inner xsep=6pt,inner ysep=3.5pt,align=left,font=\scSmall},
  % small labelled value pill
  chip/.style={rounded corners=2pt,fill=scPaper,draw=scRule,line width=0.4pt,
               inner xsep=4pt,inner ysep=2pt,font=\scSmall},
  % Circular letter badge: the same accessible panel marker used in Style/*.
  hbullet/.style={circle,fill=scInk,text=white,draw=none,
                  font=\scArial\fontsize{6}{6}\selectfont\bfseries,
                  minimum size=4.2mm,inner sep=0pt},
  ctag/.style={font=\scTitle,anchor=north west,text=scInk},
  % CRN primitives --------------------------------------------------------
  species/.style={rounded rectangle,fill=scSpF,draw=scSpE,line width=0.45pt,
                  inner xsep=2.4pt,inner ysep=1.6pt,font=\scSmall,
                  text=scInk,minimum height=3.6mm},
  cofactor/.style={species,fill=scRefF,draw=scRef},
  catalyst/.style={species,fill=scHiF,draw=scHiE},
  rxn/.style={rectangle,fill=scRxF,draw=scRxE,line width=0.45pt,
              inner sep=1.2pt,font=\scSmall,text=scInk,
              minimum width=3.0mm,minimum height=3.0mm},
  place/.style={circle,fill=scSpF,draw=scSpE,line width=0.45pt,
                minimum size=4.0mm,inner sep=0pt},
  trans/.style={rectangle,fill=scRxE,draw=scRxE,minimum width=1.1mm,
                minimum height=3.6mm,inner sep=0pt},
  token/.style={circle,fill=scInk,minimum size=1.1mm,inner sep=0pt},
  % edges -----------------------------------------------------------------
  reac/.style={->,line width=0.5pt,draw=scReac,shorten >=0.8pt,shorten <=0.8pt},
  prod/.style={->,line width=0.5pt,draw=scProd,shorten >=0.8pt,shorten <=0.8pt},
  cat/.style={->,line width=0.5pt,draw=scHiE,shorten >=0.8pt,shorten <=0.8pt},
  flow/.style={-{Stealth[round,length=2.2mm,width=1.8mm]},line width=1.1pt,
               draw=scInkM!45},
  flab/.style={font=\scSmall,text=scInkM,inner sep=2pt},
  stoi/.style={font=\scSmall,text=scInkM,inner sep=0.8pt},
  % verdict marks: shape carries the meaning, colour only reinforces it
  vpass/.style={circle,fill=scAccentF,draw=scAccent,text=scAccent,
                font=\fontsize{6}{6}\selectfont\bfseries,
                minimum size=3.2mm,inner sep=0pt},
  vfail/.style={circle,fill=scAlertF,draw=scAlert,text=scAlert,
                font=\fontsize{6}{6}\selectfont\bfseries,
                minimum size=3.2mm,inner sep=0pt},
  callout/.style={font=\scSmall,text=scInkM,align=left},
}

% \scHead[<badge colour>]{<anchor coord>}{<letter>}{<title>}{<subtitle or empty>}
% One header treatment for every panel in the paper.
\providecommand{\scHead}[5][scInk]{%
  \node[hbullet,fill=#1,anchor=north west] (scbadge#3) at ($#2+(0,-0.015)$)
    {\MakeUppercase{#3}};
  \node[ctag,anchor=west] at ($(scbadge#3.east)+(0.14,0)$) {#4};
  \ifx\relax#5\relax\else
    \node[anchor=north west,font=\scSmall,text=scInkM]
      at ($#2+(0.56,-0.245)$) {#5};
  \fi}

% --- shared pgfplots look (only if the figure loaded pgfplots) ------------
\ifdefined\pgfplotsset
  \pgfplotsset{compat=1.18}
  \pgfplotsset{
    synaxis/.style={
      scale only axis,
      tick align=outside, tick pos=left,
      axis line style={draw=scInkM!55,line width=0.4pt},
      % One grid, on the value axis only.  No minor grid: at these sizes it
      % turns the plot frame into a grey wash.
      ymajorgrids=true, xmajorgrids=false,
      major grid style={draw=scRuleL,line width=0.3pt},
      major tick length=1.6pt,
      label style={font=\scBody,color=scInk},
      tick label style={font=\scSmall,color=scInkM},
      legend style={font=\scSmall,draw=none,fill=none,inner sep=1pt,
                    row sep=0.4pt,legend columns=1},
      legend cell align=left,
      every axis plot/.append style={line width=0.7pt,mark size=1.2pt},
    },
  }
\fi
}{% ---------------------------------------------------------------------------
% synstyle.tex --- shared look for every TikZ figure of the SynKit CRN paper.
%
% Input it from a figure file:
%     \IfFileExists{synstyle.tex}{\input{synstyle}}{\input{Figure/synstyle}}
% so that the figure compiles both from Figure/ and from the paper root.
%
% The visual vocabulary follows ../../../../Style/style.tex: the
% colour-blind-safe Okabe--Ito accents, quiet neutral rules and wash, circular
% panel badges, and a compact technical type scale are shared by the companion
% papers.  Semantic aliases below keep the CRN meaning explicit.
%
% DESIGN CONTRACT -----------------------------------------------------------
% 1. Geometry.  Every figure is laid out on a canvas exactly \scW wide (the
%    sn-jnl text width, 372 pt) and is included unscaled, so one point of type
%    is the same physical size in all five figures.  Open each picture with
%    \scCanvas{<height in cm>} to pin the bounding box.
% 2. Colour is semantic, never decorative.  Four roles only:
%       scAccent (green)       structure, the normal/expected path
%       scMark   (orange)      catalyst, certificate, or highlight
%       scAlert  (vermillion)  failure, violation, or uncovered set
%       scRef    (blue)        secondary or reference series
%    A panel gets at most one accent.  Titles are always ink; panel identity
%    is also encoded by a letter inside a circular badge, not by colour alone.
% 3. Type scale.  Three sizes, no others, and nothing below 6 pt:
%       \scTitle (7 pt bold)  panel titles
%       \scBody  (6.5 pt)     labels, values, in-figure prose
%       \scSmall (6 pt)       units, footnotes, dense graph labels
% ---------------------------------------------------------------------------
% Load guard.  The paper inputs every figure source inline, so this file is
% reached once per figure; without the guard the second \newlength below
% would abort the run.
\ifdefined\SynKitStyleLoaded\endinput\fi
\newcommand{\SynKitStyleLoaded}{}

% Load in the document preamble. XeLaTeX/LuaLaTeX and an installed Arial
% family are required; fontspec reports an error instead of substituting.
\RequirePackage[no-math]{fontspec}
\newfontfamily\scArial{Arial}[NFSSFamily=scArial,Ligatures=TeX]
% Keep manuscript mathematics unchanged; select Arial letters and numerals
% only inside figures. TeX supplies specialist mathematical symbols.
\RequirePackage[subdued,italic]{mathastext}
\MTDeclareVersion{synarial}{TU}{scArial}{m}{n}
\SetMathAlphabet{\Mathsf}{synarial}{TU}{scArial}{m}{n}
\SetMathAlphabet{\Mathtt}{synarial}{TU}{scArial}{m}{n}

\usetikzlibrary{positioning,calc,fit,backgrounds,arrows.meta,shapes.geometric,
                shapes.misc,decorations.pathreplacing,patterns,math}

% --- canvas ---------------------------------------------------------------
% 372 pt text width minus the 1 pt standalone border on each side.
\newlength{\scW}\setlength{\scW}{13.0326cm}
\providecommand{\scWnum}{13.0326}
% \scCanvas{<height>} pins the picture to the full text width, y running down
% from 0, so every figure shares one coordinate convention.  It *sets* the
% bounding box rather than contributing to it, so the emitted page is exactly
% the text width and every figure is included at 1:1 -- otherwise a stray
% invisible extent (pgfplots reserves a little around an axis) rescales one
% figure relative to the others and the same \scSmall prints at two sizes.
% Anything drawn outside the canvas is therefore cropped: keep content inside,
% and check with
%     \providecommand{\scCanvas}[1]{\path[draw=red] (0,0) rectangle (\scWnum,-#1);}
\providecommand{\scCanvas}[1]{%
  \path[use as bounding box] (0,0) rectangle (\scWnum,-#1);}

% --- palette --------------------------------------------------------------
% Canonical companion-paper values (../../../../Style/style.tex).
\definecolor{npgOrange}{HTML}{E69F00}
\definecolor{npgSky}{HTML}{56B4E9}
\definecolor{npgGreen}{HTML}{009E73}
\definecolor{npgYellow}{HTML}{F0E442}
\definecolor{npgBlue}{HTML}{0072B2}
\definecolor{npgVermillion}{HTML}{E64B35}
\definecolor{npgPurple}{HTML}{CC79A7}
\definecolor{figInk}{HTML}{1A1A1A}
\definecolor{figMuted}{HTML}{666666}
\definecolor{figHair}{HTML}{D9D9D9}
\definecolor{figWash}{HTML}{F7F7F7}

\colorlet{scInk}{figInk}              % primary text
\colorlet{scInkM}{figMuted}           % secondary text
\colorlet{scPaper}{figWash}           % quiet cards and note strips
\colorlet{scSpF}{npgSky!12}           % species fill
\colorlet{scSpE}{npgBlue}             % species edge
\colorlet{scRxF}{npgOrange!12}        % reaction fill
\colorlet{scRxE}{npgOrange!88!black}  % reaction edge
\colorlet{scReac}{npgBlue!78!black}   % reactant edge
\colorlet{scProd}{npgGreen!82!black}  % product edge
\colorlet{scHiF}{npgOrange!13}        % highlight fill
\colorlet{scHiE}{npgOrange!82!black}  % highlight edge
% Stage / accent tints.  stDE remains available for the fifth plot series.
\colorlet{stAF}{npgGreen!9}\colorlet{stAE}{npgGreen}
\colorlet{stBF}{npgOrange!11}\colorlet{stBE}{npgOrange}
\colorlet{stCF}{npgBlue!9}\colorlet{stCE}{npgBlue}
\colorlet{stDF}{npgPurple!10}\colorlet{stDE}{npgPurple!85!black}
\colorlet{stEF}{npgVermillion!9}\colorlet{stEE}{npgVermillion}

% Semantic aliases.  Figures should speak in these, not in stAE/stBE/...
\colorlet{scAccent}{stAE}\colorlet{scAccentF}{stAF}   % structure / normal
\colorlet{scMark}{stBE}\colorlet{scMarkF}{stBF}       % highlight / certificate
\colorlet{scRef}{stCE}\colorlet{scRefF}{stCF}         % secondary series
\colorlet{scAlert}{stEE}\colorlet{scAlertF}{stEF}     % failure / violation
\colorlet{scRule}{figHair}                            % card and axis hairlines
\colorlet{scRuleL}{figHair!55}                        % grid, dividers

% --- text styles ----------------------------------------------------------
% Three sizes, fixed in points so they do not drift with the document class.
% Arial is shared by labels, headings, panel badges, and chemistry labels.
\providecommand{\scTitle}{\scArial\fontsize{7}{8.4}\selectfont\bfseries\color{scInk}}
\providecommand{\scBody}{\scArial\fontsize{6.5}{7.8}\selectfont\color{scInkM}}
\providecommand{\scSmall}{\scArial\fontsize{6}{7.2}\selectfont\color{scInkM}}
\providecommand{\scSub}[1]{{\scSmall\itshape #1}}
\providecommand{\scPanel}[1]{{\scTitle(#1)}}
% Numerals that carry a result: ink, not colour, so the eye reads the value.
\providecommand{\scNum}{\scArial\fontsize{10}{11}\selectfont\bfseries\color{scInk}}

% --- shared TikZ styles ---------------------------------------------------
\tikzset{
  % text never hyphenates inside a figure
  every node/.append style={execute at begin node={%
      \hyphenpenalty=10000 \exhyphenpenalty=10000 \hbadness=10000 }},
  syn/.style={font=\scArial,color=scInk,
              execute at begin picture={\MTversion*{synarial}},
              >={Stealth[round,length=2.4mm,width=1.9mm]}},
  % containers ------------------------------------------------------------
  % One card chrome for the whole paper: white, 3 pt radius, one hairline.
  panel/.style={rounded corners=2.5pt,fill=none,draw=scRule,line width=0.45pt},
  card/.style={rounded corners=2.2pt,fill=scPaper,draw=scRule,line width=0.45pt,
               inner sep=5pt,align=left,anchor=north west},
  % a quiet strip for a panel's closing fact
  note/.style={rounded corners=2.5pt,fill=scPaper,draw=scRuleL,line width=0.4pt,
               inner xsep=6pt,inner ysep=3.5pt,align=left,font=\scSmall},
  % small labelled value pill
  chip/.style={rounded corners=2pt,fill=scPaper,draw=scRule,line width=0.4pt,
               inner xsep=4pt,inner ysep=2pt,font=\scSmall},
  % Circular letter badge: the same accessible panel marker used in Style/*.
  hbullet/.style={circle,fill=scInk,text=white,draw=none,
                  font=\scArial\fontsize{6}{6}\selectfont\bfseries,
                  minimum size=4.2mm,inner sep=0pt},
  ctag/.style={font=\scTitle,anchor=north west,text=scInk},
  % CRN primitives --------------------------------------------------------
  species/.style={rounded rectangle,fill=scSpF,draw=scSpE,line width=0.45pt,
                  inner xsep=2.4pt,inner ysep=1.6pt,font=\scSmall,
                  text=scInk,minimum height=3.6mm},
  cofactor/.style={species,fill=scRefF,draw=scRef},
  catalyst/.style={species,fill=scHiF,draw=scHiE},
  rxn/.style={rectangle,fill=scRxF,draw=scRxE,line width=0.45pt,
              inner sep=1.2pt,font=\scSmall,text=scInk,
              minimum width=3.0mm,minimum height=3.0mm},
  place/.style={circle,fill=scSpF,draw=scSpE,line width=0.45pt,
                minimum size=4.0mm,inner sep=0pt},
  trans/.style={rectangle,fill=scRxE,draw=scRxE,minimum width=1.1mm,
                minimum height=3.6mm,inner sep=0pt},
  token/.style={circle,fill=scInk,minimum size=1.1mm,inner sep=0pt},
  % edges -----------------------------------------------------------------
  reac/.style={->,line width=0.5pt,draw=scReac,shorten >=0.8pt,shorten <=0.8pt},
  prod/.style={->,line width=0.5pt,draw=scProd,shorten >=0.8pt,shorten <=0.8pt},
  cat/.style={->,line width=0.5pt,draw=scHiE,shorten >=0.8pt,shorten <=0.8pt},
  flow/.style={-{Stealth[round,length=2.2mm,width=1.8mm]},line width=1.1pt,
               draw=scInkM!45},
  flab/.style={font=\scSmall,text=scInkM,inner sep=2pt},
  stoi/.style={font=\scSmall,text=scInkM,inner sep=0.8pt},
  % verdict marks: shape carries the meaning, colour only reinforces it
  vpass/.style={circle,fill=scAccentF,draw=scAccent,text=scAccent,
                font=\fontsize{6}{6}\selectfont\bfseries,
                minimum size=3.2mm,inner sep=0pt},
  vfail/.style={circle,fill=scAlertF,draw=scAlert,text=scAlert,
                font=\fontsize{6}{6}\selectfont\bfseries,
                minimum size=3.2mm,inner sep=0pt},
  callout/.style={font=\scSmall,text=scInkM,align=left},
}

% \scHead[<badge colour>]{<anchor coord>}{<letter>}{<title>}{<subtitle or empty>}
% One header treatment for every panel in the paper.
\providecommand{\scHead}[5][scInk]{%
  \node[hbullet,fill=#1,anchor=north west] (scbadge#3) at ($#2+(0,-0.015)$)
    {\MakeUppercase{#3}};
  \node[ctag,anchor=west] at ($(scbadge#3.east)+(0.14,0)$) {#4};
  \ifx\relax#5\relax\else
    \node[anchor=north west,font=\scSmall,text=scInkM]
      at ($#2+(0.56,-0.245)$) {#5};
  \fi}

% --- shared pgfplots look (only if the figure loaded pgfplots) ------------
\ifdefined\pgfplotsset
  \pgfplotsset{compat=1.18}
  \pgfplotsset{
    synaxis/.style={
      scale only axis,
      tick align=outside, tick pos=left,
      axis line style={draw=scInkM!55,line width=0.4pt},
      % One grid, on the value axis only.  No minor grid: at these sizes it
      % turns the plot frame into a grey wash.
      ymajorgrids=true, xmajorgrids=false,
      major grid style={draw=scRuleL,line width=0.3pt},
      major tick length=1.6pt,
      label style={font=\scBody,color=scInk},
      tick label style={font=\scSmall,color=scInkM},
      legend style={font=\scSmall,draw=none,fill=none,inner sep=1pt,
                    row sep=0.4pt,legend columns=1},
      legend cell align=left,
      every axis plot/.append style={line width=0.7pt,mark size=1.2pt},
    },
  }
\fi
}
% so that the figure compiles both from Figure/ and from the paper root.
%
% The visual vocabulary follows ../../../../Style/style.tex: the
% colour-blind-safe Okabe--Ito accents, quiet neutral rules and wash, circular
% panel badges, and a compact technical type scale are shared by the companion
% papers.  Semantic aliases below keep the CRN meaning explicit.
%
% DESIGN CONTRACT -----------------------------------------------------------
% 1. Geometry.  Every figure is laid out on a canvas exactly \scW wide (the
%    sn-jnl text width, 372 pt) and is included unscaled, so one point of type
%    is the same physical size in all five figures.  Open each picture with
%    \scCanvas{<height in cm>} to pin the bounding box.
% 2. Colour is semantic, never decorative.  Four roles only:
%       scAccent (green)       structure, the normal/expected path
%       scMark   (orange)      catalyst, certificate, or highlight
%       scAlert  (vermillion)  failure, violation, or uncovered set
%       scRef    (blue)        secondary or reference series
%    A panel gets at most one accent.  Titles are always ink; panel identity
%    is also encoded by a letter inside a circular badge, not by colour alone.
% 3. Type scale.  Three sizes, no others, and nothing below 6 pt:
%       \scTitle (7 pt bold)  panel titles
%       \scBody  (6.5 pt)     labels, values, in-figure prose
%       \scSmall (6 pt)       units, footnotes, dense graph labels
% ---------------------------------------------------------------------------
% Load guard.  The paper inputs every figure source inline, so this file is
% reached once per figure; without the guard the second \newlength below
% would abort the run.
\ifdefined\SynKitStyleLoaded\endinput\fi
\newcommand{\SynKitStyleLoaded}{}

% Load in the document preamble. XeLaTeX/LuaLaTeX and an installed Arial
% family are required; fontspec reports an error instead of substituting.
\RequirePackage[no-math]{fontspec}
\newfontfamily\scArial{Arial}[NFSSFamily=scArial,Ligatures=TeX]
% Keep manuscript mathematics unchanged; select Arial letters and numerals
% only inside figures. TeX supplies specialist mathematical symbols.
\RequirePackage[subdued,italic]{mathastext}
\MTDeclareVersion{synarial}{TU}{scArial}{m}{n}
\SetMathAlphabet{\Mathsf}{synarial}{TU}{scArial}{m}{n}
\SetMathAlphabet{\Mathtt}{synarial}{TU}{scArial}{m}{n}

\usetikzlibrary{positioning,calc,fit,backgrounds,arrows.meta,shapes.geometric,
                shapes.misc,decorations.pathreplacing,patterns,math}

% --- canvas ---------------------------------------------------------------
% 372 pt text width minus the 1 pt standalone border on each side.
\newlength{\scW}\setlength{\scW}{13.0326cm}
\providecommand{\scWnum}{13.0326}
% \scCanvas{<height>} pins the picture to the full text width, y running down
% from 0, so every figure shares one coordinate convention.  It *sets* the
% bounding box rather than contributing to it, so the emitted page is exactly
% the text width and every figure is included at 1:1 -- otherwise a stray
% invisible extent (pgfplots reserves a little around an axis) rescales one
% figure relative to the others and the same \scSmall prints at two sizes.
% Anything drawn outside the canvas is therefore cropped: keep content inside,
% and check with
%     \providecommand{\scCanvas}[1]{\path[draw=red] (0,0) rectangle (\scWnum,-#1);}
\providecommand{\scCanvas}[1]{%
  \path[use as bounding box] (0,0) rectangle (\scWnum,-#1);}

% --- palette --------------------------------------------------------------
% Canonical companion-paper values (../../../../Style/style.tex).
\definecolor{npgOrange}{HTML}{E69F00}
\definecolor{npgSky}{HTML}{56B4E9}
\definecolor{npgGreen}{HTML}{009E73}
\definecolor{npgYellow}{HTML}{F0E442}
\definecolor{npgBlue}{HTML}{0072B2}
\definecolor{npgVermillion}{HTML}{E64B35}
\definecolor{npgPurple}{HTML}{CC79A7}
\definecolor{figInk}{HTML}{1A1A1A}
\definecolor{figMuted}{HTML}{666666}
\definecolor{figHair}{HTML}{D9D9D9}
\definecolor{figWash}{HTML}{F7F7F7}

\colorlet{scInk}{figInk}              % primary text
\colorlet{scInkM}{figMuted}           % secondary text
\colorlet{scPaper}{figWash}           % quiet cards and note strips
\colorlet{scSpF}{npgSky!12}           % species fill
\colorlet{scSpE}{npgBlue}             % species edge
\colorlet{scRxF}{npgOrange!12}        % reaction fill
\colorlet{scRxE}{npgOrange!88!black}  % reaction edge
\colorlet{scReac}{npgBlue!78!black}   % reactant edge
\colorlet{scProd}{npgGreen!82!black}  % product edge
\colorlet{scHiF}{npgOrange!13}        % highlight fill
\colorlet{scHiE}{npgOrange!82!black}  % highlight edge
% Stage / accent tints.  stDE remains available for the fifth plot series.
\colorlet{stAF}{npgGreen!9}\colorlet{stAE}{npgGreen}
\colorlet{stBF}{npgOrange!11}\colorlet{stBE}{npgOrange}
\colorlet{stCF}{npgBlue!9}\colorlet{stCE}{npgBlue}
\colorlet{stDF}{npgPurple!10}\colorlet{stDE}{npgPurple!85!black}
\colorlet{stEF}{npgVermillion!9}\colorlet{stEE}{npgVermillion}

% Semantic aliases.  Figures should speak in these, not in stAE/stBE/...
\colorlet{scAccent}{stAE}\colorlet{scAccentF}{stAF}   % structure / normal
\colorlet{scMark}{stBE}\colorlet{scMarkF}{stBF}       % highlight / certificate
\colorlet{scRef}{stCE}\colorlet{scRefF}{stCF}         % secondary series
\colorlet{scAlert}{stEE}\colorlet{scAlertF}{stEF}     % failure / violation
\colorlet{scRule}{figHair}                            % card and axis hairlines
\colorlet{scRuleL}{figHair!55}                        % grid, dividers

% --- text styles ----------------------------------------------------------
% Three sizes, fixed in points so they do not drift with the document class.
% Arial is shared by labels, headings, panel badges, and chemistry labels.
\providecommand{\scTitle}{\scArial\fontsize{7}{8.4}\selectfont\bfseries\color{scInk}}
\providecommand{\scBody}{\scArial\fontsize{6.5}{7.8}\selectfont\color{scInkM}}
\providecommand{\scSmall}{\scArial\fontsize{6}{7.2}\selectfont\color{scInkM}}
\providecommand{\scSub}[1]{{\scSmall\itshape #1}}
\providecommand{\scPanel}[1]{{\scTitle(#1)}}
% Numerals that carry a result: ink, not colour, so the eye reads the value.
\providecommand{\scNum}{\scArial\fontsize{10}{11}\selectfont\bfseries\color{scInk}}

% --- shared TikZ styles ---------------------------------------------------
\tikzset{
  % text never hyphenates inside a figure
  every node/.append style={execute at begin node={%
      \hyphenpenalty=10000 \exhyphenpenalty=10000 \hbadness=10000 }},
  syn/.style={font=\scArial,color=scInk,
              execute at begin picture={\MTversion*{synarial}},
              >={Stealth[round,length=2.4mm,width=1.9mm]}},
  % containers ------------------------------------------------------------
  % One card chrome for the whole paper: white, 3 pt radius, one hairline.
  panel/.style={rounded corners=2.5pt,fill=none,draw=scRule,line width=0.45pt},
  card/.style={rounded corners=2.2pt,fill=scPaper,draw=scRule,line width=0.45pt,
               inner sep=5pt,align=left,anchor=north west},
  % a quiet strip for a panel's closing fact
  note/.style={rounded corners=2.5pt,fill=scPaper,draw=scRuleL,line width=0.4pt,
               inner xsep=6pt,inner ysep=3.5pt,align=left,font=\scSmall},
  % small labelled value pill
  chip/.style={rounded corners=2pt,fill=scPaper,draw=scRule,line width=0.4pt,
               inner xsep=4pt,inner ysep=2pt,font=\scSmall},
  % Circular letter badge: the same accessible panel marker used in Style/*.
  hbullet/.style={circle,fill=scInk,text=white,draw=none,
                  font=\scArial\fontsize{6}{6}\selectfont\bfseries,
                  minimum size=4.2mm,inner sep=0pt},
  ctag/.style={font=\scTitle,anchor=north west,text=scInk},
  % CRN primitives --------------------------------------------------------
  species/.style={rounded rectangle,fill=scSpF,draw=scSpE,line width=0.45pt,
                  inner xsep=2.4pt,inner ysep=1.6pt,font=\scSmall,
                  text=scInk,minimum height=3.6mm},
  cofactor/.style={species,fill=scRefF,draw=scRef},
  catalyst/.style={species,fill=scHiF,draw=scHiE},
  rxn/.style={rectangle,fill=scRxF,draw=scRxE,line width=0.45pt,
              inner sep=1.2pt,font=\scSmall,text=scInk,
              minimum width=3.0mm,minimum height=3.0mm},
  place/.style={circle,fill=scSpF,draw=scSpE,line width=0.45pt,
                minimum size=4.0mm,inner sep=0pt},
  trans/.style={rectangle,fill=scRxE,draw=scRxE,minimum width=1.1mm,
                minimum height=3.6mm,inner sep=0pt},
  token/.style={circle,fill=scInk,minimum size=1.1mm,inner sep=0pt},
  % edges -----------------------------------------------------------------
  reac/.style={->,line width=0.5pt,draw=scReac,shorten >=0.8pt,shorten <=0.8pt},
  prod/.style={->,line width=0.5pt,draw=scProd,shorten >=0.8pt,shorten <=0.8pt},
  cat/.style={->,line width=0.5pt,draw=scHiE,shorten >=0.8pt,shorten <=0.8pt},
  flow/.style={-{Stealth[round,length=2.2mm,width=1.8mm]},line width=1.1pt,
               draw=scInkM!45},
  flab/.style={font=\scSmall,text=scInkM,inner sep=2pt},
  stoi/.style={font=\scSmall,text=scInkM,inner sep=0.8pt},
  % verdict marks: shape carries the meaning, colour only reinforces it
  vpass/.style={circle,fill=scAccentF,draw=scAccent,text=scAccent,
                font=\fontsize{6}{6}\selectfont\bfseries,
                minimum size=3.2mm,inner sep=0pt},
  vfail/.style={circle,fill=scAlertF,draw=scAlert,text=scAlert,
                font=\fontsize{6}{6}\selectfont\bfseries,
                minimum size=3.2mm,inner sep=0pt},
  callout/.style={font=\scSmall,text=scInkM,align=left},
}

% \scHead[<badge colour>]{<anchor coord>}{<letter>}{<title>}{<subtitle or empty>}
% One header treatment for every panel in the paper.
\providecommand{\scHead}[5][scInk]{%
  \node[hbullet,fill=#1,anchor=north west] (scbadge#3) at ($#2+(0,-0.015)$)
    {\MakeUppercase{#3}};
  \node[ctag,anchor=west] at ($(scbadge#3.east)+(0.14,0)$) {#4};
  \ifx\relax#5\relax\else
    \node[anchor=north west,font=\scSmall,text=scInkM]
      at ($#2+(0.56,-0.245)$) {#5};
  \fi}

% --- shared pgfplots look (only if the figure loaded pgfplots) ------------
\ifdefined\pgfplotsset
  \pgfplotsset{compat=1.18}
  \pgfplotsset{
    synaxis/.style={
      scale only axis,
      tick align=outside, tick pos=left,
      axis line style={draw=scInkM!55,line width=0.4pt},
      % One grid, on the value axis only.  No minor grid: at these sizes it
      % turns the plot frame into a grey wash.
      ymajorgrids=true, xmajorgrids=false,
      major grid style={draw=scRuleL,line width=0.3pt},
      major tick length=1.6pt,
      label style={font=\scBody,color=scInk},
      tick label style={font=\scSmall,color=scInkM},
      legend style={font=\scSmall,draw=none,fill=none,inner sep=1pt,
                    row sep=0.4pt,legend columns=1},
      legend cell align=left,
      every axis plot/.append style={line width=0.7pt,mark size=1.2pt},
    },
  }
\fi
}{% ---------------------------------------------------------------------------
% synstyle.tex --- shared look for every TikZ figure of the SynKit CRN paper.
%
% Input it from a figure file:
%     \IfFileExists{synstyle.tex}{% ---------------------------------------------------------------------------
% synstyle.tex --- shared look for every TikZ figure of the SynKit CRN paper.
%
% Input it from a figure file:
%     \IfFileExists{synstyle.tex}{\input{synstyle}}{\input{Figure/synstyle}}
% so that the figure compiles both from Figure/ and from the paper root.
%
% The visual vocabulary follows ../../../../Style/style.tex: the
% colour-blind-safe Okabe--Ito accents, quiet neutral rules and wash, circular
% panel badges, and a compact technical type scale are shared by the companion
% papers.  Semantic aliases below keep the CRN meaning explicit.
%
% DESIGN CONTRACT -----------------------------------------------------------
% 1. Geometry.  Every figure is laid out on a canvas exactly \scW wide (the
%    sn-jnl text width, 372 pt) and is included unscaled, so one point of type
%    is the same physical size in all five figures.  Open each picture with
%    \scCanvas{<height in cm>} to pin the bounding box.
% 2. Colour is semantic, never decorative.  Four roles only:
%       scAccent (green)       structure, the normal/expected path
%       scMark   (orange)      catalyst, certificate, or highlight
%       scAlert  (vermillion)  failure, violation, or uncovered set
%       scRef    (blue)        secondary or reference series
%    A panel gets at most one accent.  Titles are always ink; panel identity
%    is also encoded by a letter inside a circular badge, not by colour alone.
% 3. Type scale.  Three sizes, no others, and nothing below 6 pt:
%       \scTitle (7 pt bold)  panel titles
%       \scBody  (6.5 pt)     labels, values, in-figure prose
%       \scSmall (6 pt)       units, footnotes, dense graph labels
% ---------------------------------------------------------------------------
% Load guard.  The paper inputs every figure source inline, so this file is
% reached once per figure; without the guard the second \newlength below
% would abort the run.
\ifdefined\SynKitStyleLoaded\endinput\fi
\newcommand{\SynKitStyleLoaded}{}

% Load in the document preamble. XeLaTeX/LuaLaTeX and an installed Arial
% family are required; fontspec reports an error instead of substituting.
\RequirePackage[no-math]{fontspec}
\newfontfamily\scArial{Arial}[NFSSFamily=scArial,Ligatures=TeX]
% Keep manuscript mathematics unchanged; select Arial letters and numerals
% only inside figures. TeX supplies specialist mathematical symbols.
\RequirePackage[subdued,italic]{mathastext}
\MTDeclareVersion{synarial}{TU}{scArial}{m}{n}
\SetMathAlphabet{\Mathsf}{synarial}{TU}{scArial}{m}{n}
\SetMathAlphabet{\Mathtt}{synarial}{TU}{scArial}{m}{n}

\usetikzlibrary{positioning,calc,fit,backgrounds,arrows.meta,shapes.geometric,
                shapes.misc,decorations.pathreplacing,patterns,math}

% --- canvas ---------------------------------------------------------------
% 372 pt text width minus the 1 pt standalone border on each side.
\newlength{\scW}\setlength{\scW}{13.0326cm}
\providecommand{\scWnum}{13.0326}
% \scCanvas{<height>} pins the picture to the full text width, y running down
% from 0, so every figure shares one coordinate convention.  It *sets* the
% bounding box rather than contributing to it, so the emitted page is exactly
% the text width and every figure is included at 1:1 -- otherwise a stray
% invisible extent (pgfplots reserves a little around an axis) rescales one
% figure relative to the others and the same \scSmall prints at two sizes.
% Anything drawn outside the canvas is therefore cropped: keep content inside,
% and check with
%     \providecommand{\scCanvas}[1]{\path[draw=red] (0,0) rectangle (\scWnum,-#1);}
\providecommand{\scCanvas}[1]{%
  \path[use as bounding box] (0,0) rectangle (\scWnum,-#1);}

% --- palette --------------------------------------------------------------
% Canonical companion-paper values (../../../../Style/style.tex).
\definecolor{npgOrange}{HTML}{E69F00}
\definecolor{npgSky}{HTML}{56B4E9}
\definecolor{npgGreen}{HTML}{009E73}
\definecolor{npgYellow}{HTML}{F0E442}
\definecolor{npgBlue}{HTML}{0072B2}
\definecolor{npgVermillion}{HTML}{E64B35}
\definecolor{npgPurple}{HTML}{CC79A7}
\definecolor{figInk}{HTML}{1A1A1A}
\definecolor{figMuted}{HTML}{666666}
\definecolor{figHair}{HTML}{D9D9D9}
\definecolor{figWash}{HTML}{F7F7F7}

\colorlet{scInk}{figInk}              % primary text
\colorlet{scInkM}{figMuted}           % secondary text
\colorlet{scPaper}{figWash}           % quiet cards and note strips
\colorlet{scSpF}{npgSky!12}           % species fill
\colorlet{scSpE}{npgBlue}             % species edge
\colorlet{scRxF}{npgOrange!12}        % reaction fill
\colorlet{scRxE}{npgOrange!88!black}  % reaction edge
\colorlet{scReac}{npgBlue!78!black}   % reactant edge
\colorlet{scProd}{npgGreen!82!black}  % product edge
\colorlet{scHiF}{npgOrange!13}        % highlight fill
\colorlet{scHiE}{npgOrange!82!black}  % highlight edge
% Stage / accent tints.  stDE remains available for the fifth plot series.
\colorlet{stAF}{npgGreen!9}\colorlet{stAE}{npgGreen}
\colorlet{stBF}{npgOrange!11}\colorlet{stBE}{npgOrange}
\colorlet{stCF}{npgBlue!9}\colorlet{stCE}{npgBlue}
\colorlet{stDF}{npgPurple!10}\colorlet{stDE}{npgPurple!85!black}
\colorlet{stEF}{npgVermillion!9}\colorlet{stEE}{npgVermillion}

% Semantic aliases.  Figures should speak in these, not in stAE/stBE/...
\colorlet{scAccent}{stAE}\colorlet{scAccentF}{stAF}   % structure / normal
\colorlet{scMark}{stBE}\colorlet{scMarkF}{stBF}       % highlight / certificate
\colorlet{scRef}{stCE}\colorlet{scRefF}{stCF}         % secondary series
\colorlet{scAlert}{stEE}\colorlet{scAlertF}{stEF}     % failure / violation
\colorlet{scRule}{figHair}                            % card and axis hairlines
\colorlet{scRuleL}{figHair!55}                        % grid, dividers

% --- text styles ----------------------------------------------------------
% Three sizes, fixed in points so they do not drift with the document class.
% Arial is shared by labels, headings, panel badges, and chemistry labels.
\providecommand{\scTitle}{\scArial\fontsize{7}{8.4}\selectfont\bfseries\color{scInk}}
\providecommand{\scBody}{\scArial\fontsize{6.5}{7.8}\selectfont\color{scInkM}}
\providecommand{\scSmall}{\scArial\fontsize{6}{7.2}\selectfont\color{scInkM}}
\providecommand{\scSub}[1]{{\scSmall\itshape #1}}
\providecommand{\scPanel}[1]{{\scTitle(#1)}}
% Numerals that carry a result: ink, not colour, so the eye reads the value.
\providecommand{\scNum}{\scArial\fontsize{10}{11}\selectfont\bfseries\color{scInk}}

% --- shared TikZ styles ---------------------------------------------------
\tikzset{
  % text never hyphenates inside a figure
  every node/.append style={execute at begin node={%
      \hyphenpenalty=10000 \exhyphenpenalty=10000 \hbadness=10000 }},
  syn/.style={font=\scArial,color=scInk,
              execute at begin picture={\MTversion*{synarial}},
              >={Stealth[round,length=2.4mm,width=1.9mm]}},
  % containers ------------------------------------------------------------
  % One card chrome for the whole paper: white, 3 pt radius, one hairline.
  panel/.style={rounded corners=2.5pt,fill=none,draw=scRule,line width=0.45pt},
  card/.style={rounded corners=2.2pt,fill=scPaper,draw=scRule,line width=0.45pt,
               inner sep=5pt,align=left,anchor=north west},
  % a quiet strip for a panel's closing fact
  note/.style={rounded corners=2.5pt,fill=scPaper,draw=scRuleL,line width=0.4pt,
               inner xsep=6pt,inner ysep=3.5pt,align=left,font=\scSmall},
  % small labelled value pill
  chip/.style={rounded corners=2pt,fill=scPaper,draw=scRule,line width=0.4pt,
               inner xsep=4pt,inner ysep=2pt,font=\scSmall},
  % Circular letter badge: the same accessible panel marker used in Style/*.
  hbullet/.style={circle,fill=scInk,text=white,draw=none,
                  font=\scArial\fontsize{6}{6}\selectfont\bfseries,
                  minimum size=4.2mm,inner sep=0pt},
  ctag/.style={font=\scTitle,anchor=north west,text=scInk},
  % CRN primitives --------------------------------------------------------
  species/.style={rounded rectangle,fill=scSpF,draw=scSpE,line width=0.45pt,
                  inner xsep=2.4pt,inner ysep=1.6pt,font=\scSmall,
                  text=scInk,minimum height=3.6mm},
  cofactor/.style={species,fill=scRefF,draw=scRef},
  catalyst/.style={species,fill=scHiF,draw=scHiE},
  rxn/.style={rectangle,fill=scRxF,draw=scRxE,line width=0.45pt,
              inner sep=1.2pt,font=\scSmall,text=scInk,
              minimum width=3.0mm,minimum height=3.0mm},
  place/.style={circle,fill=scSpF,draw=scSpE,line width=0.45pt,
                minimum size=4.0mm,inner sep=0pt},
  trans/.style={rectangle,fill=scRxE,draw=scRxE,minimum width=1.1mm,
                minimum height=3.6mm,inner sep=0pt},
  token/.style={circle,fill=scInk,minimum size=1.1mm,inner sep=0pt},
  % edges -----------------------------------------------------------------
  reac/.style={->,line width=0.5pt,draw=scReac,shorten >=0.8pt,shorten <=0.8pt},
  prod/.style={->,line width=0.5pt,draw=scProd,shorten >=0.8pt,shorten <=0.8pt},
  cat/.style={->,line width=0.5pt,draw=scHiE,shorten >=0.8pt,shorten <=0.8pt},
  flow/.style={-{Stealth[round,length=2.2mm,width=1.8mm]},line width=1.1pt,
               draw=scInkM!45},
  flab/.style={font=\scSmall,text=scInkM,inner sep=2pt},
  stoi/.style={font=\scSmall,text=scInkM,inner sep=0.8pt},
  % verdict marks: shape carries the meaning, colour only reinforces it
  vpass/.style={circle,fill=scAccentF,draw=scAccent,text=scAccent,
                font=\fontsize{6}{6}\selectfont\bfseries,
                minimum size=3.2mm,inner sep=0pt},
  vfail/.style={circle,fill=scAlertF,draw=scAlert,text=scAlert,
                font=\fontsize{6}{6}\selectfont\bfseries,
                minimum size=3.2mm,inner sep=0pt},
  callout/.style={font=\scSmall,text=scInkM,align=left},
}

% \scHead[<badge colour>]{<anchor coord>}{<letter>}{<title>}{<subtitle or empty>}
% One header treatment for every panel in the paper.
\providecommand{\scHead}[5][scInk]{%
  \node[hbullet,fill=#1,anchor=north west] (scbadge#3) at ($#2+(0,-0.015)$)
    {\MakeUppercase{#3}};
  \node[ctag,anchor=west] at ($(scbadge#3.east)+(0.14,0)$) {#4};
  \ifx\relax#5\relax\else
    \node[anchor=north west,font=\scSmall,text=scInkM]
      at ($#2+(0.56,-0.245)$) {#5};
  \fi}

% --- shared pgfplots look (only if the figure loaded pgfplots) ------------
\ifdefined\pgfplotsset
  \pgfplotsset{compat=1.18}
  \pgfplotsset{
    synaxis/.style={
      scale only axis,
      tick align=outside, tick pos=left,
      axis line style={draw=scInkM!55,line width=0.4pt},
      % One grid, on the value axis only.  No minor grid: at these sizes it
      % turns the plot frame into a grey wash.
      ymajorgrids=true, xmajorgrids=false,
      major grid style={draw=scRuleL,line width=0.3pt},
      major tick length=1.6pt,
      label style={font=\scBody,color=scInk},
      tick label style={font=\scSmall,color=scInkM},
      legend style={font=\scSmall,draw=none,fill=none,inner sep=1pt,
                    row sep=0.4pt,legend columns=1},
      legend cell align=left,
      every axis plot/.append style={line width=0.7pt,mark size=1.2pt},
    },
  }
\fi
}{% ---------------------------------------------------------------------------
% synstyle.tex --- shared look for every TikZ figure of the SynKit CRN paper.
%
% Input it from a figure file:
%     \IfFileExists{synstyle.tex}{\input{synstyle}}{\input{Figure/synstyle}}
% so that the figure compiles both from Figure/ and from the paper root.
%
% The visual vocabulary follows ../../../../Style/style.tex: the
% colour-blind-safe Okabe--Ito accents, quiet neutral rules and wash, circular
% panel badges, and a compact technical type scale are shared by the companion
% papers.  Semantic aliases below keep the CRN meaning explicit.
%
% DESIGN CONTRACT -----------------------------------------------------------
% 1. Geometry.  Every figure is laid out on a canvas exactly \scW wide (the
%    sn-jnl text width, 372 pt) and is included unscaled, so one point of type
%    is the same physical size in all five figures.  Open each picture with
%    \scCanvas{<height in cm>} to pin the bounding box.
% 2. Colour is semantic, never decorative.  Four roles only:
%       scAccent (green)       structure, the normal/expected path
%       scMark   (orange)      catalyst, certificate, or highlight
%       scAlert  (vermillion)  failure, violation, or uncovered set
%       scRef    (blue)        secondary or reference series
%    A panel gets at most one accent.  Titles are always ink; panel identity
%    is also encoded by a letter inside a circular badge, not by colour alone.
% 3. Type scale.  Three sizes, no others, and nothing below 6 pt:
%       \scTitle (7 pt bold)  panel titles
%       \scBody  (6.5 pt)     labels, values, in-figure prose
%       \scSmall (6 pt)       units, footnotes, dense graph labels
% ---------------------------------------------------------------------------
% Load guard.  The paper inputs every figure source inline, so this file is
% reached once per figure; without the guard the second \newlength below
% would abort the run.
\ifdefined\SynKitStyleLoaded\endinput\fi
\newcommand{\SynKitStyleLoaded}{}

% Load in the document preamble. XeLaTeX/LuaLaTeX and an installed Arial
% family are required; fontspec reports an error instead of substituting.
\RequirePackage[no-math]{fontspec}
\newfontfamily\scArial{Arial}[NFSSFamily=scArial,Ligatures=TeX]
% Keep manuscript mathematics unchanged; select Arial letters and numerals
% only inside figures. TeX supplies specialist mathematical symbols.
\RequirePackage[subdued,italic]{mathastext}
\MTDeclareVersion{synarial}{TU}{scArial}{m}{n}
\SetMathAlphabet{\Mathsf}{synarial}{TU}{scArial}{m}{n}
\SetMathAlphabet{\Mathtt}{synarial}{TU}{scArial}{m}{n}

\usetikzlibrary{positioning,calc,fit,backgrounds,arrows.meta,shapes.geometric,
                shapes.misc,decorations.pathreplacing,patterns,math}

% --- canvas ---------------------------------------------------------------
% 372 pt text width minus the 1 pt standalone border on each side.
\newlength{\scW}\setlength{\scW}{13.0326cm}
\providecommand{\scWnum}{13.0326}
% \scCanvas{<height>} pins the picture to the full text width, y running down
% from 0, so every figure shares one coordinate convention.  It *sets* the
% bounding box rather than contributing to it, so the emitted page is exactly
% the text width and every figure is included at 1:1 -- otherwise a stray
% invisible extent (pgfplots reserves a little around an axis) rescales one
% figure relative to the others and the same \scSmall prints at two sizes.
% Anything drawn outside the canvas is therefore cropped: keep content inside,
% and check with
%     \providecommand{\scCanvas}[1]{\path[draw=red] (0,0) rectangle (\scWnum,-#1);}
\providecommand{\scCanvas}[1]{%
  \path[use as bounding box] (0,0) rectangle (\scWnum,-#1);}

% --- palette --------------------------------------------------------------
% Canonical companion-paper values (../../../../Style/style.tex).
\definecolor{npgOrange}{HTML}{E69F00}
\definecolor{npgSky}{HTML}{56B4E9}
\definecolor{npgGreen}{HTML}{009E73}
\definecolor{npgYellow}{HTML}{F0E442}
\definecolor{npgBlue}{HTML}{0072B2}
\definecolor{npgVermillion}{HTML}{E64B35}
\definecolor{npgPurple}{HTML}{CC79A7}
\definecolor{figInk}{HTML}{1A1A1A}
\definecolor{figMuted}{HTML}{666666}
\definecolor{figHair}{HTML}{D9D9D9}
\definecolor{figWash}{HTML}{F7F7F7}

\colorlet{scInk}{figInk}              % primary text
\colorlet{scInkM}{figMuted}           % secondary text
\colorlet{scPaper}{figWash}           % quiet cards and note strips
\colorlet{scSpF}{npgSky!12}           % species fill
\colorlet{scSpE}{npgBlue}             % species edge
\colorlet{scRxF}{npgOrange!12}        % reaction fill
\colorlet{scRxE}{npgOrange!88!black}  % reaction edge
\colorlet{scReac}{npgBlue!78!black}   % reactant edge
\colorlet{scProd}{npgGreen!82!black}  % product edge
\colorlet{scHiF}{npgOrange!13}        % highlight fill
\colorlet{scHiE}{npgOrange!82!black}  % highlight edge
% Stage / accent tints.  stDE remains available for the fifth plot series.
\colorlet{stAF}{npgGreen!9}\colorlet{stAE}{npgGreen}
\colorlet{stBF}{npgOrange!11}\colorlet{stBE}{npgOrange}
\colorlet{stCF}{npgBlue!9}\colorlet{stCE}{npgBlue}
\colorlet{stDF}{npgPurple!10}\colorlet{stDE}{npgPurple!85!black}
\colorlet{stEF}{npgVermillion!9}\colorlet{stEE}{npgVermillion}

% Semantic aliases.  Figures should speak in these, not in stAE/stBE/...
\colorlet{scAccent}{stAE}\colorlet{scAccentF}{stAF}   % structure / normal
\colorlet{scMark}{stBE}\colorlet{scMarkF}{stBF}       % highlight / certificate
\colorlet{scRef}{stCE}\colorlet{scRefF}{stCF}         % secondary series
\colorlet{scAlert}{stEE}\colorlet{scAlertF}{stEF}     % failure / violation
\colorlet{scRule}{figHair}                            % card and axis hairlines
\colorlet{scRuleL}{figHair!55}                        % grid, dividers

% --- text styles ----------------------------------------------------------
% Three sizes, fixed in points so they do not drift with the document class.
% Arial is shared by labels, headings, panel badges, and chemistry labels.
\providecommand{\scTitle}{\scArial\fontsize{7}{8.4}\selectfont\bfseries\color{scInk}}
\providecommand{\scBody}{\scArial\fontsize{6.5}{7.8}\selectfont\color{scInkM}}
\providecommand{\scSmall}{\scArial\fontsize{6}{7.2}\selectfont\color{scInkM}}
\providecommand{\scSub}[1]{{\scSmall\itshape #1}}
\providecommand{\scPanel}[1]{{\scTitle(#1)}}
% Numerals that carry a result: ink, not colour, so the eye reads the value.
\providecommand{\scNum}{\scArial\fontsize{10}{11}\selectfont\bfseries\color{scInk}}

% --- shared TikZ styles ---------------------------------------------------
\tikzset{
  % text never hyphenates inside a figure
  every node/.append style={execute at begin node={%
      \hyphenpenalty=10000 \exhyphenpenalty=10000 \hbadness=10000 }},
  syn/.style={font=\scArial,color=scInk,
              execute at begin picture={\MTversion*{synarial}},
              >={Stealth[round,length=2.4mm,width=1.9mm]}},
  % containers ------------------------------------------------------------
  % One card chrome for the whole paper: white, 3 pt radius, one hairline.
  panel/.style={rounded corners=2.5pt,fill=none,draw=scRule,line width=0.45pt},
  card/.style={rounded corners=2.2pt,fill=scPaper,draw=scRule,line width=0.45pt,
               inner sep=5pt,align=left,anchor=north west},
  % a quiet strip for a panel's closing fact
  note/.style={rounded corners=2.5pt,fill=scPaper,draw=scRuleL,line width=0.4pt,
               inner xsep=6pt,inner ysep=3.5pt,align=left,font=\scSmall},
  % small labelled value pill
  chip/.style={rounded corners=2pt,fill=scPaper,draw=scRule,line width=0.4pt,
               inner xsep=4pt,inner ysep=2pt,font=\scSmall},
  % Circular letter badge: the same accessible panel marker used in Style/*.
  hbullet/.style={circle,fill=scInk,text=white,draw=none,
                  font=\scArial\fontsize{6}{6}\selectfont\bfseries,
                  minimum size=4.2mm,inner sep=0pt},
  ctag/.style={font=\scTitle,anchor=north west,text=scInk},
  % CRN primitives --------------------------------------------------------
  species/.style={rounded rectangle,fill=scSpF,draw=scSpE,line width=0.45pt,
                  inner xsep=2.4pt,inner ysep=1.6pt,font=\scSmall,
                  text=scInk,minimum height=3.6mm},
  cofactor/.style={species,fill=scRefF,draw=scRef},
  catalyst/.style={species,fill=scHiF,draw=scHiE},
  rxn/.style={rectangle,fill=scRxF,draw=scRxE,line width=0.45pt,
              inner sep=1.2pt,font=\scSmall,text=scInk,
              minimum width=3.0mm,minimum height=3.0mm},
  place/.style={circle,fill=scSpF,draw=scSpE,line width=0.45pt,
                minimum size=4.0mm,inner sep=0pt},
  trans/.style={rectangle,fill=scRxE,draw=scRxE,minimum width=1.1mm,
                minimum height=3.6mm,inner sep=0pt},
  token/.style={circle,fill=scInk,minimum size=1.1mm,inner sep=0pt},
  % edges -----------------------------------------------------------------
  reac/.style={->,line width=0.5pt,draw=scReac,shorten >=0.8pt,shorten <=0.8pt},
  prod/.style={->,line width=0.5pt,draw=scProd,shorten >=0.8pt,shorten <=0.8pt},
  cat/.style={->,line width=0.5pt,draw=scHiE,shorten >=0.8pt,shorten <=0.8pt},
  flow/.style={-{Stealth[round,length=2.2mm,width=1.8mm]},line width=1.1pt,
               draw=scInkM!45},
  flab/.style={font=\scSmall,text=scInkM,inner sep=2pt},
  stoi/.style={font=\scSmall,text=scInkM,inner sep=0.8pt},
  % verdict marks: shape carries the meaning, colour only reinforces it
  vpass/.style={circle,fill=scAccentF,draw=scAccent,text=scAccent,
                font=\fontsize{6}{6}\selectfont\bfseries,
                minimum size=3.2mm,inner sep=0pt},
  vfail/.style={circle,fill=scAlertF,draw=scAlert,text=scAlert,
                font=\fontsize{6}{6}\selectfont\bfseries,
                minimum size=3.2mm,inner sep=0pt},
  callout/.style={font=\scSmall,text=scInkM,align=left},
}

% \scHead[<badge colour>]{<anchor coord>}{<letter>}{<title>}{<subtitle or empty>}
% One header treatment for every panel in the paper.
\providecommand{\scHead}[5][scInk]{%
  \node[hbullet,fill=#1,anchor=north west] (scbadge#3) at ($#2+(0,-0.015)$)
    {\MakeUppercase{#3}};
  \node[ctag,anchor=west] at ($(scbadge#3.east)+(0.14,0)$) {#4};
  \ifx\relax#5\relax\else
    \node[anchor=north west,font=\scSmall,text=scInkM]
      at ($#2+(0.56,-0.245)$) {#5};
  \fi}

% --- shared pgfplots look (only if the figure loaded pgfplots) ------------
\ifdefined\pgfplotsset
  \pgfplotsset{compat=1.18}
  \pgfplotsset{
    synaxis/.style={
      scale only axis,
      tick align=outside, tick pos=left,
      axis line style={draw=scInkM!55,line width=0.4pt},
      % One grid, on the value axis only.  No minor grid: at these sizes it
      % turns the plot frame into a grey wash.
      ymajorgrids=true, xmajorgrids=false,
      major grid style={draw=scRuleL,line width=0.3pt},
      major tick length=1.6pt,
      label style={font=\scBody,color=scInk},
      tick label style={font=\scSmall,color=scInkM},
      legend style={font=\scSmall,draw=none,fill=none,inner sep=1pt,
                    row sep=0.4pt,legend columns=1},
      legend cell align=left,
      every axis plot/.append style={line width=0.7pt,mark size=1.2pt},
    },
  }
\fi
}
% so that the figure compiles both from Figure/ and from the paper root.
%
% The visual vocabulary follows ../../../../Style/style.tex: the
% colour-blind-safe Okabe--Ito accents, quiet neutral rules and wash, circular
% panel badges, and a compact technical type scale are shared by the companion
% papers.  Semantic aliases below keep the CRN meaning explicit.
%
% DESIGN CONTRACT -----------------------------------------------------------
% 1. Geometry.  Every figure is laid out on a canvas exactly \scW wide (the
%    sn-jnl text width, 372 pt) and is included unscaled, so one point of type
%    is the same physical size in all five figures.  Open each picture with
%    \scCanvas{<height in cm>} to pin the bounding box.
% 2. Colour is semantic, never decorative.  Four roles only:
%       scAccent (green)       structure, the normal/expected path
%       scMark   (orange)      catalyst, certificate, or highlight
%       scAlert  (vermillion)  failure, violation, or uncovered set
%       scRef    (blue)        secondary or reference series
%    A panel gets at most one accent.  Titles are always ink; panel identity
%    is also encoded by a letter inside a circular badge, not by colour alone.
% 3. Type scale.  Three sizes, no others, and nothing below 6 pt:
%       \scTitle (7 pt bold)  panel titles
%       \scBody  (6.5 pt)     labels, values, in-figure prose
%       \scSmall (6 pt)       units, footnotes, dense graph labels
% ---------------------------------------------------------------------------
% Load guard.  The paper inputs every figure source inline, so this file is
% reached once per figure; without the guard the second \newlength below
% would abort the run.
\ifdefined\SynKitStyleLoaded\endinput\fi
\newcommand{\SynKitStyleLoaded}{}

% Load in the document preamble. XeLaTeX/LuaLaTeX and an installed Arial
% family are required; fontspec reports an error instead of substituting.
\RequirePackage[no-math]{fontspec}
\newfontfamily\scArial{Arial}[NFSSFamily=scArial,Ligatures=TeX]
% Keep manuscript mathematics unchanged; select Arial letters and numerals
% only inside figures. TeX supplies specialist mathematical symbols.
\RequirePackage[subdued,italic]{mathastext}
\MTDeclareVersion{synarial}{TU}{scArial}{m}{n}
\SetMathAlphabet{\Mathsf}{synarial}{TU}{scArial}{m}{n}
\SetMathAlphabet{\Mathtt}{synarial}{TU}{scArial}{m}{n}

\usetikzlibrary{positioning,calc,fit,backgrounds,arrows.meta,shapes.geometric,
                shapes.misc,decorations.pathreplacing,patterns,math}

% --- canvas ---------------------------------------------------------------
% 372 pt text width minus the 1 pt standalone border on each side.
\newlength{\scW}\setlength{\scW}{13.0326cm}
\providecommand{\scWnum}{13.0326}
% \scCanvas{<height>} pins the picture to the full text width, y running down
% from 0, so every figure shares one coordinate convention.  It *sets* the
% bounding box rather than contributing to it, so the emitted page is exactly
% the text width and every figure is included at 1:1 -- otherwise a stray
% invisible extent (pgfplots reserves a little around an axis) rescales one
% figure relative to the others and the same \scSmall prints at two sizes.
% Anything drawn outside the canvas is therefore cropped: keep content inside,
% and check with
%     \providecommand{\scCanvas}[1]{\path[draw=red] (0,0) rectangle (\scWnum,-#1);}
\providecommand{\scCanvas}[1]{%
  \path[use as bounding box] (0,0) rectangle (\scWnum,-#1);}

% --- palette --------------------------------------------------------------
% Canonical companion-paper values (../../../../Style/style.tex).
\definecolor{npgOrange}{HTML}{E69F00}
\definecolor{npgSky}{HTML}{56B4E9}
\definecolor{npgGreen}{HTML}{009E73}
\definecolor{npgYellow}{HTML}{F0E442}
\definecolor{npgBlue}{HTML}{0072B2}
\definecolor{npgVermillion}{HTML}{E64B35}
\definecolor{npgPurple}{HTML}{CC79A7}
\definecolor{figInk}{HTML}{1A1A1A}
\definecolor{figMuted}{HTML}{666666}
\definecolor{figHair}{HTML}{D9D9D9}
\definecolor{figWash}{HTML}{F7F7F7}

\colorlet{scInk}{figInk}              % primary text
\colorlet{scInkM}{figMuted}           % secondary text
\colorlet{scPaper}{figWash}           % quiet cards and note strips
\colorlet{scSpF}{npgSky!12}           % species fill
\colorlet{scSpE}{npgBlue}             % species edge
\colorlet{scRxF}{npgOrange!12}        % reaction fill
\colorlet{scRxE}{npgOrange!88!black}  % reaction edge
\colorlet{scReac}{npgBlue!78!black}   % reactant edge
\colorlet{scProd}{npgGreen!82!black}  % product edge
\colorlet{scHiF}{npgOrange!13}        % highlight fill
\colorlet{scHiE}{npgOrange!82!black}  % highlight edge
% Stage / accent tints.  stDE remains available for the fifth plot series.
\colorlet{stAF}{npgGreen!9}\colorlet{stAE}{npgGreen}
\colorlet{stBF}{npgOrange!11}\colorlet{stBE}{npgOrange}
\colorlet{stCF}{npgBlue!9}\colorlet{stCE}{npgBlue}
\colorlet{stDF}{npgPurple!10}\colorlet{stDE}{npgPurple!85!black}
\colorlet{stEF}{npgVermillion!9}\colorlet{stEE}{npgVermillion}

% Semantic aliases.  Figures should speak in these, not in stAE/stBE/...
\colorlet{scAccent}{stAE}\colorlet{scAccentF}{stAF}   % structure / normal
\colorlet{scMark}{stBE}\colorlet{scMarkF}{stBF}       % highlight / certificate
\colorlet{scRef}{stCE}\colorlet{scRefF}{stCF}         % secondary series
\colorlet{scAlert}{stEE}\colorlet{scAlertF}{stEF}     % failure / violation
\colorlet{scRule}{figHair}                            % card and axis hairlines
\colorlet{scRuleL}{figHair!55}                        % grid, dividers

% --- text styles ----------------------------------------------------------
% Three sizes, fixed in points so they do not drift with the document class.
% Arial is shared by labels, headings, panel badges, and chemistry labels.
\providecommand{\scTitle}{\scArial\fontsize{7}{8.4}\selectfont\bfseries\color{scInk}}
\providecommand{\scBody}{\scArial\fontsize{6.5}{7.8}\selectfont\color{scInkM}}
\providecommand{\scSmall}{\scArial\fontsize{6}{7.2}\selectfont\color{scInkM}}
\providecommand{\scSub}[1]{{\scSmall\itshape #1}}
\providecommand{\scPanel}[1]{{\scTitle(#1)}}
% Numerals that carry a result: ink, not colour, so the eye reads the value.
\providecommand{\scNum}{\scArial\fontsize{10}{11}\selectfont\bfseries\color{scInk}}

% --- shared TikZ styles ---------------------------------------------------
\tikzset{
  % text never hyphenates inside a figure
  every node/.append style={execute at begin node={%
      \hyphenpenalty=10000 \exhyphenpenalty=10000 \hbadness=10000 }},
  syn/.style={font=\scArial,color=scInk,
              execute at begin picture={\MTversion*{synarial}},
              >={Stealth[round,length=2.4mm,width=1.9mm]}},
  % containers ------------------------------------------------------------
  % One card chrome for the whole paper: white, 3 pt radius, one hairline.
  panel/.style={rounded corners=2.5pt,fill=none,draw=scRule,line width=0.45pt},
  card/.style={rounded corners=2.2pt,fill=scPaper,draw=scRule,line width=0.45pt,
               inner sep=5pt,align=left,anchor=north west},
  % a quiet strip for a panel's closing fact
  note/.style={rounded corners=2.5pt,fill=scPaper,draw=scRuleL,line width=0.4pt,
               inner xsep=6pt,inner ysep=3.5pt,align=left,font=\scSmall},
  % small labelled value pill
  chip/.style={rounded corners=2pt,fill=scPaper,draw=scRule,line width=0.4pt,
               inner xsep=4pt,inner ysep=2pt,font=\scSmall},
  % Circular letter badge: the same accessible panel marker used in Style/*.
  hbullet/.style={circle,fill=scInk,text=white,draw=none,
                  font=\scArial\fontsize{6}{6}\selectfont\bfseries,
                  minimum size=4.2mm,inner sep=0pt},
  ctag/.style={font=\scTitle,anchor=north west,text=scInk},
  % CRN primitives --------------------------------------------------------
  species/.style={rounded rectangle,fill=scSpF,draw=scSpE,line width=0.45pt,
                  inner xsep=2.4pt,inner ysep=1.6pt,font=\scSmall,
                  text=scInk,minimum height=3.6mm},
  cofactor/.style={species,fill=scRefF,draw=scRef},
  catalyst/.style={species,fill=scHiF,draw=scHiE},
  rxn/.style={rectangle,fill=scRxF,draw=scRxE,line width=0.45pt,
              inner sep=1.2pt,font=\scSmall,text=scInk,
              minimum width=3.0mm,minimum height=3.0mm},
  place/.style={circle,fill=scSpF,draw=scSpE,line width=0.45pt,
                minimum size=4.0mm,inner sep=0pt},
  trans/.style={rectangle,fill=scRxE,draw=scRxE,minimum width=1.1mm,
                minimum height=3.6mm,inner sep=0pt},
  token/.style={circle,fill=scInk,minimum size=1.1mm,inner sep=0pt},
  % edges -----------------------------------------------------------------
  reac/.style={->,line width=0.5pt,draw=scReac,shorten >=0.8pt,shorten <=0.8pt},
  prod/.style={->,line width=0.5pt,draw=scProd,shorten >=0.8pt,shorten <=0.8pt},
  cat/.style={->,line width=0.5pt,draw=scHiE,shorten >=0.8pt,shorten <=0.8pt},
  flow/.style={-{Stealth[round,length=2.2mm,width=1.8mm]},line width=1.1pt,
               draw=scInkM!45},
  flab/.style={font=\scSmall,text=scInkM,inner sep=2pt},
  stoi/.style={font=\scSmall,text=scInkM,inner sep=0.8pt},
  % verdict marks: shape carries the meaning, colour only reinforces it
  vpass/.style={circle,fill=scAccentF,draw=scAccent,text=scAccent,
                font=\fontsize{6}{6}\selectfont\bfseries,
                minimum size=3.2mm,inner sep=0pt},
  vfail/.style={circle,fill=scAlertF,draw=scAlert,text=scAlert,
                font=\fontsize{6}{6}\selectfont\bfseries,
                minimum size=3.2mm,inner sep=0pt},
  callout/.style={font=\scSmall,text=scInkM,align=left},
}

% \scHead[<badge colour>]{<anchor coord>}{<letter>}{<title>}{<subtitle or empty>}
% One header treatment for every panel in the paper.
\providecommand{\scHead}[5][scInk]{%
  \node[hbullet,fill=#1,anchor=north west] (scbadge#3) at ($#2+(0,-0.015)$)
    {\MakeUppercase{#3}};
  \node[ctag,anchor=west] at ($(scbadge#3.east)+(0.14,0)$) {#4};
  \ifx\relax#5\relax\else
    \node[anchor=north west,font=\scSmall,text=scInkM]
      at ($#2+(0.56,-0.245)$) {#5};
  \fi}

% --- shared pgfplots look (only if the figure loaded pgfplots) ------------
\ifdefined\pgfplotsset
  \pgfplotsset{compat=1.18}
  \pgfplotsset{
    synaxis/.style={
      scale only axis,
      tick align=outside, tick pos=left,
      axis line style={draw=scInkM!55,line width=0.4pt},
      % One grid, on the value axis only.  No minor grid: at these sizes it
      % turns the plot frame into a grey wash.
      ymajorgrids=true, xmajorgrids=false,
      major grid style={draw=scRuleL,line width=0.3pt},
      major tick length=1.6pt,
      label style={font=\scBody,color=scInk},
      tick label style={font=\scSmall,color=scInkM},
      legend style={font=\scSmall,draw=none,fill=none,inner sep=1pt,
                    row sep=0.4pt,legend columns=1},
      legend cell align=left,
      every axis plot/.append style={line width=0.7pt,mark size=1.2pt},
    },
  }
\fi
}
% so that the figure compiles both from Figure/ and from the paper root.
%
% The visual vocabulary follows ../../../../Style/style.tex: the
% colour-blind-safe Okabe--Ito accents, quiet neutral rules and wash, circular
% panel badges, and a compact technical type scale are shared by the companion
% papers.  Semantic aliases below keep the CRN meaning explicit.
%
% DESIGN CONTRACT -----------------------------------------------------------
% 1. Geometry.  Every figure is laid out on a canvas exactly \scW wide (the
%    sn-jnl text width, 372 pt) and is included unscaled, so one point of type
%    is the same physical size in all five figures.  Open each picture with
%    \scCanvas{<height in cm>} to pin the bounding box.
% 2. Colour is semantic, never decorative.  Four roles only:
%       scAccent (green)       structure, the normal/expected path
%       scMark   (orange)      catalyst, certificate, or highlight
%       scAlert  (vermillion)  failure, violation, or uncovered set
%       scRef    (blue)        secondary or reference series
%    A panel gets at most one accent.  Titles are always ink; panel identity
%    is also encoded by a letter inside a circular badge, not by colour alone.
% 3. Type scale.  Three sizes, no others, and nothing below 6 pt:
%       \scTitle (7 pt bold)  panel titles
%       \scBody  (6.5 pt)     labels, values, in-figure prose
%       \scSmall (6 pt)       units, footnotes, dense graph labels
% ---------------------------------------------------------------------------
% Load guard.  The paper inputs every figure source inline, so this file is
% reached once per figure; without the guard the second \newlength below
% would abort the run.
\ifdefined\SynKitStyleLoaded\endinput\fi
\newcommand{\SynKitStyleLoaded}{}

% Load in the document preamble. XeLaTeX/LuaLaTeX and an installed Arial
% family are required; fontspec reports an error instead of substituting.
\RequirePackage[no-math]{fontspec}
\newfontfamily\scArial{Arial}[NFSSFamily=scArial,Ligatures=TeX]
% Keep manuscript mathematics unchanged; select Arial letters and numerals
% only inside figures. TeX supplies specialist mathematical symbols.
\RequirePackage[subdued,italic]{mathastext}
\MTDeclareVersion{synarial}{TU}{scArial}{m}{n}
\SetMathAlphabet{\Mathsf}{synarial}{TU}{scArial}{m}{n}
\SetMathAlphabet{\Mathtt}{synarial}{TU}{scArial}{m}{n}

\usetikzlibrary{positioning,calc,fit,backgrounds,arrows.meta,shapes.geometric,
                shapes.misc,decorations.pathreplacing,patterns,math}

% --- canvas ---------------------------------------------------------------
% 372 pt text width minus the 1 pt standalone border on each side.
\newlength{\scW}\setlength{\scW}{13.0326cm}
\providecommand{\scWnum}{13.0326}
% \scCanvas{<height>} pins the picture to the full text width, y running down
% from 0, so every figure shares one coordinate convention.  It *sets* the
% bounding box rather than contributing to it, so the emitted page is exactly
% the text width and every figure is included at 1:1 -- otherwise a stray
% invisible extent (pgfplots reserves a little around an axis) rescales one
% figure relative to the others and the same \scSmall prints at two sizes.
% Anything drawn outside the canvas is therefore cropped: keep content inside,
% and check with
%     \providecommand{\scCanvas}[1]{\path[draw=red] (0,0) rectangle (\scWnum,-#1);}
\providecommand{\scCanvas}[1]{%
  \path[use as bounding box] (0,0) rectangle (\scWnum,-#1);}

% --- palette --------------------------------------------------------------
% Canonical companion-paper values (../../../../Style/style.tex).
\definecolor{npgOrange}{HTML}{E69F00}
\definecolor{npgSky}{HTML}{56B4E9}
\definecolor{npgGreen}{HTML}{009E73}
\definecolor{npgYellow}{HTML}{F0E442}
\definecolor{npgBlue}{HTML}{0072B2}
\definecolor{npgVermillion}{HTML}{E64B35}
\definecolor{npgPurple}{HTML}{CC79A7}
\definecolor{figInk}{HTML}{1A1A1A}
\definecolor{figMuted}{HTML}{666666}
\definecolor{figHair}{HTML}{D9D9D9}
\definecolor{figWash}{HTML}{F7F7F7}

\colorlet{scInk}{figInk}              % primary text
\colorlet{scInkM}{figMuted}           % secondary text
\colorlet{scPaper}{figWash}           % quiet cards and note strips
\colorlet{scSpF}{npgSky!12}           % species fill
\colorlet{scSpE}{npgBlue}             % species edge
\colorlet{scRxF}{npgOrange!12}        % reaction fill
\colorlet{scRxE}{npgOrange!88!black}  % reaction edge
\colorlet{scReac}{npgBlue!78!black}   % reactant edge
\colorlet{scProd}{npgGreen!82!black}  % product edge
\colorlet{scHiF}{npgOrange!13}        % highlight fill
\colorlet{scHiE}{npgOrange!82!black}  % highlight edge
% Stage / accent tints.  stDE remains available for the fifth plot series.
\colorlet{stAF}{npgGreen!9}\colorlet{stAE}{npgGreen}
\colorlet{stBF}{npgOrange!11}\colorlet{stBE}{npgOrange}
\colorlet{stCF}{npgBlue!9}\colorlet{stCE}{npgBlue}
\colorlet{stDF}{npgPurple!10}\colorlet{stDE}{npgPurple!85!black}
\colorlet{stEF}{npgVermillion!9}\colorlet{stEE}{npgVermillion}

% Semantic aliases.  Figures should speak in these, not in stAE/stBE/...
\colorlet{scAccent}{stAE}\colorlet{scAccentF}{stAF}   % structure / normal
\colorlet{scMark}{stBE}\colorlet{scMarkF}{stBF}       % highlight / certificate
\colorlet{scRef}{stCE}\colorlet{scRefF}{stCF}         % secondary series
\colorlet{scAlert}{stEE}\colorlet{scAlertF}{stEF}     % failure / violation
\colorlet{scRule}{figHair}                            % card and axis hairlines
\colorlet{scRuleL}{figHair!55}                        % grid, dividers

% --- text styles ----------------------------------------------------------
% Three sizes, fixed in points so they do not drift with the document class.
% Arial is shared by labels, headings, panel badges, and chemistry labels.
\providecommand{\scTitle}{\scArial\fontsize{7}{8.4}\selectfont\bfseries\color{scInk}}
\providecommand{\scBody}{\scArial\fontsize{6.5}{7.8}\selectfont\color{scInkM}}
\providecommand{\scSmall}{\scArial\fontsize{6}{7.2}\selectfont\color{scInkM}}
\providecommand{\scSub}[1]{{\scSmall\itshape #1}}
\providecommand{\scPanel}[1]{{\scTitle(#1)}}
% Numerals that carry a result: ink, not colour, so the eye reads the value.
\providecommand{\scNum}{\scArial\fontsize{10}{11}\selectfont\bfseries\color{scInk}}

% --- shared TikZ styles ---------------------------------------------------
\tikzset{
  % text never hyphenates inside a figure
  every node/.append style={execute at begin node={%
      \hyphenpenalty=10000 \exhyphenpenalty=10000 \hbadness=10000 }},
  syn/.style={font=\scArial,color=scInk,
              execute at begin picture={\MTversion*{synarial}},
              >={Stealth[round,length=2.4mm,width=1.9mm]}},
  % containers ------------------------------------------------------------
  % One card chrome for the whole paper: white, 3 pt radius, one hairline.
  panel/.style={rounded corners=2.5pt,fill=none,draw=scRule,line width=0.45pt},
  card/.style={rounded corners=2.2pt,fill=scPaper,draw=scRule,line width=0.45pt,
               inner sep=5pt,align=left,anchor=north west},
  % a quiet strip for a panel's closing fact
  note/.style={rounded corners=2.5pt,fill=scPaper,draw=scRuleL,line width=0.4pt,
               inner xsep=6pt,inner ysep=3.5pt,align=left,font=\scSmall},
  % small labelled value pill
  chip/.style={rounded corners=2pt,fill=scPaper,draw=scRule,line width=0.4pt,
               inner xsep=4pt,inner ysep=2pt,font=\scSmall},
  % Circular letter badge: the same accessible panel marker used in Style/*.
  hbullet/.style={circle,fill=scInk,text=white,draw=none,
                  font=\scArial\fontsize{6}{6}\selectfont\bfseries,
                  minimum size=4.2mm,inner sep=0pt},
  ctag/.style={font=\scTitle,anchor=north west,text=scInk},
  % CRN primitives --------------------------------------------------------
  species/.style={rounded rectangle,fill=scSpF,draw=scSpE,line width=0.45pt,
                  inner xsep=2.4pt,inner ysep=1.6pt,font=\scSmall,
                  text=scInk,minimum height=3.6mm},
  cofactor/.style={species,fill=scRefF,draw=scRef},
  catalyst/.style={species,fill=scHiF,draw=scHiE},
  rxn/.style={rectangle,fill=scRxF,draw=scRxE,line width=0.45pt,
              inner sep=1.2pt,font=\scSmall,text=scInk,
              minimum width=3.0mm,minimum height=3.0mm},
  place/.style={circle,fill=scSpF,draw=scSpE,line width=0.45pt,
                minimum size=4.0mm,inner sep=0pt},
  trans/.style={rectangle,fill=scRxE,draw=scRxE,minimum width=1.1mm,
                minimum height=3.6mm,inner sep=0pt},
  token/.style={circle,fill=scInk,minimum size=1.1mm,inner sep=0pt},
  % edges -----------------------------------------------------------------
  reac/.style={->,line width=0.5pt,draw=scReac,shorten >=0.8pt,shorten <=0.8pt},
  prod/.style={->,line width=0.5pt,draw=scProd,shorten >=0.8pt,shorten <=0.8pt},
  cat/.style={->,line width=0.5pt,draw=scHiE,shorten >=0.8pt,shorten <=0.8pt},
  flow/.style={-{Stealth[round,length=2.2mm,width=1.8mm]},line width=1.1pt,
               draw=scInkM!45},
  flab/.style={font=\scSmall,text=scInkM,inner sep=2pt},
  stoi/.style={font=\scSmall,text=scInkM,inner sep=0.8pt},
  % verdict marks: shape carries the meaning, colour only reinforces it
  vpass/.style={circle,fill=scAccentF,draw=scAccent,text=scAccent,
                font=\fontsize{6}{6}\selectfont\bfseries,
                minimum size=3.2mm,inner sep=0pt},
  vfail/.style={circle,fill=scAlertF,draw=scAlert,text=scAlert,
                font=\fontsize{6}{6}\selectfont\bfseries,
                minimum size=3.2mm,inner sep=0pt},
  callout/.style={font=\scSmall,text=scInkM,align=left},
}

% \scHead[<badge colour>]{<anchor coord>}{<letter>}{<title>}{<subtitle or empty>}
% One header treatment for every panel in the paper.
\providecommand{\scHead}[5][scInk]{%
  \node[hbullet,fill=#1,anchor=north west] (scbadge#3) at ($#2+(0,-0.015)$)
    {\MakeUppercase{#3}};
  \node[ctag,anchor=west] at ($(scbadge#3.east)+(0.14,0)$) {#4};
  \ifx\relax#5\relax\else
    \node[anchor=north west,font=\scSmall,text=scInkM]
      at ($#2+(0.56,-0.245)$) {#5};
  \fi}

% --- shared pgfplots look (only if the figure loaded pgfplots) ------------
\ifdefined\pgfplotsset
  \pgfplotsset{compat=1.18}
  \pgfplotsset{
    synaxis/.style={
      scale only axis,
      tick align=outside, tick pos=left,
      axis line style={draw=scInkM!55,line width=0.4pt},
      % One grid, on the value axis only.  No minor grid: at these sizes it
      % turns the plot frame into a grey wash.
      ymajorgrids=true, xmajorgrids=false,
      major grid style={draw=scRuleL,line width=0.3pt},
      major tick length=1.6pt,
      label style={font=\scBody,color=scInk},
      tick label style={font=\scSmall,color=scInkM},
      legend style={font=\scSmall,draw=none,fill=none,inner sep=1pt,
                    row sep=0.4pt,legend columns=1},
      legend cell align=left,
      every axis plot/.append style={line width=0.7pt,mark size=1.2pt},
    },
  }
\fi
}}

\tikzset{
  minisp/.style={circle,fill=scSpF,draw=scSpE,line width=0.45pt,
    minimum size=3.2mm,inner sep=0pt},
  minirx/.style={rectangle,fill=scRxF,draw=scRxE,line width=0.45pt,
    minimum size=2.7mm,inner sep=0pt},
  % the route arrows between panels
  pipe/.style={-{Stealth[round,length=1.9mm,width=1.6mm]},draw=scInkM!38,
    line width=1.0pt},
  % one row of an adapter list: icon column, then name, then format
  rowname/.style={font=\scBody,text=scInk,anchor=base west,
    execute at begin node={\bfseries}},
  rowgloss/.style={font=\scSmall,text=scInkM,anchor=base west},
  % a row whose name is optically centred on its icon (no second line)
  rowmid/.style={font=\scBody,text=scInk,anchor=west,
    execute at begin node={\bfseries}},
}

% Column geometry, in cm, on the 16.40 canvas.
\def\W{16.40}
\def\ax{0}\def\aw{4.05}
\def\bx{4.55}\def\bw{7.30}
\def\cx{12.35}\def\cw{4.05}
\def\rowh{3.40}          % height of the top row of panels
\def\dy{-3.82}\def\dh{2.55}  % analysis band
% Icon column: the widest glyph reaches 0.62 cm from its centre, so names start
% 0.78 cm out.  At the previous 0.62/1.06 the hypergraph ran 1.8 mm into its name.
\def\icon{0.80}\def\name{1.58}
% Row centres, shared by panels (a) and (c) so the two columns align.
\def\rya{-1.18}\def\ryb{-2.05}\def\ryc{-2.92}

\begin{tikzpicture}[syn]
\path[use as bounding box] (0,0) rectangle (\W,-6.37);

% =========================================================================
% (a) input adapters
% =========================================================================
\path[panel] (\ax,0) rectangle (\ax+\aw,-\rowh);
\scHead{(\ax+0.22,-0.16)}{a}{input adapters}{three routes in}

% --- Query: a database cylinder
\begin{scope}[shift={(\ax+\icon,\rya+0.16)},draw=scAccent,line width=0.5pt]
  \draw[fill=scAccentF] (0,0) ellipse (0.17 and 0.065);
  \draw[fill=scAccentF] (-0.17,0) -- (-0.17,-0.26)
    arc[start angle=180,end angle=360,x radius=0.17,y radius=0.065] -- (0.17,0);
  \draw (-0.17,-0.13) arc[start angle=180,end angle=360,x radius=0.17,y radius=0.065];
\end{scope}
\node[rowname]  at (\ax+\name,\rya+0.04) {Query};
\node[rowgloss] at (\ax+\name,\rya-0.22) {KEGG modules};

% --- IO: an XML tag
\node[rounded corners=1pt,fill=scRefF,draw=scRef,line width=0.4pt,
      font=\scSmall,text=scRef,inner xsep=2.5pt,inner ysep=1.5pt]
  at (\ax+\icon,\ryb) {<xml>};
\node[rowname]  at (\ax+\name,\ryb+0.04) {IO};
\node[rowgloss] at (\ax+\name,\ryb-0.22) {SBML L3 Core};

% --- Construct: a rule firing
\begin{scope}[shift={(\ax+\icon,\ryc)}]
  \node[minisp] (ia) at (-0.42,0) {};
  \node[minirx,fill=scMarkF,draw=scMark] (ir) at (0,0) {};
  \node[minisp] (ib) at (0.42,0) {};
  \draw[reac] (ia)--(ir); \draw[prod] (ir)--(ib);
\end{scope}
\node[rowname]  at (\ax+\name,\ryc+0.04) {Construct};
\node[rowgloss] at (\ax+\name,\ryc-0.22) {SMARTS / DPO rules};

% =========================================================================
% (b) the canonical object
% =========================================================================
\draw[pipe] (\ax+\aw+0.10,-1.70) -- (\bx-0.10,-1.70);

\path[panel] (\bx,0) rectangle (\bx+\bw,-\rowh);
\scHead[scMark]{(\bx+0.22,-0.16)}{b}{reaction records}{}
\node[anchor=north east,font=\scBody,text=scInk] at (\bx+\bw-0.22,-0.155)
  {\textbf{SynCRN}};

% One reaction with two substrates, one with a catalyst that is consumed and
% returned.  Both are exactly what a net-stoichiometry view would lose.
\begin{scope}[shift={(\bx+0.95,-1.85)}]
  \node[minisp] (a)  at (0,0.58) {};
  \node[minisp] (b)  at (0,-0.58) {};
  \node[minirx] (r1) at (1.35,0) {};
  \node[minisp] (c)  at (2.70,0) {};
  \node[minisp] (d)  at (2.70,-0.92) {};
  \node[minirx] (r2) at (4.05,-0.38) {};
  \node[minisp] (e)  at (5.40,0.14) {};
  \node[minisp] (f)  at (5.40,-0.88) {};
  \node[minisp,fill=scHiF,draw=scHiE] (cat) at (4.05,0.78) {};
  \draw[reac] (a)--(r1); \draw[reac] (b)--(r1);
  \draw[prod] (r1)--(c);
  \draw[reac] (c)--(r2); \draw[reac] (d)--(r2);
  \draw[prod] (r2)--(e); \draw[prod] (r2)--(f);
  % the catalyst is consumed and returned by the same reaction
  \draw[cat] (cat) to[out=200,in=110] (r2);
  \draw[cat] (r2) to[out=70,in=-20] (cat);
  \node[stoi,fill=white,inner sep=1pt] at (0.46,0.33) {2};
  \node[font=\scSmall,text=scHiE,anchor=west] at (4.35,0.78) {catalyst};
\end{scope}

\node[anchor=north,font=\scBody,text=scInkM] at (\bx+0.5*\bw,-3.02)
  {$\mathcal N=(X,\mathscr R,\mathbf S^{-},\mathbf S^{+})$};

% =========================================================================
% (c) derived views
% =========================================================================
\draw[pipe] (\bx+\bw+0.10,-1.70) -- (\cx-0.10,-1.70);

\path[panel] (\cx,0) rectangle (\cx+\cw,-\rowh);
\scHead[scRef]{(\cx+0.22,-0.16)}{c}{derived views}{on demand}

% --- stoichiometry: a small sign pattern
\begin{scope}[shift={(\cx+\icon,\rya)}]
  \foreach \x/\y/\c in {-0.19/0.19/scAlertF,0/0.19/scAccentF,0.19/0.19/white,
                        -0.19/0/white,0/0/scRefF,0.19/0/scAlertF,
                        -0.19/-0.19/scAccentF,0/-0.19/white,0.19/-0.19/scAccentF}
    \node[draw=scRule,fill=\c,minimum size=1.8mm,inner sep=0pt] at (\x,\y) {};
\end{scope}
\node[rowmid] at (\cx+\name,\rya) {stoichiometry};

% --- Petri net: a marked place firing
\begin{scope}[shift={(\cx+\icon,\ryb)}]
  \node[place,minimum size=3.0mm] (p1) at (-0.42,0) {};
  \node[token] at (-0.42,0) {};
  \node[trans,minimum height=3.0mm] (pt) at (0,0) {};
  \node[place,minimum size=3.0mm] (p2) at (0.42,0) {};
  \draw[reac] (p1)--(pt); \draw[prod] (pt)--(p2);
\end{scope}
\node[rowmid] at (\cx+\name,\ryb) {Petri net};

% --- hypergraph: a two-tail edge
\begin{scope}[shift={(\cx+\icon,\ryc)}]
  \node[minisp] (h1) at (-0.46,0.20) {};
  \node[minisp] (h2) at (-0.46,-0.20) {};
  \node[rounded corners=1pt,draw=scRxE,fill=scRxF,minimum width=1.8mm,
        minimum height=5.0mm,inner sep=0pt] (hr) at (0,0) {};
  \node[minisp] (h3) at (0.46,0) {};
  \draw[reac] (h1)--(hr); \draw[reac] (h2)--(hr); \draw[prod] (hr)--(h3);
\end{scope}
\node[rowmid] at (\cx+\name,\ryc) {hypergraph};

% =========================================================================
% (d) analysis outputs --- one band, four cards, one accent
% =========================================================================
\path[panel] (0,\dy) rectangle (\W,\dy-\dh);
\scHead{(0.22,\dy-0.16)}{d}{analysis outputs}{}
\node[anchor=north east,font=\scSmall,text=scInkM] at (\W-0.22,\dy-0.17)
  {exact arithmetic\quad$\cdot$\quad bounded search};

\def\dcw{3.825}\def\dcy{\dy-0.60}\def\dch{1.80}
% card centres, precomputed: TikZ would read the parentheses of an inline
% expression as a nested coordinate.
\pgfmathsetmacro{\cdA}{0.22+0.5*\dcw}
\pgfmathsetmacro{\cdB}{0.22+1*(\dcw+0.22)+0.5*\dcw}
\pgfmathsetmacro{\cdC}{0.22+2*(\dcw+0.22)+0.5*\dcw}
\pgfmathsetmacro{\cdD}{0.22+3*(\dcw+0.22)+0.5*\dcw}
\foreach \i/\ttl in {0/{CRNT structure},1/{conservation},
                     2/{Petri structure},3/{realizability}} {
  \pgfmathsetmacro{\cxx}{0.22+\i*(\dcw+0.22)}
  \path[rounded corners=2.5pt,fill=scPaper,draw=scRuleL,line width=0.4pt]
    (\cxx,\dcy) rectangle (\cxx+\dcw,\dcy-\dch);
  \node[hbullet,anchor=north west,minimum size=1.8mm] at (\cxx+0.16,\dcy-0.16) {};
  \node[anchor=north west,font=\scBody,text=scInk]
    at (\cxx+0.42,\dcy-0.115) {\bfseries\ttl};
}

% --- card 1: the deficiency identity
\node[font=\fontsize{9.5}{11}\selectfont,text=scInk] at (\cdA,\dy-1.42)
  {$\delta=n-\ell-\operatorname{rank}\mathbf S$};
\node[font=\scSmall,text=scInkM,align=center,text width=3.45cm]
  at (\cdA,\dy-2.02) {exact rational rank};

% --- card 2: a conservation law spans the chain
\begin{scope}[shift={(\cdB,\dy-1.36)}]
  \node[minisp] (c1) at (-1.30,0) {};
  \node[minirx] (cr) at (-0.65,0) {};
  \node[minisp] (c2) at (0,0) {};
  \node[minirx] (cs) at (0.65,0) {};
  \node[minisp] (c3) at (1.30,0) {};
  \draw[reac] (c1)--(cr); \draw[prod] (cr)--(c2);
  \draw[reac] (c2)--(cs); \draw[prod] (cs)--(c3);
  \draw[decorate,decoration={brace,amplitude=2.4pt,raise=1.6pt},
        draw=scAccent,line width=0.4pt] (1.44,-0.19) -- (-1.44,-0.19);
\end{scope}
\node[font=\scSmall,text=scInkM,align=center,text width=3.45cm]
  at (\cdB,\dy-2.02) {minimal $y\geq0$ with $y^{\mathsf T}\mathbf S=0$};

% --- card 3: a siphon --- once empty, never refilled
\begin{scope}[shift={(\cdC,\dy-1.36)}]
  \node[minisp,fill=scAlertF,draw=scAlert] (s1) at (-1.20,0) {};
  \node[minirx] (sr) at (-0.60,0) {};
  \node[minisp] (s2) at (0,0) {};
  \node[minirx] (ss) at (0.60,0) {};
  \node[minisp] (s3) at (1.20,0) {};
  \draw[reac] (s1)--(sr); \draw[prod] (sr)--(s2);
  \draw[reac] (s2)--(ss); \draw[prod] (ss)--(s3);
  \draw[scAlert,line width=0.55pt] (-1.32,-0.12)--(-1.08,0.12);
\end{scope}
\node[font=\scSmall,text=scInkM,align=center,text width=3.45cm]
  at (\cdC,\dy-2.02) {minimal siphons and traps};

% --- card 4: a replayable firing sequence
\begin{scope}[shift={(\cdD,\dy-1.30)}]
  \foreach \x/\n in {-1.15/1,-0.575/1,0/2,0.575/1,1.15/2} {
    \node[rounded corners=1pt,draw=scSpE,fill=scSpF,minimum width=3.8mm,
          minimum height=2.3mm,inner sep=0pt] at (\x,0) {};
    \ifnum\n=1 \node[token] at (\x,0) {};
    \else \node[token] at (\x-0.06,0) {}; \node[token] at (\x+0.06,0) {}; \fi
  }
  \draw[-{Stealth[round,length=1.6mm,width=1.3mm]},draw=scMark,line width=0.5pt]
    (-1.38,-0.30)--(1.38,-0.30);
\end{scope}
\node[font=\scSmall,text=scInkM,align=center,text width=3.45cm]
  at (\cdD,\dy-2.02) {a firing certificate, or inconclusive};

\end{tikzpicture}
